% Leçon 5 — Vocabulaire et compréhension de texte : héritage — Chinois Tle A — Cours signé OnBuch+
% Programme officiel MINESEC. Compiler avec : tectonic ZH05-vocabulaire-heritage.tex
\newcommand{\DOCMATIERE}{Chinois}
\newcommand{\DOCNIVEAU}{Tle A}
\newcommand{\DOCMODULE}{Module 1 : Vie familiale et intégration sociale}
\newcommand{\DOCLECON}{ZH05}
\newcommand{\DOCTITRE}{Leçon 5 — Vocabulaire et compréhension de texte : héritage}
% OnBuch+ — préambule « cours signés » ENRICHI (compilé avec tectonic / XeLaTeX)
% Système visuel « Pop » d'OnBuch+ : fond crème (jamais blanc), bordures encre
% épaisses, ombres « dures » décalées, titres Archivo Black en capitales.
% Version MATHÉMATIQUES : graphes (pgfplots/tikz), boîtes pédagogiques
% (définition, propriété, méthode, attention, le savais-tu, exercice résolu,
% exercice, TP, bilan), page de garde et signature OnBuch+ ancrée en bas.
%
% Macros attendues dans main.tex AVANT \input{preamble} :
%   \DOCMATIERE  \DOCNIVEAU  \DOCMODULE  \DOCLECON (numéro)  \DOCTITRE
\documentclass[11pt]{article}

\usepackage{fontspec}
\usepackage[french]{babel} % français = langue principale ; le chinois via \zh{..} / textebox (xeCJK)
\usepackage{xeCJK}
\usepackage{pifont} % \ding{51} \ding{55} utilisés par les modèles
\usepackage[a4paper,margin=2cm,top=3.4cm,bottom=2.8cm,headheight=1.4cm,footskip=1.3cm]{geometry}
\usepackage{amsmath,amssymb,amsfonts}
\usepackage{mathtools} % \xmapsto, \coloneqq... souvent utilisés par les modèles
\usepackage{cancel} % simplifications barrées (bilans de forces)
% \abs{z} (module, valeur absolue) et \norm{u} (norme), utilisés par les modèles
\providecommand{\abs}{}\let\abs\relax\DeclarePairedDelimiter{\abs}{\lvert}{\rvert}
\providecommand{\norm}{}\let\norm\relax\DeclarePairedDelimiter{\norm}{\lVert}{\rVert}
\usepackage[table]{xcolor} % [table] : fournit aussi \rowcolors (alternance de lignes)
\usepackage{pagecolor}
\usepackage{tikz}
\usetikzlibrary{arrows.meta,positioning,calc,shapes.geometric,shapes.misc,shapes.arrows,decorations.pathmorphing,decorations.pathreplacing,decorations.markings,patterns,shadows,mindmap,trees,fit,backgrounds,matrix,intersections,angles,quotes}
\usepackage{pgfplots}
\usepackage[version=4]{mhchem} % équations nucléaires \ce{^{235}_{92}U}
\usepackage[european,siunitx]{circuitikz} % schémas électriques
\pgfplotsset{compat=1.18}
\usepgfplotslibrary{fillbetween,groupplots} % aires entre courbes (calcul intégral)
\usepackage{enumitem}
\usepackage{fancyhdr}
% [most] charge toutes les bibliothèques tcolorbox (listings, minted, raster,
% documentation...) et gonfle inutilement la mémoire TeX sur un document long
% (~300 boîtes) : on ne charge que ce qui est réellement utilisé.
\usepackage[skins,breakable]{tcolorbox}
\usepackage{graphicx}
\usepackage{colortbl}
\usepackage{booktabs}
\usepackage{tabularx}
\usepackage{diagbox} % case d'en-tête diagonale des tableaux à double entrée
% Colonnes extensibles centrée / gauche / droite (C, L, R), souvent utilisées par les modèles
\newcolumntype{C}{>{\centering\arraybackslash}X}
\newcolumntype{L}{>{\raggedright\arraybackslash}X}
\newcolumntype{R}{>{\raggedleft\arraybackslash}X}
\usepackage{array}
\usepackage{multicol}
\usepackage{multirow}
\usepackage{siunitx}
% parse-numbers=false : accepte aussi \SI{2,5 \times 10^{-3}}{...}
\sisetup{locale=FR,output-decimal-marker={,},group-minimum-digits=4,parse-numbers=false}
\DeclareSIUnit\ohm{\ensuremath{\symup{\Omega}}} % U+2126 absent de la police
\DeclareSIUnit\division{div} % réglages d'oscilloscope (ms/div, V/div)
\DeclareSIUnit\tr{tr} % tours (vitesse de rotation, ex. \SI{1500}{\tr\per\minute})
\DeclareSIUnit\molar{M} % concentration molaire (ex. \SI{3}{\molar}, \SI{1}{\micro\molar})
\DeclareSIUnit\an{an} % année (le modèle confond parfois avec \year, un compteur TeX) — ex. \SI{5}{\kilo\meter\per\an}
\DeclareSIUnit\FCFA{FCFA} % franc CFA (ex. \SI{500}{\FCFA\per\kilo\gram})
\DeclareSIUnit\degreeBrix{\textdegree Brix} % teneur en sucre d'un sirop/jus (réfractométrie)
\DeclareSIUnit\month{mois} % le modèle utilise parfois \month, un compteur TeX, comme unité de durée
\usepackage{qrcode}
% \degree : commande fréquemment utilisée par les modèles pour un angle ou une
% température en dehors de \SI{}{} (non fournie par les paquets chargés).
\providecommand{\degree}{^{\circ}}
% \qed (fin de démonstration), utilisé par les modèles sans amsthm
\providecommand{\qed}{\hfill$\square$}
% \barre : double barre des valeurs interdites dans un tableau de variations (tkz-tab)
\providecommand{\barre}{\ensuremath{\|}}
% environnement proof (amsthm non chargé) utilisé par les modèles
\ifdefined\proof\else\newenvironment{proof}[1][Démonstration]{\par\noindent\textit{#1.}\ }{\qed\par}\fi
\usepackage{lastpage}
\usepackage{hyperref}
\hypersetup{hidelinks}

% ============================================================
% Palette « Pop » OnBuch+ — mêmes hex que la classe P (plus_ui.dart)
% ============================================================
\definecolor{poporange}{HTML}{F59321}
\definecolor{popdark}{HTML}{E07A0C}
\definecolor{popcream}{HTML}{FFF6E8}
\definecolor{popink}{HTML}{1C1714}
\definecolor{poppurple}{HTML}{7A5AE0}
\definecolor{popgreen}{HTML}{2E9E6A}
\definecolor{poppink}{HTML}{E0457B}
\definecolor{popblue}{HTML}{2F7DE1}
\definecolor{popgold}{HTML}{C98A12}
\definecolor{popmuted}{HTML}{8A7F76}
\definecolor{popline}{HTML}{EADFD2}
\definecolor{popgray}{HTML}{8A7F76}
\colorlet{popgrey}{popgray}
% Alias tolérants (le modèle nomme parfois les couleurs différemment)
\colorlet{popwhite}{white}
\colorlet{popblack}{popink}
\colorlet{popred}{poppink}
\colorlet{popyellow}{popgold}
% Teintes claires (fonds de tableaux/encadrés) — le modèle nomme parfois
% les couleurs avec un suffixe L (ex. popblueL) sans les avoir définies.
\colorlet{poporangeL}{poporange!20}
\colorlet{popdarkL}{popdark!20}
\colorlet{poppurpleL}{poppurple!20}
\colorlet{popgreenL}{popgreen!20}
\colorlet{poppinkL}{poppink!20}
\colorlet{popblueL}{popblue!20}
\colorlet{popgoldL}{popgold!20}
\colorlet{popredL}{popred!20}
\colorlet{popgrayL}{popgray!20}
% Teintes claires pour fonds de tableaux / zones
\colorlet{popblueL}{popblue!10!white}
\colorlet{poporangeL}{poporange!14!white}
\colorlet{popgreenL}{popgreen!12!white}
\colorlet{poppurpleL}{poppurple!12!white}
\colorlet{poppinkL}{poppink!12!white}
\colorlet{popgoldL}{popgold!14!white}

\pagecolor{popcream}

% ============================================================
% Typographie
% ============================================================
\defaultfontfeatures{Ligatures=TeX}
\setmainfont{PlusJakartaSans-Medium.ttf}[Path=../../../fonts/,BoldFont=PlusJakartaSans-Bold.ttf]
% Symboles, API et emojis absents de la police principale : repli sur DejaVu Sans (caractères actifs)
\newfontface\symfont{DejaVuSans.ttf}[Path=../../../fonts/]
\makeatletter
\newcommand\symactive[1]{\catcode#1=13 \begingroup\lccode`\~=#1 \lowercase{\endgroup\def~}{{\symfont\char#1}}}
\newcommand\symblank[1]{\catcode#1=13 \begingroup\lccode`\~=#1 \lowercase{\endgroup\def~}{}}
\newcommand\symhyph[1]{\catcode#1=13 \begingroup\lccode`\~=#1 \lowercase{\endgroup\def~}{\mbox{-}}}
\newcount\symc
\newcommand\symrange[2]{\symc=#1 \@whilenum\symc<#2 \do{\expandafter\symactive\expandafter{\the\symc}\advance\symc by 1 }}
\newcommand\symblankrange[2]{\symc=#1 \@whilenum\symc<#2 \do{\expandafter\symblank\expandafter{\the\symc}\advance\symc by 1 }}
\makeatother
\symrange{"0250}{"02FF}\symactive{"2713}\symactive{"2714}\symactive{"2717}\symactive{"2718}\symactive{"2610}\symactive{"2611}\symactive{"2612}
\symrange{"2600}{"27BF}
\catcode"2705=13 \begingroup\lccode`\~="2705 \lowercase{\endgroup\def~}{{\symfont\char"2713}}\symblank{"26BD}\symblank{"26A1}
\symrange{"2460}{"257F}\symactive{"223C}\symactive{"2248}\symactive{"2260}
\symhyph{"2011}\symhyph{"2E1A}\symblankrange{"1F300}{"1FAFF}\symblank{"FE0F}
% Caractères chinois : WenQuanYi Zen Hei (sous-ensemble GB2312 + pinyin), gras/italique simulés
\setCJKmainfont{WenQuanYiZenHei-subset.ttf}[Path=../../../fonts/,AutoFakeBold=3,AutoFakeSlant=0.2]
\xeCJKsetup{CJKspace=false}
% Pinyin : voyelles à caron (ǎ ǐ ǒ ǔ ǖ ǘ ǚ ǜ) absentes de la police principale -> caractères actifs
\newfontface\pyfont{WenQuanYiZenHei-subset.ttf}[Path=../../../fonts/]
\newcommand\pyactive[1]{\catcode#1=13 \begingroup\lccode`\~=#1 \lowercase{\endgroup\def~}{{\pyfont\char#1}}}
\pyactive{"01CD}\pyactive{"01CE}\pyactive{"01CF}\pyactive{"01D0}\pyactive{"01D1}\pyactive{"01D2}\pyactive{"01D3}\pyactive{"01D4}\pyactive{"01D5}\pyactive{"01D6}\pyactive{"01D7}\pyactive{"01D8}\pyactive{"01D9}\pyactive{"01DA}\pyactive{"01DB}\pyactive{"01DC}
% Maths en sans-serif (Fira Math) avec chiffres et lettres droites de Plus
% Jakarta, pour que les formules se fondent dans le texte.
\usepackage[math-style=ISO,bold-style=ISO]{unicode-math}
\setmathfont{FiraMath-Regular.otf}[Path=../../../fonts/]
\setmathfont{PlusJakartaSans-Medium.ttf}[Path=../../../fonts/,range={up/num,up/{Latin,latin},\mathup}]
\usepackage{icomma} % virgule décimale sans espace en mode math
% Caractères mathématiques Unicode absents des polices de texte (Plus Jakarta,
% Archivo Black) : rendus par leur équivalent mathématique, en texte comme en
% formule — sinon ils s'affichent en carré vide (ex. « Arithmétique dans ℤ »).
\usepackage{newunicodechar}
\newunicodechar{ℝ}{\ensuremath{\mathbb{R}}}
\newunicodechar{ℤ}{\ensuremath{\mathbb{Z}}}
\newunicodechar{ℂ}{\ensuremath{\mathbb{C}}}
\newunicodechar{ℕ}{\ensuremath{\mathbb{N}}}
\newunicodechar{ℚ}{\ensuremath{\mathbb{Q}}}
\newunicodechar{𝔻}{\ensuremath{\mathbb{D}}}
\newunicodechar{α}{\ensuremath{\alpha}}
\newunicodechar{β}{\ensuremath{\beta}}
\newunicodechar{γ}{\ensuremath{\gamma}}
\newunicodechar{δ}{\ensuremath{\delta}}
\newunicodechar{ε}{\ensuremath{\varepsilon}}
\newunicodechar{θ}{\ensuremath{\theta}}
\newunicodechar{λ}{\ensuremath{\lambda}}
\newunicodechar{μ}{\ensuremath{\mu}}
\newunicodechar{σ}{\ensuremath{\sigma}}
\newunicodechar{τ}{\ensuremath{\tau}}
\newunicodechar{φ}{\ensuremath{\varphi}}
\newunicodechar{ψ}{\ensuremath{\psi}}
\newunicodechar{ω}{\ensuremath{\omega}}
\newunicodechar{Σ}{\ensuremath{\Sigma}}
% majuscules produites par \MakeUppercase dans les titres (α-aminés -> Α)
\newunicodechar{Α}{\ensuremath{\alpha}}
\newunicodechar{Β}{\ensuremath{\beta}}
\newunicodechar{Γ}{\ensuremath{\Gamma}}
\newunicodechar{Δ}{\ensuremath{\Delta}}
\newunicodechar{Ε}{\ensuremath{\varepsilon}}
\newunicodechar{Θ}{\ensuremath{\Theta}}
\newunicodechar{Λ}{\ensuremath{\Lambda}}
\newunicodechar{Μ}{\ensuremath{\mu}}
\newunicodechar{Τ}{\ensuremath{\tau}}
\newunicodechar{Φ}{\ensuremath{\Phi}}
\newunicodechar{Ψ}{\ensuremath{\Psi}}
\newunicodechar{Ω}{\ensuremath{\Omega}}
\newunicodechar{↦}{\ensuremath{\mapsto}}
\newunicodechar{∈}{\ensuremath{\in}}
\newunicodechar{∉}{\ensuremath{\notin}}
\newunicodechar{∧}{\ensuremath{\wedge}}
\newunicodechar{⇔}{\ensuremath{\Leftrightarrow}}
\newunicodechar{⟺}{\ensuremath{\Longleftrightarrow}}
\newunicodechar{⇒}{\ensuremath{\Rightarrow}}
\newunicodechar{⟹}{\ensuremath{\Longrightarrow}}
\newunicodechar{∘}{\ensuremath{\circ}}
\newunicodechar{∪}{\ensuremath{\cup}}
\newunicodechar{∩}{\ensuremath{\cap}}
\newunicodechar{≡}{\ensuremath{\equiv}}
\newunicodechar{⊂}{\ensuremath{\subset}}
\newunicodechar{⊥}{\ensuremath{\perp}}
\newunicodechar{∃}{\ensuremath{\exists}}
\newunicodechar{∀}{\ensuremath{\forall}}
\newunicodechar{∛}{\ensuremath{\sqrt[3]{\,}}}
\newunicodechar{∜}{\ensuremath{\sqrt[4]{\,}}}
\newunicodechar{⃗}{\ensuremath{{}^{\to}}}
\newunicodechar{ⁿ}{\ifmmode^{n}\else\textsuperscript{\ensuremath{n}}\fi}
\newunicodechar{ˣ}{\ifmmode^{x}\else\textsuperscript{\ensuremath{x}}\fi}
\newunicodechar{ᵇ}{\ifmmode^{b}\else\textsuperscript{\ensuremath{b}}\fi}
\newunicodechar{ʳ}{\ifmmode^{r}\else\textsuperscript{\ensuremath{r}}\fi}
\newunicodechar{ᵘ}{\ifmmode^{u}\else\textsuperscript{\ensuremath{u}}\fi}
\newunicodechar{ʸ}{\ifmmode^{y}\else\textsuperscript{\ensuremath{y}}\fi}
\newunicodechar{ᵃ}{\ifmmode^{a}\else\textsuperscript{\ensuremath{a}}\fi}
\newunicodechar{ᵉ}{\ifmmode^{e}\else\textsuperscript{\ensuremath{e}}\fi}
\newunicodechar{ᵏ}{\ifmmode^{k}\else\textsuperscript{\ensuremath{k}}\fi}
\newunicodechar{ᵖ}{\ifmmode^{p}\else\textsuperscript{\ensuremath{p}}\fi}
\newunicodechar{ᴺ}{\ifmmode^{N}\else\textsuperscript{\ensuremath{N}}\fi}
\newunicodechar{ᶜ}{\ifmmode^{c}\else\textsuperscript{\ensuremath{c}}\fi}
\newunicodechar{ʰ}{\ifmmode^{h}\else\textsuperscript{\ensuremath{h}}\fi}
\newunicodechar{ᵠ}{\ifmmode^{q}\else\textsuperscript{\ensuremath{q}}\fi}
\newunicodechar{ₙ}{\ifmmode_{n}\else\textsubscript{\ensuremath{n}}\fi}
\newunicodechar{ₐ}{\ifmmode_{a}\else\textsubscript{\ensuremath{a}}\fi}
\newunicodechar{ᵢ}{\ifmmode_{i}\else\textsubscript{\ensuremath{i}}\fi}
\newunicodechar{ᵣ}{\ifmmode_{r}\else\textsubscript{\ensuremath{r}}\fi}
\newunicodechar{ₖ}{\ifmmode_{k}\else\textsubscript{\ensuremath{k}}\fi}
\newunicodechar{ᵦ}{\ifmmode_{\beta}\else\textsubscript{\ensuremath{\beta}}\fi}
\newunicodechar{ₓ}{\ifmmode_{x}\else\textsubscript{\ensuremath{x}}\fi}
\newfontfamily\popdisplay{ArchivoBlack-Regular.ttf}[Path=../../../fonts/]
\newfontfamily\poplogo{SpaceGrotesk-Bold.ttf}[Path=../../../fonts/]
\newfontfamily\popsemi{PlusJakartaSans-SemiBold.ttf}[Path=../../../fonts/]
\newfontfamily\popxbold{PlusJakartaSans-ExtraBold.ttf}[Path=../../../fonts/]
\newfontfamily\popxboldls{PlusJakartaSans-ExtraBold.ttf}[Path=../../../fonts/,LetterSpace=3.0]

\setlist[enumerate]{leftmargin=1.7em,itemsep=5pt,topsep=4pt}
\setlist[itemize]{leftmargin=1.7em,itemsep=4pt,topsep=4pt,label={\color{poporange}\textbullet}}
\setlength{\parindent}{0pt}
\setlength{\parskip}{5pt}
\renewcommand{\arraystretch}{1.25}

% pgfplots : style Pop par défaut
\pgfplotsset{
  every axis/.append style={axis line style={popink,line width=1pt},tick style={popink},
    grid style={popline,line width=0.5pt},label style={font=\small},tick label style={font=\footnotesize}},
  popcurve/.style={poporange,line width=1.8pt,no markers,smooth},
}
% Même style disponible dans un \draw TikZ simple (hors pgfplots)
\tikzset{popcurve/.style={poporange,line width=1.8pt}}
\tikzset{popgrid/.style={popline,very thin}}
\colorlet{popcurve}{poporange} % le nom du style est parfois utilisé comme couleur

% ============================================================
% Pill / squiggle
% ============================================================
\newcommand{\poppill}[3]{%
  \tikz[baseline=-0.55ex]\node[draw=popink,line width=1.2pt,rounded corners=2.6pt,%
    fill=#2,inner xsep=6.5pt,inner ysep=3pt]{%
    \popxboldls\fontsize{7}{7}\selectfont\color{#3}\MakeUppercase{#1}};%
}
\newcommand{\popsquiggle}[1]{%
  \tikz[baseline=-0.4ex]{
    \draw[#1,line width=2.6pt,line cap=round]
      plot [smooth] coordinates {(0,0.09) (0.34,0.01) (0.68,0.21) (1.02,0.01) (1.36,0.10)};
  }%
}
% Mot-clé surligné (marqueur orange sous le texte)
\newcommand{\cle}[1]{{\popxbold\color{popdark}#1}}

% ============================================================
% En-tête / pied
% ============================================================
\pagestyle{fancy}
\fancyhf{}
\renewcommand{\headrulewidth}{0pt}
\renewcommand{\footrulewidth}{0pt}
\lhead{\raisebox{-0.15ex}{\poplogo\fontsize{13}{13}\selectfont\color{popink}OnBuch\color{poporange}+}}
\rhead{\poppill{\DOCMATIERE\ \textbullet\ \DOCNIVEAU\ \textbullet\ Leçon \DOCLECON}{white}{popink}}
\lfoot{\popxbold\fontsize{7.6}{7.6}\selectfont\color{popmuted}\MakeUppercase{Cours signé OnBuch+}\ \textbullet\ {\popsemi\DOCTITRE}}
\rfoot{\tikz[baseline=-0.6ex]\node[rounded corners=8.5pt,fill=popink,minimum height=17pt,minimum width=17pt,inner xsep=5pt,inner sep=0]{\popxbold\fontsize{8}{8}\selectfont\color{popcream}\thepage\,/\,\pageref*{LastPage}};}

% ============================================================
% Titres
% ============================================================
\newcommand{\chaptitre}[2]{%
  \begin{tcolorbox}[enhanced,colback=poporange,colframe=popink,boxrule=2.6pt,arc=11pt,
      drop shadow=popink,left=15pt,right=15pt,top=12pt,bottom=11pt,before skip=3pt,after skip=18pt]
    \poppill{#2}{popink}{popcream}\par\vspace{9pt}
    {\raggedright\hyphenpenalty=10000\exhyphenpenalty=10000\popdisplay\fontsize{22}{24}\selectfont\color{popcream}\MakeUppercase{#1}\par}
  \end{tcolorbox}
}
% Section numérotée : pastille orange avec le numéro + titre Archivo
\newcounter{popsec}
\newcommand{\coursec}[1]{%
  \stepcounter{popsec}\setcounter{popsub}{0}%
  \par\vspace{16pt}\noindent
  \begin{minipage}{\linewidth}
  \tikz[baseline=-0.7ex]\node[draw=popink,line width=1.6pt,fill=poporange,rounded corners=4pt,minimum size=22pt,inner sep=0]{\popdisplay\fontsize{12}{12}\selectfont\color{popcream}\thepopsec};%
  \hspace{8pt}{\popdisplay\fontsize{15}{17}\selectfont\color{popink}\MakeUppercase{#1}}\par
  \vspace{4pt}\noindent{\color{poporange}\rule{36pt}{3.6pt}}
  \end{minipage}\par\nopagebreak\vspace{8pt}
}
\newcounter{popsub}
\newcommand{\courssub}[1]{%
  \stepcounter{popsub}%
  \par\vspace{9pt}\noindent{\popxbold\fontsize{11.5}{13}\selectfont\color{popink}{\color{poporange}\thepopsec.\thepopsub}\hspace{6pt}#1}\par\nopagebreak\vspace{4pt}
}

% ============================================================
% Boîtes pédagogiques — même recette : fond blanc, bordure épaisse colorée,
% ombre dure encre, badge plein. Titre optionnel [#1].
% ============================================================
\tcbset{popbase/.style={breakable,enhanced,colback=white,boxrule=2.3pt,arc=9pt,
  shadow={3pt}{-3pt}{0pt}{popink},left=14pt,right=14pt,top=11pt,bottom=11pt,before skip=11pt,after skip=15pt,
  fontupper=\popsemi\fontsize{10.6}{14.5}\selectfont\color{popink},
  fontlower=\popsemi\fontsize{10.6}{14.5}\selectfont\color{popink}}}
% Coche dessinée (le glyphe ✓ n'existe pas dans les polices du document)
\newcommand{\popcheck}[1]{\tikz[baseline=-0.5ex]\draw[#1,line width=1.6pt,line cap=round,line join=round](-0.13,0.0)--(-0.04,-0.09)--(0.13,0.1);}
\providecommand{\coche}{\popcheck{popgreen}}
\providecommand{\resultat}[1]{\par\smallskip\noindent\fcolorbox{popgreen}{popgreenL}{\parbox{\dimexpr\linewidth-2\fboxsep-2\fboxrule\relax}{\textbf{Résultat.} #1}}\par\smallskip}
\newcommand{\popboxhead}[3]{\poppill{#1}{#2}{white}\ifx\relax#3\relax\else\hspace{6pt}{\popxbold\fontsize{10.6}{12}\selectfont\color{popink}#3}\fi\par\vspace{7pt}}

\newtcolorbox{aretenir}[1][]{popbase,colframe=popgreen,before upper={\popboxhead{À retenir}{popgreen}{#1}}}
% Texte en chinois (document de lecture, dialogue, extrait) : cadre
% d'étude. Un titre optionnel (source, auteur) s'affiche dans l'en-tête.
\newtcolorbox{textebox}[1][]{popbase,colframe=popblue,before upper={\popboxhead{文本 · Texte}{popblue}{#1}},after upper={}}
% Marques ✓ / ✗ (forme correcte / fautive) souvent utilisées par les modèles
\providecommand{\cross}{\textcolor{poppink}{\ensuremath{\times}}}
\providecommand{\xmark}{\cross}\providecommand{\crossmark}{\cross}\providecommand{\cmark}{\ensuremath{\checkmark}}
% Ligne de réponse / justification à compléter (inventée par les modèles)
\providecommand{\justification}[1]{\par\noindent\textit{Justification~:} \dotfill\par}
% \textipa (package tipa non chargé) : la police ne couvre pas l'API -> simple passage
\providecommand{\textipa}[1]{#1}
% Soulignement (ulem non chargé) : \uline, \ul
\providecommand{\uline}[1]{\underline{#1}}\providecommand{\ul}[1]{\underline{#1}}
% \glossaire{\item ...} inventée par les modèles : liste de glossaire
\providecommand{\glossaire}[1]{\begin{itemize}#1\end{itemize}}
% \alert (beamer) inventée par les modèles : texte mis en évidence
\providecommand{\alert}[1]{\textbf{\textcolor{poppink}{#1}}}
% variantes ASCII d'ordinaux écrites en macro
\providecommand{\iere}{\textsuperscript{re}}\providecommand{\ieme}{\textsuperscript{e}}\providecommand{\ere}{\textsuperscript{re}}\providecommand{\eme}{\textsuperscript{e}}\providecommand{\er}{\textsuperscript{er}}\providecommand{\ie}{\textsuperscript{e}}
% Ordinaux français écrits en macro par les modèles : 1\ière, 3\ième
\newcommand{\ière}{\textsuperscript{re}}\newcommand{\ième}{\textsuperscript{e}}
% \exemple (inventée par les modèles) : étiquette « Exemple : »
\providecommand{\exemple}{\textbf{Exemple~:}\ }
% \square utilisable aussi hors mode math (cases à cocher)
\AtBeginDocument{\let\squareMath\square\renewcommand{\square}{\ensuremath{\squareMath}}}
% Mot ou phrase en chinois à l'intérieur d'un paragraphe français
\newcommand{\zh}[1]{\ifmmode\text{#1}\else{#1}\fi}
\let\eng\zh \let\esp\zh \let\all\zh % alias tolérés
% flèches d'intonation inventées par les modèles
\providecommand{\textsearrow}{}\renewcommand{\textsearrow}{\ensuremath{\searrow}}\providecommand{\textnearrow}{}\renewcommand{\textnearrow}{\ensuremath{\nearrow}}
% \fr{..} : mot français (inventé par les modèles)
\providecommand{\fr}[1]{\textit{#1}}
% zhbox : encadré de texte chinois inventé par les modèles
\newenvironment{zhbox}{\begin{quote}}{\end{quote}}
% \vs : « versus » inventé par les modèles
\providecommand{\vs}{\textit{vs}\ }
% \blank : case à compléter inventée par les modèles
\providecommand{\blank}[1][]{\underline{\hspace{1.6em}}}
\newtcolorbox{exemplebox}[1][]{popbase,colframe=poporange,before upper={\popboxhead{Exemple}{poporange}{#1}}}
\newtcolorbox{definition}[1][]{popbase,colframe=popblue,before upper={\popboxhead{Définition}{popblue}{#1}}}
\newtcolorbox{propriete}[1][]{popbase,colframe=poppurple,before upper={\popboxhead{Propriété}{poppurple}{#1}}}
\newtcolorbox{methode}[1][]{popbase,colframe=popgold,before upper={\popboxhead{Méthode}{popgold}{#1}}}
\newtcolorbox{attention}[1][]{popbase,colframe=poppink,before upper={\popboxhead{Attention}{poppink}{#1}}}
\newtcolorbox{experience}[1][]{popbase,colframe=popblue,colback=popblueL,before upper={\popboxhead{Expérience}{popblue}{#1}}}
% « Le savais-tu ? » : carte violette pleine (contexte, culture, Cameroun)
\newtcolorbox{savaistu}[1][]{popbase,colback=poppurple,colframe=popink,
  fontupper=\popsemi\fontsize{10.4}{14.2}\selectfont\color{white},
  before upper={\poppill{Le savais-tu\,?}{white}{poppurple}\ifx\relax#1\relax\else\hspace{6pt}{\popxbold\color{white}#1}\fi\par\vspace{7pt}}}
% Exercice résolu : énoncé en haut, \tcblower puis la solution sur fond crème
\newcounter{popexo}\newcounter{popexr}
\newtcolorbox{exoresolu}[1][]{popbase,colframe=poporange,colbacklower=popcream,
  segmentation style={popink,line width=1.2pt,dashed},
  before upper={\stepcounter{popexr}\popboxhead{Exercice résolu \thepopexr}{poporange}{#1}},
  before lower={\poppill{Solution}{popink}{popcream}\par\vspace{6pt}}}
% Exercice (non résolu en place) — la correction est dans la partie Corrigés
\newtcolorbox{exercice}[1][]{popbase,colframe=popink,
  before upper={\stepcounter{popexo}\popboxhead{Exercice \thepopexo}{popink}{#1}}}
% Corrigé
\newtcolorbox{corrige}[1][]{popbase,colframe=popgreen,colback=popgreenL,
  before upper={\popboxhead{Corrigé}{popgreen}{#1}}}
% Objectifs / prérequis / bilan
\newtcolorbox{objectifs}{popbase,colframe=popink,colback=poporangeL,
  before upper={\popboxhead{Objectifs de la leçon}{popink}{}}}
\newtcolorbox{prerequis}{popbase,colframe=popink,colback=popblueL,
  before upper={\popboxhead{Prérequis}{popblue}{}}}
\newtcolorbox{bilan}{popbase,colframe=popink,colback=white,boxrule=2.8pt,
  before upper={\popboxhead{Fiche bilan}{poporange}{L'essentiel à connaître pour l'examen}}}
% Figure encadrée avec légende
\newtcolorbox{popfigure}[1][]{enhanced,breakable=false,colback=white,colframe=popline,boxrule=1.4pt,arc=8pt,
  left=10pt,right=10pt,top=10pt,bottom=8pt,before skip=10pt,after skip=12pt,halign=center,
  lower separated=false,halign lower=center,
  fontlower=\popsemi\fontsize{9}{11.5}\selectfont\color{popmuted},
  #1}
\newcommand{\legende}[1]{\par\vspace{6pt}{\popsemi\fontsize{9}{11.5}\selectfont\color{popmuted}#1}}

% ============================================================
% Signature OnBuch+
% ============================================================
\newcommand{\poplogoinline}[1]{{\poplogo\fontsize{#1}{#1}\selectfont\color{popink}OnBuch\color{poporange}+}}
\newcommand{\signatureonbuch}{%
  \begin{tcolorbox}[enhanced,colback=popink,colframe=popink,boxrule=2.2pt,arc=10pt,
      drop shadow=poporange,left=14pt,right=14pt,top=10pt,bottom=10pt,before skip=0pt,after skip=0pt]
    \tikz[baseline=-0.7ex]\node[circle,fill=popgreen,draw=popcream,line width=1.3pt,minimum size=22pt,inner sep=0]{\popcheck{white}};%
    \hspace{9pt}\begin{minipage}[c]{0.62\linewidth}
      {\popxbold\fontsize{12}{13}\selectfont\color{popcream}Signé\ \textbullet\ {\poplogo OnBuch\color{poporange}+}}\par\vspace{2pt}
      {\raggedright\popsemi\fontsize{8}{10.5}\selectfont\color{popline}\MakeUppercase{Équipe pédagogique OnBuch+}\\\MakeUppercase{Conforme au programme officiel MINESEC}\par}
    \end{minipage}\hfill
    {\poplogo\fontsize{22}{22}\selectfont\color{popcream}OnBuch\color{poporange}+}
  \end{tcolorbox}%
}

% ============================================================
% Page de garde
% #1 titre · #2 matière · #3 niveau · #4 module · #5 n° leçon · #6 durée ·
% #7 description / objectifs (texte ou itemize)
% ============================================================
\newcommand{\pagegarde}[7]{%
  \thispagestyle{empty}
  \newgeometry{a4paper,margin=1.7cm,top=1.6cm,bottom=1.4cm}%
  \begin{tikzpicture}[remember picture,overlay]
    \fill[poporange!18] ([xshift=-3.2cm,yshift=-3.6cm]current page.north east) circle (4.6cm);
    \fill[poppurple!14] ([xshift=2.4cm,yshift=7.6cm]current page.south west) circle (3.8cm);
  \end{tikzpicture}%
  {\poplogo\fontsize{30}{30}\selectfont\color{popink}OnBuch\color{poporange}+}\hfill
  \raisebox{0.6ex}{\poppill{Cours signé \textbullet\ Premium}{popink}{popcream}}\par
  \vspace{1pt}\popsquiggle{poppurple}\par
  \vspace{8pt}
  \begin{tcolorbox}[enhanced,colback=poporange,colframe=popink,boxrule=2.8pt,arc=13pt,
      drop shadow=popink,left=18pt,right=18pt,top=12pt,bottom=13pt,after skip=10pt]
    \poppill{#2 \textbullet\ #3}{popink}{popcream}\hspace{5pt}\poppill{#4}{white}{popink}\par\vspace{8pt}
    {\popdisplay\fontsize{14}{14}\selectfont\color{popink}LEÇON #5}\par\vspace{4pt}
    {\raggedright\hyphenpenalty=10000\exhyphenpenalty=10000\popdisplay\fontsize{25}{27}\selectfont\color{popcream}\MakeUppercase{#1}\par}
    \vspace{8pt}
    {\popxbold\fontsize{9.5}{9.5}\selectfont\color{popink}\MakeUppercase{Durée conseillée : #6}}
  \end{tcolorbox}
  \begin{tcolorbox}[enhanced,colback=white,colframe=popink,boxrule=2.2pt,arc=10pt,
      drop shadow=popink,left=16pt,right=16pt,top=10pt,bottom=10pt,after skip=8pt]
    \poppill{Dans cette leçon}{poppurple}{white}\par\vspace{5pt}
    \popsemi\fontsize{9.3}{12.6}\selectfont\color{popink}#7
  \end{tcolorbox}
  \noindent\begin{minipage}[t]{0.487\linewidth}
    \begin{tcolorbox}[enhanced,colback=popgreen,colframe=popink,boxrule=2.2pt,arc=10pt,
        drop shadow=popink,left=11pt,right=11pt,top=10pt,bottom=10pt,height=3.1cm,valign=center]
      \tcbox[colback=white,colframe=popink,boxrule=1.3pt,arc=5pt,boxsep=0pt,nobeforeafter,
        left=3pt,right=3pt,top=3pt,bottom=3pt,valign=center]{\qrcode[height=1.9cm]{https://play.google.com/store/apps/details?id=com.luvvix.onbuch&hl=fr}}%
      \hspace{8pt}\begin{minipage}[c]{0.5\linewidth}
        {\popdisplay\fontsize{11}{12.5}\selectfont\color{white}\MakeUppercase{Télécharge OnBuch}}\par\vspace{4pt}
        {\raggedright\popsemi\fontsize{8.4}{11}\selectfont\color{white}Gratuit \textbullet\ Google Play\\Tuteur IA, annales, concours, résultats\par}
      \end{minipage}
    \end{tcolorbox}
  \end{minipage}\hfill
  \begin{minipage}[t]{0.487\linewidth}
    \begin{tcolorbox}[enhanced,colback=poppurple,colframe=popink,boxrule=2.2pt,arc=10pt,
        drop shadow=popink,left=11pt,right=11pt,top=10pt,bottom=10pt,height=3.1cm,valign=center]
      \tikz[baseline=-0.7ex]\node[circle,fill=white,draw=popink,line width=1.3pt,minimum size=1.6cm,inner sep=0]{\popdisplay\fontsize{24}{24}\selectfont\color{poppurple}+};%
      \hspace{8pt}\begin{minipage}[c]{0.6\linewidth}
        {\popdisplay\fontsize{11}{12.5}\selectfont\color{white}\MakeUppercase{Passe à OnBuch+}}\par\vspace{4pt}
        {\raggedright\popsemi\fontsize{8.4}{11}\selectfont\color{white}Tous les cours signés, Tuteur IA illimité, fiches PDF — onglet OnBuch+ de l'app\par}
      \end{minipage}
    \end{tcolorbox}
  \end{minipage}
  \vspace{8pt}
  \popcontenu
  \vfill
  \signatureonbuch
  \restoregeometry
  \newpage
}

% Rangée « Ce cours contient » (page de garde)
\newcommand{\poptile}[3]{% #1 couleur #2 gros texte #3 légende
  \begin{tcolorbox}[enhanced,colback=#1,colframe=popink,boxrule=2pt,arc=9pt,shadow={2.5pt}{-2.5pt}{0pt}{popink},
      left=6pt,right=6pt,top=9pt,bottom=9pt,halign=center,valign=center,height=1.75cm,width=\linewidth,nobeforeafter]
    {\popdisplay\fontsize{12.5}{13}\selectfont\color{white}\MakeUppercase{#2}}\par\vspace{3pt}
    {\centering\popsemi\fontsize{7.8}{9.5}\selectfont\color{white}#3\par}
  \end{tcolorbox}}
\newcommand{\popcontenu}{%
  \par\noindent{\popxboldls\fontsize{7.5}{7.5}\selectfont\color{popmuted}CE COURS SIGNÉ CONTIENT}\par\vspace{7pt}
  \noindent\begin{minipage}[t]{0.235\linewidth}\poptile{popblue}{Cours}{complet, illustré, pas à pas}\end{minipage}\hfill
  \begin{minipage}[t]{0.235\linewidth}\poptile{poporange}{Méthodes}{et exercices résolus}\end{minipage}\hfill
  \begin{minipage}[t]{0.235\linewidth}\poptile{poppink}{Exercices}{type examen + corrigés}\end{minipage}\hfill
  \begin{minipage}[t]{0.235\linewidth}\poptile{popgreen}{Fiche bilan}{l'essentiel à retenir}\end{minipage}\par}

% Sommaire
\newtcolorbox{sommaire}{enhanced,colback=white,colframe=popink,boxrule=2.2pt,arc=10pt,
  drop shadow=popink,left=18pt,right=18pt,top=13pt,bottom=13pt,before skip=6pt,after skip=20pt,
  before upper={\poppill{Sommaire}{poppurple}{white}\par\vspace{9pt}\popsemi\fontsize{10.6}{14.5}\selectfont\color{popink}}}

% Signature de fin de cours — poussée tout en bas de la dernière page
\newcommand{\findecours}{%
  \par\vfill\nopagebreak
  \begin{center}\popsquiggle{poporange}\end{center}\vspace{4pt}
  \signatureonbuch
}
\AtBeginDocument{\def\llbracket{\mathopen{[\mkern-3.5mu[}}\def\rrbracket{\mathclose{]\mkern-3.5mu]}}} % intervalles d'entiers (glyphe ⟦ absent de la police)
% Glyphes absents de Fira Math : redéfinis à partir de glyphes disponibles
\AtBeginDocument{%
  \def\setminus{\mathbin{\backslash}}\let\smallsetminus\setminus
  \def\vdots{\mathord{\vbox{\baselineskip3.2pt\lineskiplimit0pt\kern4pt\hbox{.}\hbox{.}\hbox{.}}}}%
  \def\ddots{\mathinner{\mkern1mu\raise6.5pt\hbox{.}\mkern2mu\raise3.6pt\hbox{.}\mkern2mu\raise0.7pt\hbox{.}\mkern1mu}}%
  \def\triangle{\mathord{\scalebox{0.9}{$\symup{\Delta}$}}}%
  \def\nabla{\mathord{\raisebox{\depth}{\scalebox{1}[-1]{$\symup{\Delta}$}}}}%
  \def\checkmark{\popcheck{popdark}}%
  \def\star{\mathbin{\ast}}\def\bigstar{\mathord{\ast}}%
  \def\diamond{\mathbin{\scalebox{0.6}{$\Diamond$}}}%
  \def\bigcup{\mathop{\mathchoice{\raisebox{-0.3em}{\scalebox{1.7}{$\cup$}}}{\scalebox{1.25}{$\cup$}}{\cup}{\cup}}\displaylimits}%
  \def\bigcap{\mathop{\mathchoice{\raisebox{-0.3em}{\scalebox{1.7}{$\cap$}}}{\scalebox{1.25}{$\cap$}}{\cap}{\cap}}\displaylimits}%
  \def\lesseqqgtr{\mathrel{\lessgtr}}\def\bigoplus{\mathop{\oplus}\displaylimits}%
}
\providecommand{\ssi}{si et seulement si} % abréviation fréquente
\AtBeginDocument{\providecommand{\wideparen}{\overparen}} % arc de cercle

%% FIGFIX-BEGIN v5 — correctifs globaux des figures (docs/onbuch-generation/tools/figfix)
\usepackage{adjustbox}\usepackage{etoolbox}\tcbuselibrary{raster}
% toute figure TikZ/pgfplots plus large que la ligne est réduite à la largeur disponible
\BeforeBeginEnvironment{tikzpicture}{\begin{adjustbox}{max width=\linewidth}}
\AfterEndEnvironment{tikzpicture}{\end{adjustbox}}
% légendes pgfplots lisibles par-dessus les courbes
\pgfplotsset{every axis/.append style={legend style={fill=white,fill opacity=0.92,text opacity=1,draw=black!15,inner sep=3pt}}}
% étiquettes d'arêtes (syntaxe edge["..."]) avec fond blanc
\tikzset{every edge quotes/.append style={fill=white,inner sep=1.2pt}}
%% FIGFIX-END
\begin{document}

\pagegarde{Leçon 5 — Vocabulaire et compréhension de texte : héritage}{Chinois}{Tle A}{Module 1 : Vie familiale et intégration sociale}{ZH05}{2 h}{Cette leçon de vocabulaire et compréhension de texte porte sur le thème de l'héritage (遗产). À partir d'un texte support original, les élèves enrichissent leur lexique, étudient les clés et l'ordre des traits de caractères clés, puis répondent à des questions de compréhension.
\vspace{4pt}
\begin{multicols}{2}\small
\begin{itemize}
\item À la fin de cette leçon, je sais lire et comprendre un court texte chinois sur le thème de l'héritage
\item À la fin de cette leçon, je sais reconnaître, prononcer (pinyin et tons) et écrire le vocabulaire du thème de l'héritage
\item À la fin de cette leçon, je sais identifier les clés (radicaux) et tracer l'ordre des traits de quelques caractères du thème
\item À la fin de cette leçon, je sais employer les collocations essentielles : 留下遗产, 继承遗产, 分财产, 立遗嘱
\item À la fin de cette leçon, je sais répondre à des questions de compréhension sur un texte en phrases complètes
\item À la fin de cette leçon, je sais distinguer héritage matériel (财产) et héritage culturel (文化遗产)
\end{itemize}\end{multicols}}

\begin{sommaire}
\begin{multicols}{2}
\begin{enumerate}
\item Mise en route et découverte du texte
\item Texte support : lecture, pinyin, traduction
\item Glossaire du thème en tableau
\item Clés (radicaux) et ordre des traits
\item Compréhension du texte et collocations
\item Production guidée et réinvestissement
\item Activité d'intégration
\item Méthodes et exercices résolus
\item Exercices
\item Corrigés des exercices
\item Fiche bilan
\end{enumerate}
\end{multicols}
\end{sommaire}

\begin{objectifs}
\begin{itemize}
\item À la fin de cette leçon, je sais lire et comprendre un court texte chinois sur le thème de l'héritage
\item À la fin de cette leçon, je sais reconnaître, prononcer (pinyin et tons) et écrire le vocabulaire du thème de l'héritage
\item À la fin de cette leçon, je sais identifier les clés (radicaux) et tracer l'ordre des traits de quelques caractères du thème
\item À la fin de cette leçon, je sais employer les collocations essentielles : 留下遗产, 继承遗产, 分财产, 立遗嘱
\item À la fin de cette leçon, je sais répondre à des questions de compréhension sur un texte en phrases complètes
\item À la fin de cette leçon, je sais distinguer héritage matériel (财产) et héritage culturel (文化遗产)
\item À la fin de cette leçon, je sais rédiger un court paragraphe guidé sur l'héritage dans ma famille ou mon pays
\end{itemize}
\end{objectifs}

\begin{prerequis}
\begin{itemize}
\item Vocabulaire de la famille : 爸爸, 妈妈, 爷爷, 奶奶, 儿子, 女儿 (classes antérieures)
\item Structure de base sujet-verbe-objet et particule 的 pour la possession
\item Verbes fréquents : 有, 给, 是, 在
\item Lecture du pinyin et maîtrise des quatre tons
\item Notion de mesure et nombres (pour parler de biens et de parts)
\end{itemize}
\end{prerequis}

\newpage

\coursec{Mise en route et découverte du texte}

\courssub{Situation d'accroche : un partage de famille}

Imaginez la scène. À Bamenda comme à Yaoundé, quand un grand-père décède, toute la famille se réunit : oncle de Douala, tante de Bafoussam, cousins arrivés de loin. On s'assoit sous le manguier ou dans le salon, et on partage ce que le défunt a laissé : la maison familiale, les champs de manioc, parfois un petit commerce au marché, quelques têtes de bétail. Des discussions parfois vives, des souvenirs, des larmes aussi.

En Chine, la scène est étonnamment semblable. La question de l'\cle{héritage} — \zh{遗产} (\emph{yíchǎn}) — est au cœur de la vie familiale, et le droit chinois moderne encadre ce partage avec des règles précises. Mais attention : l'héritage, ce n'est pas seulement l'argent et les maisons ! C'est aussi la langue que l'on parle, les traditions que l'on transmet, les recettes de grand-mère, les chansons, les valeurs. En chinois, on parle d'ailleurs de \zh{文化遗产} (\emph{wénhuà yíchǎn}), le « patrimoine culturel », comme la Grande Muraille ou l'opéra de Pékin.

\textbf{Problématique de la leçon :} comment parler de l'héritage — biens matériels et richesses culturelles — en chinois ? Pour le découvrir, nous allons lire et comprendre un court texte sur l'héritage d'une famille chinoise. Mais d'abord, activons ce que vous savez déjà.

\courssub{Brainstorming : que laisse-t-on en héritage au Cameroun ?}

Avant d'ouvrir le texte chinois, réfléchissons ensemble en français. Posez-vous la question : \emph{quand quelqu'un meurt, que laisse-t-il derrière lui ?}

\begin{experience}[Brainstorming collectif — 5 minutes]
L'enseignant écrit au tableau toutes les réponses des élèves, en français, puis ajoute peu à peu le mot chinois correspondant (sans exiger la mémorisation à ce stade). Pistes attendues :
\begin{itemize}
\item les biens matériels : la \textbf{maison}, le \textbf{champ}, le \textbf{bétail}, l'\textbf{argent}, un \textbf{commerce}, des meubles ;
\item les biens immatériels : les \textbf{traditions}, la \textbf{langue}, les \textbf{recettes de cuisine} (le ndolé de grand-mère !), les \textbf{histoires} de la famille, les \textbf{valeurs} (le travail, le respect) ;
\item les documents : le \textbf{testament}, les papiers de propriété.
\end{itemize}
Question relance : « Chez vous, qui hérite ? Les fils ? Les filles ? Les deux ? » — gardez les réponses en tête, nous y reviendrons avec la Chine.
\end{experience}

Ce brainstorming a un but précis : quand vous lirez le texte chinois, vous y retrouverez ces mêmes idées. Votre cerveau est maintenant prêt à les reconnaître sous leur forme chinoise.

\courssub{Le mot-clé de la leçon : 遗产 (yíchǎn)}

Observons d'abord le mot central de la leçon, avant même de le définir : \zh{遗产}. Il est composé de deux caractères que l'on peut analyser séparément.

\begin{definition}[遗产 yíchǎn — héritage]
\zh{遗产} (\emph{yíchǎn}) est un \textbf{nom} qui signifie « héritage » : ce qu'une personne laisse après sa mort, qu'il s'agisse de biens matériels (maison, argent, terres) ou de biens culturels (langue, traditions, œuvres).\\
Composition du mot :
\begin{itemize}
\item \zh{遗} (\emph{yí}) : laisser (derrière soi), léguer ; on le retrouve dans \zh{遗忘} (\emph{yíwàng}, oublier) ;
\item \zh{产} (\emph{chǎn}) : bien, propriété, produit ; on le retrouve dans \zh{财产} (\emph{cáichǎn}, biens, fortune) et \zh{生产} (\emph{shēngchǎn}, produire).
\end{itemize}
Littéralement : « ce qui est laissé comme bien » → l'héritage.
\end{definition}

\begin{textebox}[遗产 (yíchǎn) — héritage]
遗产是一个人去世后留下的东西。\\
\emph{Yíchǎn shì yí ge rén qùshì hòu liúxià de dōngxi.}\\
« L'héritage, c'est ce qu'une personne laisse après sa mort. »\\[4pt]
遗产可以是房子、钱，也可以是文化和传统。\\
\emph{Yíchǎn kěyǐ shì fángzi, qián, yě kěyǐ shì wénhuà hé chuántǒng.}\\
« L'héritage peut être une maison, de l'argent, mais aussi la culture et les traditions. »
\end{textebox}

Remarquez la structure très utile de la deuxième phrase : \zh{可以是……也可以是……} (\emph{kěyǐ shì... yě kěyǐ shì...}), « cela peut être... mais aussi... ». Vous pourrez la réutiliser dans vos propres productions.

\begin{attention}[Ne pas confondre : 遗产 et 遗传]
Un seul caractère change tout !
\begin{itemize}
\item \zh{遗产} (\emph{yíchǎn}) = l'\textbf{héritage} laissé par un défunt (biens, culture) ;
\item \zh{遗传} (\emph{yíchuán}) = l'\textbf{hérédité} génétique, ce qui se transmet biologiquement (la couleur des yeux, une maladie héréditaire).
\end{itemize}
Comparez : \zh{他继承了父亲的遗产。} (\emph{Tā jìchéng le fùqin de yíchǎn.}) « Il a hérité des biens de son père. » et \zh{这个病是遗传的。} (\emph{Zhège bìng shì yíchuán de.}) « Cette maladie est héréditaire. » Piège classique du francophone : les deux mots commencent par \zh{遗} et se prononcent presque pareil — vérifiez toujours le deuxième caractère : \zh{产} (bien) contre \zh{传} (transmettre).
\end{attention}

\courssub{Carte mentale du vocabulaire de l'héritage}

Pour mémoriser durablement un mot nouveau, la meilleure stratégie est de le relier à sa « famille » de mots. Voici la carte mentale du champ lexical de l'héritage ; recopiez-la dans votre cahier et enrichissez-la au fil de la leçon.

\begin{popfigure}
\begin{tikzpicture}[mindmap, grow cyclic, every node/.style={concept, align=center, font=\small},
root concept/.append style={concept color=popblue, font=\large\bfseries},
level 1/.append style={level distance=3.2cm, sibling angle=90},
level 1 concept/.append style={concept color=poporangeL}]
\node[root concept, align=center]{遗产\\ \emph{yíchǎn}\\ héritage}
  child { node[concept, align=center]{财产\\ \emph{cáichǎn}\\ biens, fortune} }
  child { node[concept, align=center]{遗嘱\\ \emph{yízhǔ}\\ testament} }
  child { node[concept, align=center]{继承\\ \emph{jìchéng}\\ hériter} }
  child { node[concept, align=center]{文化\\ \emph{wénhuà}\\ culture} };
\end{tikzpicture}
\legende{Carte mentale du champ lexical de l'héritage : le mot-clé 遗产 au centre et ses quatre mots associés.}
\end{popfigure}

Lisons chaque branche :
\begin{itemize}
\item \zh{财产} (\emph{cáichǎn}) : les biens, la fortune — c'est souvent le contenu matériel de l'héritage ;
\item \zh{遗嘱} (\emph{yízhǔ}) : le testament, le document où l'on écrit ses dernières volontés ;
\item \zh{继承} (\emph{jìchéng}) : hériter (verbe) — \zh{继承遗产}, « hériter d'un héritage » ;
\item \zh{文化} (\emph{wénhuà}) : la culture — avec \zh{遗产}, elle forme \zh{文化遗产}, le patrimoine culturel.
\end{itemize}

\courssub{Première écoute du texte : les oreilles d'abord}

\begin{experience}[Écoute sans support écrit — 3 minutes]
Fermez vos cahiers. L'enseignant lit le texte support à voix haute, deux fois, à vitesse naturelle. Votre seule tâche : \textbf{écouter les sons}, sans chercher à tout comprendre.
\begin{enumerate}
\item Première écoute : combien de fois entendez-vous le mot \zh{遗产} (\emph{yíchǎn}) ? Levez le doigt à chaque occurrence.
\item Deuxième écoute : notez en pinyin (même approximatif) deux ou trois mots que vous reconnaissez : \zh{jiā} (famille), \zh{fùqin} (père), \zh{fángzi} (maison)...
\end{enumerate}
But de l'activité : habituer l'oreille aux tons avant de voir les caractères. En chinois, entendre la mélodie de la phrase est la première étape de la compréhension.
\end{experience}

\courssub{Deuxième lecture : le texte avec caractères et pinyin}

Voici maintenant le texte support de la leçon. Lisez-le d'abord en silence, en soulignant les mots que vous connaissez déjà (\zh{家}, \zh{爸爸}, \zh{房子}, \zh{文化}...). La traduction guidée viendra ensuite — ne la lisez pas tout de suite, faites d'abord vos hypothèses !

\begin{textebox}[爷爷的遗产 — L'héritage de grand-père]
我的爷爷去年去世了。他给我们家留下了很多遗产。\\
\emph{Wǒ de yéye qùnián qùshì le. Tā gěi wǒmen jiā liúxià le hěn duō yíchǎn.}\\
« Mon grand-père est mort l'année dernière. Il a laissé beaucoup d'héritage à notre famille. »\\[4pt]
爷爷有一份遗嘱。遗嘱里说：房子给我的爸爸，钱分给三个孩子。\\
\emph{Yéye yǒu yí fèn yízhǔ. Yízhǔ lǐ shuō: fángzi gěi wǒ de bàba, qián fēn gěi sān ge háizi.}\\
« Grand-père avait un testament. Le testament disait : la maison pour mon père, l'argent partagé entre les trois enfants. »\\[4pt]
我的姑姑也继承了遗产。在中国，女儿和儿子有一样的权利。\\
\emph{Wǒ de gūgu yě jìchéng le yíchǎn. Zài Zhōngguó, nǚ'ér hé érzi yǒu yíyàng de quánlì.}\\
« Ma tante (sœur du père) a aussi hérité. En Chine, la fille et le fils ont les mêmes droits. »\\[4pt]
可是爷爷留下的最重要的东西不是钱，也不是房子。\\
\emph{Kěshì yéye liúxià de zuì zhòngyào de dōngxi bú shì qián, yě bú shì fángzi.}\\
« Mais la chose la plus importante que grand-père a laissée, ce n'est ni l'argent ni la maison. »\\[4pt]
他教我们写汉字，给我们讲中国的历史和故事。\\
\emph{Tā jiāo wǒmen xiě Hànzì, gěi wǒmen jiǎng Zhōngguó de lìshǐ hé gùshi.}\\
« Il nous a appris à écrire les caractères chinois, il nous racontait l'histoire et les contes de Chine. »\\[4pt]
这些文化遗产，比钱更重要。我们永远不会忘记爷爷。\\
\emph{Zhèxiē wénhuà yíchǎn, bǐ qián gèng zhòngyào. Wǒmen yǒngyuǎn bú huì wàngjì yéye.}\\
« Ce patrimoine culturel est plus important que l'argent. Nous n'oublierons jamais grand-père. »
\end{textebox}

\textbf{Repérage des mots connus.} Avant toute traduction, entourez dans le texte :
\begin{itemize}
\item les mots de la famille : \zh{爷爷} (grand-père paternel), \zh{爸爸} (papa), \zh{姑姑} (tante, sœur du père), \zh{女儿} (fille), \zh{儿子} (fils) ;
\item les mots de la vie quotidienne : \zh{房子} (maison), \zh{钱} (argent), \zh{东西} (chose) ;
\item les mots nouveaux de la carte mentale : \zh{遗产}, \zh{遗嘱}, \zh{继承}, \zh{文化}.
\end{itemize}

\courssub{Hypothèses de sens : deviner avant de traduire}

\begin{methode}[Comprendre un texte chinois en quatre étapes]
\begin{enumerate}
\item \textbf{Écouter} la mélodie globale (tons, mots répétés) sans support écrit.
\item \textbf{Repérer} les mots connus dans le texte écrit : ils donnent le thème (ici : famille + maison + argent → un partage).
\item \textbf{Formuler des hypothèses} : de quoi parle le texte ? Qui ? Quoi ? Ici : un petit-fils raconte ce que son grand-père a laissé à la famille.
\item \textbf{Vérifier} avec la traduction guidée et le glossaire, puis corriger ses hypothèses.
\end{enumerate}
Ne jamais traduire mot à mot dès le départ : le sens global d'abord, les détails ensuite.
\end{methode}

\begin{exoresolu}[Hypothèses de sens global]
Énoncé : sans regarder la traduction, répondez en français à ces trois questions à partir des mots que vous avez repérés. 1) De qui parle le texte ? 2) Quels biens matériels sont mentionnés ? 3) Le texte parle-t-il seulement d'argent ?
\tcblower
Solution détaillée :
\begin{enumerate}
\item Le mot \zh{爷爷} (\emph{yéye}, grand-père) revient dans presque chaque phrase, ainsi que \zh{我们} (nous) et \zh{家} (famille) : le texte parle du \textbf{grand-père du narrateur} et de sa famille.
\item On reconnaît \zh{房子} (\emph{fángzi}, maison) et \zh{钱} (\emph{qián}, argent) : les biens matériels sont \textbf{la maison et l'argent}.
\item Le mot \zh{文化} (\emph{wénhuà}, culture) et la structure \zh{不是……也不是……} (« ce n'est ni... ni... ») montrent que le texte parle \textbf{aussi d'autre chose} : un héritage culturel (les caractères, l'histoire, les histoires racontées).
\end{enumerate}
Hypothèse globale validée : le texte raconte le partage de l'héritage d'un grand-père chinois, et montre que l'héritage culturel compte plus que l'argent.
\end{exoresolu}

\courssub{Tableau récapitulatif : les premiers mots du thème}

Voici les mots essentiels rencontrés dans cette première approche. Le glossaire complet sera étudié dans la section suivante.

\begin{center}
\begin{tabularx}{\linewidth}{|X|X|X|X|}
\hline
\rowcolor{popblueL} Caractères & Pinyin & Nature & Traduction \\
\hline
遗产 & yíchǎn & nom & héritage \\
\hline
遗嘱 & yízhǔ & nom & testament \\
\hline
继承 & jìchéng & verbe & hériter \\
\hline
财产 & cáichǎn & nom & biens, fortune \\
\hline
文化 & wénhuà & nom & culture \\
\hline
留下 & liúxià & verbe & laisser (derrière soi) \\
\hline
去世 & qùshì & verbe & décéder \\
\hline
权利 & quánlì & nom & droit (d'une personne) \\
\hline
\end{tabularx}
\end{center}

\begin{aretenir}[L'essentiel de cette mise en route]
\begin{itemize}
\item \zh{遗产} (\emph{yíchǎn}) = \zh{遗} (laisser) + \zh{产} (bien) : l'héritage, matériel ou culturel.
\item Sa famille de mots : \zh{遗嘱} (testament), \zh{继承} (hériter), \zh{财产} (biens), \zh{文化遗产} (patrimoine culturel).
\item Piège : \zh{遗产} (héritage) ≠ \zh{遗传} (\emph{yíchuán}, hérédité génétique).
\item Méthode de lecture : écouter → repérer les mots connus → faire des hypothèses → vérifier avec la traduction.
\end{itemize}
\end{aretenir}

\begin{savaistu}[Filles et fils : les mêmes droits en Chine]
En Chine, la loi sur l'héritage, la \zh{继承法} (\emph{jìchéngfǎ}), entrée en vigueur en 1985, donne \textbf{les mêmes droits aux filles et aux fils} : \zh{女儿和儿子有一样的权利} (\emph{nǚ'ér hé érzi yǒu yíyàng de quánlì}), « la fille et le fils ont les mêmes droits ». C'est ce que le texte rappelle quand la tante (\zh{姑姑}) hérite elle aussi. Au Cameroun, les pratiques successorales varient selon les régions et les traditions, et le droit moderne coexiste avec les coutumes : un beau sujet de débat interculturel pour la section orale ! Autre fait intéressant : la Chine compte de nombreux sites classés au patrimoine mondial — on parle alors de \zh{世界文化遗产} (\emph{shìjiè wénhuà yíchǎn}), le « patrimoine culturel mondial ».
\end{savaistu}

\begin{exercice}[À faire avant la prochaine section]
\begin{enumerate}
\item Recopiez la carte mentale de \zh{遗产} et ajoutez deux branches de votre choix (par exemple \zh{房子} « maison » et \zh{传统} « tradition »).
\item Traduisez en français : \zh{爷爷留下了一份遗嘱。} (\emph{Yéye liúxià le yí fèn yízhǔ.})
\item Vrai ou faux : \zh{遗传} signifie « héritage ». Justifiez.
\end{enumerate}
\end{exercice}

\begin{corrige}[Exercice de mise en route]
\begin{enumerate}
\item Carte mentale : vérifier que les deux branches ajoutées sont écrites en caractères corrects avec le pinyin et la traduction.
\item « Grand-père a laissé un testament. » Noter le classificateur \zh{份} (\emph{fèn}) utilisé pour les documents.
\item \textbf{Faux.} \zh{遗传} (\emph{yíchuán}) signifie « hérédité (génétique) » ; « héritage » se dit \zh{遗产} (\emph{yíchǎn}). Le deuxième caractère fait toute la différence : \zh{产} (bien) contre \zh{传} (transmettre).
\end{enumerate}
\end{corrige}

\coursec{Glossaire complet du thème : l'héritage}

\courssub{Tableau de vocabulaire officiel (niveau HSK 3-4)}

Ce tableau regroupe \textbf{tous} les mots lexicaux que vous devez maîtriser pour cette leçon (compréhension, expression écrite, traduction). Apprenez-les par cœur : caractères, pinyin (tons exacts), nature grammaticale et traduction.

\begin{center}
\begin{tabularx}{\linewidth}{|X|X|X|X|}
\hline
\rowcolor{popblueL} Caractères & Pinyin & Nature & Traduction \\
\hline
遗产 & yíchǎn & nom & héritage (biens, culture) \\
\hline
遗嘱 & yízhǔ & nom & testament \\
\hline
继承 & jìchéng & verbe & hériter, succéder à \\
\hline
财产 & cáichǎn & nom & biens, fortune, patrimoine matériel \\
\hline
文化 & wénhuà & nom & culture \\
\hline
文化遗产 & wénhuà yíchǎn & nom & patrimoine culturel \\
\hline
留下 & liúxià & verbe & laisser, laisser derrière soi \\
\hline
去世 & qùshì & verbe & décéder, mourir (euphémisme poli) \\
\hline
权利 & quánlì & nom & droit, privilège \\
\hline
分配 & fēnpèi & verbe & répartir, distribuer, partager \\
\hline
公平 & gōngpíng & adj. & juste, équitable, impartial \\
\hline
房子 & fángzi & nom & maison, bâtiment \\
\hline
钱 & qián & nom & argent, monnaie \\
\hline
东西 & dōngxi & nom & chose, objet, affaire \\
\hline
重要 & zhòngyào & adj. & important \\
\hline
历史 & lìshǐ & nom & histoire (discipline, passé) \\
\hline
故事 & gùshi & nom & histoire, conte, récit \\
\hline
汉字 & Hànzì & nom & caractère chinois \\
\hline
教 & jiāo & verbe & enseigner, apprendre à \\
\hline
讲 & jiǎng & verbe & raconter, expliquer, parler de \\
\hline
永远 & yǒngyuǎn & adv. & toujours, éternellement, jamais (avec négation) \\
\hline
忘记 & wàngjì & verbe & oublier \\
\hline
爷爷 & yéye & nom & grand-père paternel \\
\hline
姑姑 & gūgu & nom & tante (sœur du père) \\
\hline
女儿 & nǚ'ér & nom & fille \\
\hline
儿子 & érzi & nom & fils \\
\hline
一样 & yíyàng & adj./adv. & même, identique, pareil \\
\hline
份 & fèn & classif. & classificateur pour documents, journaux, cadeaux \\
\hline
\end{tabularx}
\end{center}

\courssub{Collocations incontournables (associations de mots)}

En chinois, les mots ne vivent pas seuls. Mémorisez ces \cle{collocations} (associations fixes) pour parler naturellement de l'héritage.

\begin{center}
\begin{tabularx}{\linewidth}{|X|X|X|}
\hline
\rowcolor{popblueL} Collocation (Caractères + Pinyin) & Structure & Traduction \\
\hline
\zh{留下遗产} (\emph{liúxià yíchǎn}) & Verbe + Objet & laisser un héritage \\
\hline
\zh{立遗嘱 / 写遗嘱} (\emph{lì yízhǔ / xiě yízhǔ}) & Verbe + Objet & faire un testament / écrire un testament \\
\hline
\zh{继承遗产} (\emph{jìchéng yíchǎn}) & Verbe + Objet & hériter (d'un héritage) \\
\hline
\zh{分配财产} (\emph{fēnpèi cáichǎn}) & Verbe + Objet & partager les biens / répartir la fortune \\
\hline
\zh{享有权利} (\emph{xiǎngyǒu quánlì}) & Verbe + Objet & jouir d'un droit / bénéficier d'un droit \\
\hline
\zh{保护文化遗产} (\emph{bǎohù wénhuà yíchǎn}) & Verbe + Objet & protéger le patrimoine culturel \\
\hline
\zh{传承文化} (\emph{chuánchéng wénhuà}) & Verbe + Objet & transmettre / perpétuer la culture \\
\hline
\end{tabularx}
\end{center}

\begin{attention}[Pièges de traduction : 留下 vs 留]
\begin{itemize}
\item \zh{留下} (\emph{liúxià}) : verbe composé \textbf{résultatif} = « laisser (derrière soi), rester ». Complément d'objet obligatoire. Ex: \zh{他留下了一本书。} (Il a laissé un livre.)
\item \zh{留} (\emph{liú}) seul : souvent dans \zh{留学} (\emph{liúxué}, étudier à l'étranger) ou \zh{保留} (\emph{bǎoliú}, conserver/réserver). Ne dites pas \zh{他留了一本书} pour « il a laissé un livre » (cela sonne comme « il a gardé/réservé un livre »).
\item \zh{去世} (\emph{qùshì}) : verbe \textbf{intransitif}, registre soutenu/poli. Ne pas dire \zh{死} (\emph{sǐ}, mourir — cru) dans un contexte familial respectueux. Pas de COD : \zh{爷爷去世了} (Grand-père est décédé), PAS \zh{爷爷去世了奶奶}.
\end{itemize}
\end{attention}

\begin{exercice}[Entraînement vocabulaire — 10 min]
\begin{enumerate}
\item Associez le verbe à son objet naturel : \zh{立} / \zh{继承} / \zh{分配} / \zh{保护} — \zh{遗嘱} / \zh{遗产} / \zh{财产} / \zh{文化遗产}.
\item Traduisez en chinois (caractères + pinyin) :
      a) Grand-père a laissé une maison. \\
      b) Les enfants partagent l'argent équitablement. (\emph{équitablement} = \zh{公平地} \emph{gōngpíng de}) \\
      c) Nous devons protéger le patrimoine culturel.
\item Complétez avec le bon mot : \zh{遗产 / 遗嘱 / 遗传 / 财产}
      1. 这个病是\_\_\_\_\_\_的。 (Cette maladie est héréditaire.)
      2. 他\_\_\_\_\_\_了父亲的\_\_\_\_\_\_. (Il a hérité des biens de son père.)
      3. 请写一份\_\_\_\_\_\_. (Veuillez écrire un testament.)
\end{enumerate}
\end{exercice}

\begin{corrige}[Entraînement vocabulaire]
\begin{enumerate}
\item \zh{立遗嘱} (faire un testament), \zh{继承遗产} (hériter), \zh{分配财产} (partager les biens), \zh{保护文化遗产} (protéger le patrimoine).
\item a) \zh{爷爷留下了一所房子。} (\emph{Yéye liúxià le yí suǒ fángzi.}) — classif. \zh{所} pour maisons. \\
      b) \zh{孩子们公平地分钱。} (\emph{Háizi men gōngpíng de fēn qián.}) ou \zh{公平地分配钱。} \\
      c) \zh{我们必须保护文化遗产。} (\emph{Wǒmen bìxū bǎohù wénhuà yíchǎn.})
\item 1. \zh{遗传} (\emph{yíchuán}) 2. \zh{继承 / 财产} (\emph{jìchéng / cáichǎn}) 3. \zh{遗嘱} (\emph{yízhǔ})
\end{enumerate}
\end{corrige}

\coursec{Clés (radicaux) et ordre des traits}

\courssub{Pourquoi analyser les radicaux ?}

Chaque caractère chinois est construit autour d'une \cle{clé} (radical, \zh{部首} \emph{bùshǒu}) qui donne souvent une indication sémantique (le sens) ou phonétique. Connaître les clés permet de : (1) retrouver un caractère dans un dictionnaire papier, (2) deviner le sens d'un caractère inconnu, (3) mémoriser l'écriture par logique et non par cœur.

\courssub{Analyse de 7 caractères clés de la leçon}

\begin{center}
\begin{tabularx}{\linewidth}{|c|X|X|X|X|}
\hline
\rowcolor{popblueL} Car. & Pinyin & Clé (Radical) & Sens de la clé & Ordre des traits (règles principales) \\
\hline
\zh{遗} & yí & \zh{辶} (\emph{chuò}) « marche » & Action de laisser derrière soi en partant & 1. Haut \zh{艹} (herbe) : horiz., vert., horiz. 2. \zh{贵} : haut→bas, gauche→droite. 3. \zh{辶} : point, horizontale courbe, trait final descendant. (17 traits) \\
\hline
\zh{产} & chǎn & \zh{生} (\emph{shēng}) « vie, naître » & Ce qui est produit, né, engendré & 1. \zh{立} (debout) : horiz., vert., horiz., vert. central. 2. \zh{二} : deux horizontales. 3. \zh{厂} (cliff) : horiz., verticale cassée. (6 traits) \\
\hline
\zh{嘱} & zhǔ & \zh{口} (\emph{kǒu}) « bouche » & Action de la parole, ordre verbal & 1. \zh{口} : vert., horiz. cassée, horiz. (boîte ouverte). 2. \zh{主} (maître) : point, horiz., vert. 3. \zh{寸} (pouce) : horiz., horiz., vert., horiz. (15 traits) \\
\hline
\zh{继} & jì & \zh{纟} (\emph{sī}) « fil de soie » & Continuer, relier comme un fil & 1. \zh{纟} : diagonale montante, diagonale descendante, point. 2. \zh{米} (riz) : point, horiz., vert., diagonales, horiz. (10 traits) \\
\hline
\zh{承} & chéng & \zh{受} (\emph{shòu}) / \zh{手} (\emph{shǒu}) « main » & Recevoir, porter, soutenir & 1. \zh{禾} (riz) : diagonale, horiz., vert., diagonale. 2. \zh{受} : \zh{爪} (griffe) + \zh{又} (main droite). (8 traits) \\
\hline
\zh{权} & quán & \zh{木} (\emph{mù}) « bois, arbre » & Poids en bois pour balance → autorité & 1. \zh{木} : horiz., vert., diagonale g., diagonale d. 2. \zh{泉} (source) : \zh{白} + \zh{水} (eau). (22 traits — complexe, reconnaître d'abord) \\
\hline
\zh{利} & lì & \zh{刂} (\emph{dāo}) « couteau » (à droite) & Tranchant, avantage, profit & 1. \zh{禾} (riz). 2. \zh{刂} : verticale, verticale crochue. (7 traits) \\
\hline
\end{tabularx}
\end{center}

\begin{methode}[Écrire un caractère correctement : 3 règles d'or]
\begin{enumerate}
\item \textbf{Haut → Bas} (\zh{上到下}) : \zh{产} s'écrit le haut (\zh{立}) puis le bas (\zh{厂}).
\item \textbf{Gauche → Droite} (\zh{左到右}) : \zh{嘱} s'écrit \zh{口} (gauche) puis \zh{主} + \zh{寸} (droite).
\item \textbf{Horizontal avant Vertical} (\zh{横竖}) : quand un trait horizontal croise un vertical, l'horizontale d'abord (sauf si le vertical est le trait principal central, ex: \zh{中}).
\end{enumerate}
Entraînez-vous sur papier quadrillé (grille \zh{田字格} \emph{tiánzìgé}) : un caractère par case, centré, proportions équilibrées.
\end{methode}

\begin{experience}[Atelier écriture — 10 minutes]
\begin{enumerate}
\item L'enseignant montre l'ordre des traits de \zh{遗} et \zh{产} au tableau (ou vidéo).
\item Les élèves écrivent chaque caractère 5 fois en disant le nom des traits à voix haute (\emph{héng, shù, piě, nà...}).
\item Échange de cahiers : vérification mutuelle de l'ordre et de la proportion (clé \zh{辶} bien en bas à gauche, pas trop grande).
\item Défi : écrire \zh{遗产} trois fois de suite sans modèle, en moins d'une minute.
\end{enumerate}
\end{experience}

\begin{aretenir}[Radicaux à retenir pour cette leçon]
\begin{itemize}
\item \zh{辶} (\emph{chuò}) : mouvement, départ, laisser derrière → \zh{遗}, \zh{过}, \zh{进}, \zh{近}.
\item \zh{生} (\emph{shēng}) : vie, naissance, production → \zh{产}, \zh{生}, \zh{姓}, \zh{牲}.
\item \zh{口} (\emph{kǒu}) : bouche, parole → \zh{嘱}, \zh{吃}, \zh{说}, \zh{问}.
\item \zh{纟} (\emph{sī}) : fil, soie, suite → \zh{继}, \zh{续}, \zh{约}, \zh{级}.
\item \zh{刂} (\emph{dāo}) : couteau, action tranchante → \zh{利}, \zh{刀}, \zh{切}, \zh{别}.
\end{itemize}
\end{aretenir}

\coursec{Compréhension du texte et collocations}

\courssub{Questions de compréhension détaillée (niveau examen)}

Relisez le texte \zh{爷爷的遗产} et répondez par écrit en phrases complètes (français ou chinois selon la consigne).

\begin{exercice}[Compréhension écrite — 15 min]
\textbf{Partie A : Questions globales (réponse en français)}
\begin{enumerate}
\item Qui est le narrateur du texte ? Justifiez par un mot du texte.
\item Combien d'enfants héritent de l'argent ? Qui sont-ils ?
\item Quelle est la nationalité implicite de la famille ? Citez l'indice textuel.
\item Selon le narrateur, quel est l'héritage \emph{le plus important} ? Pourquoi ?
\end{enumerate}

\textbf{Partie B : Questions de détail (réponse en chinois — caractères + pinyin)}
\begin{enumerate}
\item 爷爷什么时候去世的？ (\emph{Yéye shénme shíhou qùshì de?})
\item 遗嘱里怎么分配房子和钱？ (\emph{Yízhǔ lǐ zěnme fēnpèi fángzi hé qián?})
\item 姑姑有没有继承遗产？为什么？ (\emph{Gūgu yǒu méiyǒu jìchéng yíchǎn? Wèishénme?})
\item 爷爷教了孩子们什么？ (\emph{Yéye jiāo le háizi men shénme?})
\item 翻译这句话：\zh{这些文化遗产，比钱更重要。}
\end{enumerate}
\end{exercice}

\begin{corrige}[Compréhension écrite]
\textbf{Partie A}
\begin{enumerate}
\item Le narrateur est un \textbf{petit-fils / petite-fille} (mot-clé : \zh{我的爷爷} « mon grand-père »).
\item \textbf{Trois enfants} héritent de l'argent (mot-clé : \zh{三个孩子}).
\item La famille est \textbf{chinoise} (indice : \zh{在中国} « en Chine », \zh{写汉字} « écrire des caractères chinois », \zh{中国的历史} « l'histoire de Chine »).
\item L'héritage le plus important est le \textbf{patrimoine culturel} (\zh{文化遗产}) : les caractères chinois, l'histoire et les contes de Chine. Le texte dit explicitement : \zh{这些文化遗产，比钱更重要}.
\end{enumerate}

\textbf{Partie B}
\begin{enumerate}
\item \zh{爷爷去年去世了。} (\emph{Yéye qùnián qùshì le.})
\item \zh{房子给爸爸，钱分给三个孩子。} (\emph{Fángzi gěi bàba, qián fēn gěi sān ge háizi.})
\item \zh{姑姑也继承了遗产。因为在中国，女儿和儿子有一样的权利。} (\emph{Gūgu yě jìchéng le yíchǎn. Yīnwèi zài Zhōngguó, nǚ'ér hé érzi yǒu yíyàng de quánlì.})
\item \zh{爷爷教我们写汉字，讲中国的历史和故事。} (\emph{Yéye jiāo wǒmen xiě Hànzì, jiǎng Zhōngguó de lìshǐ hé gùshi.})
\item « Ce patrimoine culturel est plus important que l'argent. »
\end{enumerate}
\end{corrige}

\courssub{Travail sur les structures utiles du texte}

Le texte contient trois structures grammaticales de niveau HSK 3-4, réutilisables à l'oral et à l'écrit.

\begin{propriete}[Structure 1 : 把 (bǎ) — Disposition de l'objet (révision Terminale)]
Bien que le texte n'utilise pas \zh{把} explicitement, la phrase \zh{房子给我的爸爸，钱分给三个孩子} suit la logique \zh{把} : \textbf{Sujet + 把 + Objet + Verbe + Complément de résultat/direction}.
\begin{itemize}
\item \zh{爷爷把房子给了爸爸。} (\emph{Yéye bǎ fángzi gěi le bàba.}) — Grand-père a donné la maison à papa.
\item \zh{他把钱分给了三个孩子。} (\emph{Tā bǎ qián fēn gěi le sān ge háizi.}) — Il a partagé l'argent entre les trois enfants.
\end{itemize}
\emph{Rappel :} \zh{把} met l'accent sur \textbf{ce qui arrive à l'objet} (manipulation, changement de lieu/propriétaire). L'objet doit être \textbf{connu/spécifique} (défini).
\end{propriete}

\begin{propriete}[Structure 2 : 比 (bǐ) — Comparatif (niveau HSK 3)]
\zh{这些文化遗产，比钱更重要。} (\emph{Zhèxiē wénhuà yíchǎn, bǐ qián gèng zhòngyào.})
\textbf{Schéma :} \textbf{A + 比 + B + Adjectif (souvent + 更/还/得多)}.
\begin{itemize}
\item A = \zh{这些文化遗产} (sujet comparé), B = \zh{钱} (référence), Adj = \zh{重要}.
\item \zh{更} (\emph{gèng}) = « encore plus », renforce le comparatif.
\end{itemize}
\emph{Variantes :} \zh{比钱重要} (plus important que l'argent) ; \zh{比钱重要得多} (beaucoup plus important que l'argent).
\end{propriete}

\begin{propriete}[Structure 3 : 不是……也不是…… (bú shì... yě bú shì...) — Double négation pour insister sur un tiers exclu]
\zh{最重要的东西不是钱，也不是房子。} (\emph{Zuì zhòngyào de dōngxi bú shì qián, yě bú shì fángzi.})
\textbf{Emploi :} Énumérer ce que ce n'est \textbf{PAS}, pour préparer l'annonce de ce que c'est \textbf{VRAIMENT} (souvent introduit par \zh{而是} \emph{érshì} « mais c'est »).
\begin{itemize}
\item Ex: \zh{他不是中国人，也不是日本人，而是韩国人。} (Il n'est ni Chinois ni Japonais, mais Coréen.)
\end{itemize}
\end{propriete}

\begin{exoresolu}[Transformation de phrases]
Énoncé : Réécrivez les phrases suivantes en utilisant la structure indiquée.
\begin{enumerate}
\item \zh{爷爷留下了房子。} → Utilisez \zh{把} (\emph{bǎ}).
\item \zh{文化遗产重要。钱不重要。} → Combinez avec \zh{比} (\emph{bǐ}) + \zh{更} (\emph{gèng}).
\item \zh{遗产不是钱。} → Ajoutez \zh{也不是房子} (\emph{yě bú shì fángzi}) et concluez par \zh{而是文化} (\emph{érshì wénhuà}).
\end{enumerate}
\tcblower
Solution détaillée :
\begin{enumerate}
\item \zh{爷爷把房子留下了。} (\emph{Yéye bǎ fángzi liúxià le.}) — Attention : \zh{留下} verbe résultatif, \zh{了} en fin de phrase.
\item \zh{文化遗产比钱更重要。} (\emph{Wénhuà yíchǎn bǐ qián gèng zhòngyào.})
\item \zh{遗产不是钱，也不是房子，而是文化。} (\emph{Yíchǎn bú shì qián, yě bú shì fángzi, érshì wénhuà.})
\end{enumerate}
\end{exoresolu}

\coursec{Production guidée et réinvestissement}

\courssub{Expression écrite : « L'héritage de ma famille »}

\begin{methode}[Rédiger un paragraphe narratif en chinois (niveau Terminale)]
\begin{enumerate}
\item \textbf{Préparez le vocabulaire :} listez 5-6 mots clés de la leçon que vous allez utiliser (\zh{遗产}, \zh{留下}, \zh{继承}, \zh{文化}, \zh{比}, \zh{不是...而是...}).
\item \textbf{Structurez :} 1. Contexte (qui, quand) → 2. Biens matériels (\zh{把} structure) → 3. Héritage culturel (\zh{比} structure) → 4. Conclusion sentimentale (\zh{永远不忘记}).
\item \textbf{Rédigez en caractères} directement (pas de pinyin d'abord). Vérifiez l'ordre des mots : Sujet - Temps - Lieu - Verbe - Objet.
\item \textbf{Relisez :} tons (marquez-les au crayon), classificateurs, particules (\zh{了}, \zh{的}, \zh{地}).
\end{enumerate}
\end{methode}

\begin{exercice}[Production écrite guidée — 20 min]
Rédigez un court texte (80-100 caractères chinois) sur le thème : \textbf{« L'héritage le plus précieux de ma famille »} (\zh{我家最宝贵的遗产}).
\begin{itemize}
\item \textbf{Contraintes obligatoires :}
      \begin{enumerate}
      \item Utilisez \zh{把} au moins une fois (pour un bien matériel).
      \item Utilisez \zh{比...更...} pour comparer matériel vs culturel.
      \item Utilisez \zh{不是...而是...} pour opposer.
      \item Intégrez au moins 5 mots du glossaire de la leçon.
      \item Terminez par \zh{我们永远不会忘记……} (\emph{Wǒmen yǒngyuǎn bú huì wàngjì...}).
      \end{enumerate}
\item \textbf{Amorce (à compléter) :}
\zh{我的爷爷/奶奶去年去世了。他/她把\_\_\_\_\_\_给了\_\_\_\_\_\_。可是，\_\_\_\_\_\_比\_\_\_\_\_\_更重要。因为\_\_\_\_\_\_不是\_\_\_\_\_\_，而是\_\_\_\_\_\_。我们永远不会忘记\_\_\_\_\_\_。}
\end{itemize}
\end{exercice}

\begin{exemplebox}[Modèle de production attendue (niveau excellence)]
我的奶奶去年去世了。她把老房子给了我的爸爸。可是，奶奶教我们的做菜方法，比房子更重要。因为遗产不是钱，也不是房子，而是爱和传统。我们永远不会忘记奶奶的味道。
\emph{Wǒ de nǎinai qùnián qùshì le. Tā bǎ lǎo fángzi gěi le wǒ de bàba. Kěshì, nǎinai jiāo wǒmen de zuò cài fāngfǎ, bǐ fángzi gèng zhòngyào. Yīnwèi yíchǎn bú shì qián, yě bú shì fángzi, érshì ài hé chuántǒng. Wǒmen yǒngyuǎn bú huì wàngjì nǎinai de wèidào.}
« Ma grand-mère est décédée l'an dernier. Elle a donné la vieille maison à mon père. Mais la façon de cuisiner que grand-mère nous a apprise est plus importante que la maison. Car l'héritage n'est ni l'argent ni la maison, c'est l'amour et la tradition. Nous n'oublierons jamais le goût de grand-mère. »
\end{exemplebox}

\courssub{Expression orale : Débat et jeu de rôle}

\begin{experience}[Débat : L'argent ou la culture ? — 15 min]
\textbf{Sujet :} \zh{遗产里，钱重要，还是文化重要？} (\emph{Yíchǎn lǐ, qián zhòngyào, háishi wénhuà zhòngyào?})
\textbf{Organisation :} Deux groupes (Groupe A : \zh{钱重要} / Groupe B : \zh{文化重要}). Préparation 5 min (arguments en chinois sur fiche), débat 10 min.
\textbf{Phrases-outils (à afficher) :}
\begin{itemize}
\item \zh{我认为……因为……} (\emph{Wǒ rènwéi... yīnwèi...}) — Je pense que... parce que...
\item \zh{没有钱，生活很难。} (\emph{Méiyǒu qián, shēnghuó hěn nán.}) — Sans argent, la vie est dure.
\item \zh{文化是根，钱是叶。} (\emph{Wénhuà shì gēn, qián shì yè.}) — La culture est la racine, l'argent les feuilles (proverbe adapté).
\item \zh{两者都重要，但……} (\emph{Liǎng zhě dōu zhòngyào, dàn...}) — Les deux sont importants, mais...
\end{itemize}
\textbf{Évaluation :} Fluidité, justesse des structures (\zh{比}, \zh{因为/所以}), vocabulaire de la leçon, prononciation des tons.
\end{experience}

\begin{experience}[Jeu de rôle : Chez le notaire — 10 min]
\textbf{Scénario :} Trois héritiers (Grand frère, Petite sœur, Tante) et un Notaire (\zh{公证员} \emph{gōngzhèngyuán}) se réunissent pour lire le testament (\zh{宣读遗嘱} \emph{xuāndú yízhǔ}).
\textbf{Rôles :}
\begin{itemize}
\item Notaire : lit le testament, explique la loi (\zh{根据法律……} \emph{gēnjù fǎlǜ...}).
\item Héritiers : posent des questions, négocient (\zh{我想要……} \emph{Wǒ xiǎng yào...}, \zh{这不公平！} \emph{Zhè bù gōngpíng!}).
\end{itemize}
\textbf{Contraintes linguistiques :} Utiliser \zh{分配} (\emph{fēnpèi}), \zh{公平} (\emph{gōngpíng}), \zh{继承} (\emph{jìchéng}), \zh{份} (\emph{fèn}), \zh{权利} (\emph{quánlì}).
\end{experience}

\courssub{Traduction Thème / Version (entraînement examen)}

\begin{exercice}[Traduction — Version (Chinois → Français)]
Traduisez en français soigné :
\begin{enumerate}
\item 爷爷生前立了一份遗嘱，把财产公平地分给了三个孩子。
\item 这些文化遗产包括汉字、历史和传统故事，比金钱更有价值。
\item 女儿和儿子享有同样的继承权利，这是法律规定的。
\end{enumerate}
\end{exercice}

\begin{exercice}[Traduction — Thème (Français → Chinois)]
Traduisez en chinois (caractères + pinyin) :
\begin{enumerate}
\item Grand-père a laissé un patrimoine culturel très important à la famille.
\item La tante a aussi hérité de la maison, car la loi protège les droits des filles.
\item Nous devons transmettre la culture chinoise, ne pas l'oublier.
\end{enumerate}
\end{exercice}

\begin{corrige}[Traduction Thème / Version]
\textbf{Version}
\begin{enumerate}
\item « De son vivant, grand-père a établi un testament et a réparti équitablement la fortune entre les trois enfants. »
\item « Ce patrimoine culturel comprend les caractères chinois, l'histoire et les contes traditionnels ; il a plus de valeur que l'argent. » (\emph{注意 : 更有价值} \emph{gèng yǒu jiàzhí}).
\item « Les filles et les fils jouissent des mêmes droits d'héritage, c'est ce que stipule la loi. »
\end{enumerate}

\textbf{Thème}
\begin{enumerate}
\item \zh{爷爷给家族留下了很重要的文化遗产。} (\emph{Yéye gěi jiāzú liúxià le hěn zhòngyào de wénhuà yíchǎn.})
\item \zh{姑姑也继承了房子，因为法律保护女儿的权利。} (\emph{Gūgu yě jìchéng le fángzi, yīnwèi fǎlǜ bǎohù nǚ'ér de quánlì.})
\item \zh{我们必须传承中国文化，永远不忘记。} (\emph{Wǒmen bìxū chuánchéng Zhōngguó wénhuà, yǒngyuǎn bù wàngjì.})
\end{enumerate}
\end{corrige}

\begin{aretenir}[Bilan de la leçon ZH05 — Ce qu'il faut maîtriser pour l'examen]
\begin{itemize}
\item \textbf{Vocabulaire :} 25 mots du glossaire (caractères, pinyin, tons, nature, traduction).
\item \textbf{Écriture :} 7 caractères clés (\zh{遗, 产, 嘱, 继, 承, 权, 利}) : clé, ordre des traits, composition.
\item \textbf{Grammaire :} 3 structures : \zh{把} (disposition), \zh{比...更...} (comparatif), \zh{不是...也不是...而是...} (double négation + correction).
\item \textbf{Compréhension :} Savoir répondre aux questions globales et détaillées sur le texte \zh{爷爷的遗产}.
\item \textbf{Production :} Rédiger un paragraphe (80-100 car.) sur l'héritage familial en utilisant les structures imposées.
\item \textbf{Culture :} Droit successoral chinois (égalité filles/fils), notion de \zh{文化遗产} (patrimoine culturel), comparaison respectueuse Cameroun/Chine.
\end{itemize}
\end{aretenir}

\begin{savaistu}[Le saviez-vous ? Le classificateur 份 (fèn)]
Dans le texte : \zh{一份遗嘱} (\emph{yí fèn yízhǔ}). \zh{份} s'utilise pour : les documents officiels (\zh{合同} contrat, \zh{报纸} journal, \zh{文件} document), les cadeaux (\zh{礼物}), les médicaments (\zh{药}), les repas (\zh{早餐}). Astuce mnémotechnique : \zh{份} = \zh{亻} (homme) + \zh{分} (partager) → « la part d'un homme » → une unité comptable pour des choses qu'on distribue/partage.
\end{savaistu}

\coursec{Texte support : lecture, pinyin, traduction}

Dans la section précédente, nous avons découvert le thème de l'héritage (\zh{遗产} \emph{yíchǎn}) à travers une situation de communication et quelques images mentales : une famille réunie, un testament, une maison transmise. Il est temps maintenant d'entrer dans le cœur de la leçon : la lecture méthodique du texte support. Nous allons d'abord lire le texte en entier, phrase par phrase, avec le pinyin et la traduction, puis nous apprendrons une méthode de compréhension réutilisable pour tout texte chinois.

\courssub{Lecture globale du texte}

\begin{textebox}[L'héritage de grand-père — 爷爷的遗产]
我的爷爷去年去世了。\\
\emph{Wǒ de yéye qùnián qùshì le.}\\
Mon grand-père est décédé l'année dernière.\\[4pt]
他给我们留下了很多遗产：一栋房子、一些钱和一块地。\\
\emph{Tā gěi wǒmen liúxià le hěn duō yíchǎn : yī dòng fángzi, yīxiē qián hé yī kuài dì.}\\
Il nous a laissé beaucoup de biens : une maison, un peu d'argent et un terrain.\\[4pt]
爷爷生前立了遗嘱，说房子留给爸爸，钱分给我和妹妹。\\
\emph{Yéye shēngqián lì le yízhǔ, shuō fángzi liú gěi bàba, qián fēn gěi wǒ hé mèimei.}\\
Grand-père a rédigé un testament de son vivant : la maison revient à papa, l'argent est partagé entre ma petite sœur et moi.\\[4pt]
奶奶说，最重要的遗产不是钱，而是爷爷教我们的中国书法。\\
\emph{Nǎinai shuō, zuì zhòngyào de yíchǎn bú shì qián, ér shì yéye jiāo wǒmen de Zhōngguó shūfǎ.}\\
Grand-mère dit que l'héritage le plus important n'est pas l'argent, mais la calligraphie chinoise que grand-père nous a enseignée.\\[4pt]
我觉得她说得对。\\
\emph{Wǒ juéde tā shuō de duì.}\\
Je pense qu'elle a raison.\\[4pt]
文化遗产比钱更宝贵。\\
\emph{Wénhuà yíchǎn bǐ qián gèng bǎoguì.}\\
Le patrimoine culturel est plus précieux que l'argent.\\[4pt]
我们要把爷爷的书法传下去。\\
\emph{Wǒmen yào bǎ yéye de shūfǎ chuán xiàqù.}\\
Nous devons transmettre la calligraphie de grand-père aux générations suivantes.
\end{textebox}

\begin{attention}[Registre de langue : 去世 ou 死 ?]
Pour annoncer la mort d'un proche, on utilise \cle{去世} (\emph{qùshì}, « décéder, s'en aller »), terme poli et respectueux. Le mot \zh{死} (\emph{sǐ}, « mourir ») est direct, voire brutal : dire \zh{我爷爷死了} choquerait un interlocuteur chinois, comme dire en français « mon grand-père a crevé ». Réservez \zh{死} aux animaux, aux plantes ou aux faits divers. Autres formules respectueuses : \zh{过世} (\emph{guòshì}) et \zh{离开了我们} (\emph{líkāi le wǒmen}, « il nous a quittés »).
\end{attention}

\courssub{Lecture phrase par phrase : pinyin et traduction}

Lisons maintenant le texte phrase par phrase. Pour chaque phrase : d'abord les caractères, puis le pinyin, puis une traduction \textbf{littérale} (mot à mot, pour comprendre la mécanique chinoise), puis une traduction \textbf{naturelle} (en bon français).

\textbf{Phrase 1.} \zh{我的爷爷去年去世了。}\\
\emph{Wǒ de yéye qùnián qùshì le.}\\
Littéral : « Mon grand-père, année dernière, décédé + \zh{了}. »\\
Naturel : « Mon grand-père est décédé l'année dernière. »\\
Observez : le complément de temps \zh{去年} se place \textbf{avant le verbe}, contrairement au français où « l'année dernière » vient à la fin.

\textbf{Phrase 2.} \zh{他给我们留下了很多遗产：一栋房子、一些钱和一块地。}\\
\emph{Tā gěi wǒmen liúxià le hěn duō yíchǎn : yī dòng fángzi, yīxiē qián hé yī kuài dì.}\\
Littéral : « Lui, à nous, laissé + beaucoup de patrimoine : un [classificateur \zh{栋}] maison, un peu d'argent et un [classificateur \zh{块}] terrain. »\\
Naturel : « Il nous a laissé beaucoup de biens : une maison, un peu d'argent et un terrain. »\\
Observez les classificateurs : \zh{栋} (\emph{dòng}) pour les bâtiments, \zh{块} (\emph{kuài}) pour les terrains et les objets en forme de bloc. En français, on dit simplement « une maison », « un terrain » : le chinois, lui, exige un classificateur entre le nombre et le nom.

\textbf{Phrase 3.} \zh{爷爷生前立了遗嘱，说房子留给爸爸，钱分给我和妹妹。}\\
\emph{Yéye shēngqián lì le yízhǔ, shuō fángzi liú gěi bàba, qián fēn gěi wǒ hé mèimei.}\\
Littéral : « Grand-père, durant sa vie, établi + testament, dit : maison laisser à papa, argent partager à moi et petite sœur. »\\
Naturel : « De son vivant, grand-père a rédigé un testament : la maison revient à papa, l'argent est partagé entre ma petite sœur et moi. »\\
Observez : \zh{生前} (\emph{shēngqián}) signifie « de son vivant » ; \zh{立遗嘱} (\emph{lì yízhǔ}) est la collocation exacte pour « rédiger un testament ».

\textbf{Phrase 4.} \zh{奶奶说，最重要的遗产不是钱，而是爷爷教我们的中国书法。}\\
\emph{Nǎinai shuō, zuì zhòngyào de yíchǎn bú shì qián, ér shì yéye jiāo wǒmen de Zhōngguó shūfǎ.}\\
Littéral : « Grand-mère dit : le plus important patrimoine n'est pas argent, mais est grand-père enseigner-nous + \zh{de} calligraphie chinoise. »\\
Naturel : « Grand-mère dit que l'héritage le plus important n'est pas l'argent, mais la calligraphie chinoise que grand-père nous a enseignée. »\\
Observez la structure \zh{不是……而是……} (\emph{bú shì... ér shì...}) : « ce n'est pas... mais c'est... ». Notez aussi la mutation de ton : \zh{不} se prononce \emph{bú} (2\up{e} ton) devant \zh{是} (\emph{shì}, 4\up{e} ton).

\textbf{Phrase 5.} \zh{我觉得她说得对。}\\
\emph{Wǒ juéde tā shuō de duì.}\\
Littéral : « Moi trouver, elle parler + \zh{de} + juste. »\\
Naturel : « Je pense qu'elle a raison. »\\
Observez le complément de degré avec \zh{得} : \zh{说得对} = « parler de manière juste », c'est-à-dire « avoir raison ».

\textbf{Phrase 6.} \zh{文化遗产比钱更宝贵。}\\
\emph{Wénhuà yíchǎn bǐ qián gèng bǎoguì.}\\
Littéral : « Culture patrimoine comparé à argent plus précieux. »\\
Naturel : « Le patrimoine culturel est plus précieux que l'argent. »\\
Observez la comparaison avec \zh{比} : A + \zh{比} + B + adjectif. Ici A = \zh{文化遗产}, B = \zh{钱}.

\textbf{Phrase 7.} \zh{我们要把爷爷的书法传下去。}\\
\emph{Wǒmen yào bǎ yéye de shūfǎ chuán xiàqù.}\\
Littéral : « Nous devons + \zh{把} + calligraphie de grand-père + transmettre-vers-le-bas. »\\
Naturel : « Nous devons transmettre la calligraphie de grand-père aux générations futures. »\\
Observez la structure en \zh{把} : l'objet (\zh{爷爷的书法}) est placé \textbf{avant le verbe}. \zh{传下去} est un verbe à complément directionnel : « transmettre en descendant (dans le temps) ».

\courssub{Tableau récapitulatif des trois premières phrases}

\begin{popfigure}
\begin{tikzpicture}
\matrix[matrix of nodes, nodes={draw=popline, fill=popcream, text width=4.6cm, align=center, minimum height=1.15cm, font=\small}, row sep=-\pgflinewidth, column sep=-\pgflinewidth] {
|[fill=popblueL, font=\small\bfseries]| Caractères & |[fill=popblueL, font=\small\bfseries]| Pinyin & |[fill=popblueL, font=\small\bfseries]| Français \\
我的爷爷去年去世了。 & Wǒ de yéye qùnián qùshì le. & Mon grand-père est décédé l'année dernière. \\
他给我们留下了很多遗产。 & Tā gěi wǒmen liúxià le hěn duō yíchǎn. & Il nous a laissé beaucoup de biens. \\
爷爷生前立了遗嘱。 & Yéye shēngqián lì le yízhǔ. & De son vivant, grand-père a rédigé un testament. \\
};
\end{tikzpicture}
\legende{Les trois premières phrases du texte : caractères, pinyin, traduction.}
\end{popfigure}

\courssub{Observation grammaticale : la structure 把 + objet + 留给}

Dans la phrase 3, nous avons rencontré \zh{留给} (\emph{liú gěi}, « laisser à »). Reprenons cette structure, très fréquente dans le thème de l'héritage.

\begin{exemplebox}[La structure 把 + objet + 留给 + personne]
\zh{他把房子留给爸爸。}\\
\emph{Tā bǎ fángzi liú gěi bàba.}\\
« Il a laissé la maison à papa. »\\[4pt]
\zh{奶奶把她的菜谱留给我妈妈。}\\
\emph{Nǎinai bǎ tā de càipǔ liú gěi wǒ māma.}\\
« Grand-mère a laissé son livre de recettes à ma maman. »\\[4pt]
\zh{老师把知识留给学生。}\\
\emph{Lǎoshī bǎ zhīshi liú gěi xuésheng.}\\
« Le professeur transmet son savoir à ses élèves. »
\end{exemplebox}

\begin{propriete}[Règle : la construction en 把]
Structure : \textbf{Sujet + 把 + objet + verbe + complément}.\\
La construction en \zh{把} (\emph{bǎ}) s'emploie quand on veut insister sur ce que le sujet \emph{fait subir} à l'objet (le déplacer, le transformer, le transmettre). L'objet, normalement après le verbe, est avancé avant le verbe grâce à \zh{把}. Le verbe ne peut \textbf{jamais} rester seul : il doit être suivi d'un complément (\zh{给爸爸}, \zh{下去}, \zh{了}...).\\
Comparaison avec le français : le français dit « Il a laissé la maison à papa » (verbe puis objet) ; le chinois dit littéralement « Lui + \zh{把} + maison + laisser-à-papa » (objet puis verbe).
\end{propriete}

\begin{attention}[Erreur fréquente des francophones]
Ne dites jamais \zh{他把房子留} : un verbe nu après \zh{把} est impossible. Il faut toujours un complément : \zh{他把房子留给爸爸} ou au minimum \zh{他把房子留了}. Autre piège : ne placez pas le complément de temps après le verbe comme en français ; dites \zh{爷爷去年去世了} et non \zh{爷爷去世了去年}.
\end{attention}

\courssub{Méthode de lecture et activités de prononciation}

\begin{methode}[Comprendre un texte chinois en quatre étapes]
\begin{enumerate}
\item \textbf{Lire une fois pour le sens global} : sans s'arrêter sur les mots inconnus, identifier le sujet général (ici : un héritage familial).
\item \textbf{Repérer les mots connus} : souligner les mots déjà appris (\zh{爷爷}, \zh{房子}, \zh{钱}...) pour construire une première hypothèse de sens.
\item \textbf{Deviner les mots inconnus par le contexte} : par exemple, \zh{遗嘱} est inconnu, mais « grand-père + de son vivant + maison à papa, argent aux enfants » permet de deviner « testament ».
\item \textbf{Vérifier avec la traduction} : comparer son hypothèse avec la traduction française et noter les écarts dans son cahier.
\end{enumerate}
\end{methode}

\begin{experience}[Lecture chorale puis individuelle]
\textbf{Étape 1 — Lecture chorale.} Le professeur lit chaque phrase ; la classe répète en chœur deux fois, en suivant le pinyin. Insistez sur les tons : \zh{遗产} (\emph{yíchǎn}, 2\up{e} + 3\up{e} ton), \zh{去世} (\emph{qùshì}, 4\up{e} + 4\up{e} ton), \zh{宝贵} (\emph{bǎoguì}, 3\up{e} + 4\up{e} ton).\\
\textbf{Étape 2 — Lecture individuelle.} Chaque élève lit une phrase à voix haute ; la classe corrige les tons à l'aide du pinyin affiché. Le professeur note au tableau les tons les plus souvent confondus.\\
\textbf{Étape 3 — Repérage.} Surlignez dans le texte, en couleur, tous les mots du glossaire de la leçon : \zh{遗产}, \zh{去世}, \zh{遗嘱}, \zh{生前}, \zh{书法}, \zh{宝贵}, \zh{传下去}.
\end{experience}

\begin{exoresolu}[Retraduire en chinois]
Énoncé : traduisez en chinois la phrase « Il a laissé la maison à mon père. »
\tcblower
Solution détaillée : on identifie d'abord la structure : Sujet + \zh{把} + objet + \zh{留给} + personne. Sujet : \zh{他} ; objet : \zh{房子} ; personne : \zh{我爸爸}. On obtient : \zh{他把房子留给我爸爸。} (\emph{Tā bǎ fángzi liú gěi wǒ bàba.}). Vérification : le verbe \zh{留} n'est pas nu (complément \zh{给} + personne) : la phrase est correcte. Erreur à éviter : \zh{他留房子给我爸爸} est compréhensible mais moins naturel ; avec \zh{把}, on insiste sur la transmission, ce qui convient parfaitement au thème de l'héritage.
\end{exoresolu}

\begin{aretenir}[L'essentiel de cette section]
\begin{itemize}
\item Le texte raconte un héritage double : matériel (maison, argent, terrain) et culturel (la calligraphie).
\item \zh{去世} est le terme poli pour « décéder » ; jamais \zh{死} pour un proche.
\item Le complément de temps se place avant le verbe : \zh{去年去世了}.
\item Chaque nom compté exige son classificateur : \zh{一栋房子}, \zh{一块地}.
\item La structure \zh{把} + objet + \zh{留给} + personne exprime la transmission ; le verbe après \zh{把} porte toujours un complément.
\item Méthode de lecture en quatre étapes : sens global, mots connus, déduction par le contexte, vérification.
\end{itemize}
\end{aretenir}

Maintenant que le texte est lu, compris et prononcé correctement, nous pouvons passer à l'étude systématique du vocabulaire : la section suivante présente le glossaire complet du thème de l'héritage sous forme de tableau.

\coursec{Glossaire du thème en tableau}

Après avoir lu, écouté et traduit le texte support sur l'héritage laissé par le grand-père à sa famille, nous disposons maintenant d'un premier contact vivant avec les mots du thème. Mais un mot rencontré une seule fois dans un texte reste fragile en mémoire. Cette troisième étape de la leçon a donc pour but de \cle{fixer le lexique} : nous allons reprendre chaque mot-clé, l'observer dans son contexte, préciser sa nature grammaticale, sa prononciation exacte (pinyin et tons), ses emplois et ses combinaisons fréquentes (les \cle{collocations}). C'est ce travail méthodique qui transforme un mot « vu » en un mot « su ».

\courssub{Observation : les mots du thème dans le texte}

Avant tout tableau, reprenons les phrases du texte support où apparaissent nos mots-clés. Lisez-les à voix haute en suivant le pinyin, puis vérifiez le sens grâce à la traduction.

\begin{exemplebox}[Les mots-clés en contexte]
爷爷去世后，留下了一份遗嘱。\\
\emph{Yéye qùshì hòu, liúxià le yī fèn yízhǔ.}\\
« Après le décès de grand-père, il a laissé un testament. »\\[4pt]
他把财产分给了三个孩子。\\
\emph{Tā bǎ cáichǎn fēn gěi le sān ge háizi.}\\
« Il a partagé ses biens entre ses trois enfants. »\\[4pt]
最宝贵的遗产不是钱，而是家风。\\
\emph{Zuì bǎoguì de yíchǎn bú shì qián, ér shì jiāfēng.}\\
« L'héritage le plus précieux n'est pas l'argent, mais l'esprit de la famille. »\\[4pt]
我们要把这个传统传下去。\\
\emph{Wǒmen yào bǎ zhège chuántǒng chuán xiàqu.}\\
« Nous devons transmettre cette tradition aux générations suivantes. »
\end{exemplebox}

Observez bien : le mot \zh{遗产} apparaît tantôt avec un sens concret (des biens matériels), tantôt avec un sens abstrait (des valeurs, des traditions). C'est exactement comme le mot français « héritage », qui désigne aussi bien la maison du grand-père que « l'héritage culturel camerounais ». Le chinois et le français fonctionnent ici de la même manière : bonne nouvelle pour vous !

\courssub{Tableau récapitulatif du glossaire}

Voici le glossaire complet du thème « héritage ». Apprenez-le en trois temps : d'abord la reconnaissance (caractères vers français), puis la prononciation (pinyin et tons), enfin la production (français vers chinois).

\begin{center}
\begin{tabularx}{\linewidth}{|X|X|X|X|}
\hline
\rowcolor{popblueL} Caractères & Pinyin & Nature & Traduction \\
\hline
遗产 & yíchǎn & n. & héritage, patrimoine \\
\hline
继承 & jìchéng & v. & hériter, succéder, perpétuer \\
\hline
遗嘱 & yízhǔ & n. & testament \\
\hline
财产 & cáichǎn & n. & biens, fortune, patrimoine matériel \\
\hline
去世 & qùshì & v. & décéder (registre poli) \\
\hline
留下 & liúxià & v. & laisser (derrière soi) \\
\hline
分 & fēn & v. & partager, diviser \\
\hline
份 & fèn & n. / classif. & part, portion \\
\hline
宝贵 & bǎoguì & adj. & précieux \\
\hline
传下去 & chuán xiàqu & v. dir. & transmettre (aux générations suivantes) \\
\hline
文化遗产 & wénhuà yíchǎn & loc. n. & patrimoine culturel \\
\hline
\end{tabularx}
\end{center}

\begin{attention}[Trois pièges de prononciation pour les francophones]
\begin{itemize}
\item \zh{遗产} \emph{yíchǎn} : les deux syllabes portent le \textbf{2\textsuperscript{e} ton suivi du 3\textsuperscript{e} ton}. Attention : \zh{遗} seul est en 2\textsuperscript{e} ton montant ; ne le laissez pas retomber.
\item \zh{遗嘱} \emph{yízhǔ} : le \zh{zh} est une consonne rétroflexe (langue recourbée), très différente du « z » français. Dites \emph{zhǔ} comme dans \zh{主} \emph{zhǔ} (maître).
\item \zh{去世} \emph{qùshì} : deux 4\textsuperscript{e} tons d'affilée. En français, on a tendance à adoucir la fin de phrase ; en chinois, chaque ton doit rester franc et descendant.
\end{itemize}
\end{attention}

\courssub{Étude approfondie des mots-clés}

\begin{definition}[继承 jìchéng : hériter, succéder, perpétuer]
Le verbe \zh{继承} signifie \cle{recevoir les biens de quelqu'un après sa mort}, mais aussi, au sens figuré, \cle{continuer l'œuvre, la tradition ou l'esprit de quelqu'un}. C'est un verbe formel, utilisé à l'écrit comme à l'oral soigné. Ses deux collocations essentielles sont :
\begin{itemize}
\item \zh{继承遗产} \emph{jìchéng yíchǎn} : hériter d'un patrimoine ;
\item \zh{继承传统} \emph{jìchéng chuántǒng} : perpétuer une tradition.
\end{itemize}
Structure : \textbf{Sujet + 继承 + complément}. Le complément est directement collé au verbe, sans préposition, contrairement au français « hériter \emph{de} ».
\end{definition}

\begin{exemplebox}[继承 en action]
大儿子继承了父亲的商店。\\
\emph{Dà érzi jìchéng le fùqin de shāngdiàn.}\\
« Le fils aîné a hérité du magasin de son père. »\\[4pt]
年轻人应该继承好传统。\\
\emph{Niánqīngrén yīnggāi jìchéng hǎo chuántǒng.}\\
« Les jeunes doivent perpétuer les bonnes traditions. »
\end{exemplebox}

\begin{attention}[Le piège du « de » français]
En français on dit « hériter \textbf{de} quelque chose ». En chinois, il n'y a \textbf{aucune préposition} après \zh{继承} : on dit directement \zh{继承遗产} et jamais \zh{继承的遗产}. De même, \zh{留下} (laisser) prend son complément directement : \zh{留下遗嘱} « laisser un testament ».
\end{attention}

\begin{propriete}[分 fēn (1\textsuperscript{er} ton) et 份 fèn (4\textsuperscript{e} ton) : ne pas confondre !]
\zh{分} \emph{fēn}, au \textbf{1\textsuperscript{er} ton}, est le \textbf{verbe} « partager, diviser » : \zh{分财产} « partager les biens ».\\
\zh{份} \emph{fèn}, au \textbf{4\textsuperscript{e} ton}, est un \textbf{nom / classificateur} signifiant « part, portion » : \zh{一份遗产} \emph{yī fèn yíchǎn} « une part d'héritage », \zh{一份遗嘱} « un testament ».\\
Les deux caractères partagent le composant \zh{分}, mais \zh{份} ajoute la clé de l'homme \zh{亻} à gauche. Retenez : le \textbf{verbe} monte haut et plat (1\textsuperscript{er} ton), la \textbf{part} tombe (4\textsuperscript{e} ton).
\end{propriete}

\begin{exemplebox}[分 et 份 dans la même phrase]
钱分给我和妹妹。\\
\emph{Qián fēn gěi wǒ hé mèimei.}\\
« L'argent est partagé entre ma petite sœur et moi. »\\[4pt]
每个人都有一份遗产。\\
\emph{Měi ge rén dōu yǒu yī fèn yíchǎn.}\\
« Chacun a une part d'héritage. »\\[4pt]
爸爸把土地分成了三份。\\
\emph{Bàba bǎ tǔdì fēn chéng le sān fèn.}\\
« Papa a divisé le terrain en trois parts. »
\end{exemplebox}

Remarquez la structure très utile \textbf{分 + 给 + personne} : « partager et donner à quelqu'un », qui correspond au français « partager entre / donner sa part à ». Notez aussi la structure résultative \textbf{分成 + nombre + 份} : « diviser en X parts ».

\begin{definition}[去世 qùshì : décéder, avec délicatesse]
\zh{去世} est le terme \textbf{poli et respectueux} pour dire « mourir ». On l'emploie pour parler des personnes âgées ou respectées, exactement comme le français préfère « décéder » ou « s'éteindre » à « mourir ». Le mot direct \zh{死} \emph{sǐ} existe, mais il est brutal et à éviter dans un contexte familial respectueux. Structure : \textbf{Sujet + 去世了}.
\end{definition}

\begin{exemplebox}[Nuances de registre]
奶奶去年去世了。\\
\emph{Nǎinai qùnián qùshì le.}\\
« Grand-mère est décédée l'année dernière. » (registre respectueux, à retenir)\\[4pt]
他留下了很多宝贵的回忆。\\
\emph{Tā liúxià le hěn duō bǎoguì de huíyì.}\\
« Il a laissé beaucoup de précieux souvenirs. »
\end{exemplebox}

\begin{definition}[传下去 chuán xiàqu : transmettre aux générations suivantes]
\zh{传下去} est un \textbf{verbe directionnel} : le verbe \zh{传} « transmettre » est suivi du complément de direction \zh{下去} « vers le bas, en continuant ». L'image est celle d'un objet qui descend le long des générations, comme un flambeau. Structure : \textbf{Sujet + 把 + objet + 传下去}.
\end{definition}

\begin{exemplebox}[传下去 et le patrimoine]
这个故事要一代一代传下去。\\
\emph{Zhège gùshi yào yī dài yī dài chuán xiàqu.}\\
« Cette histoire doit être transmise de génération en génération. »\\[4pt]
中国把文化遗产传下去。\\
\emph{Zhōngguó bǎ wénhuà yíchǎn chuán xiàqu.}\\
« La Chine transmet son patrimoine culturel. »
\end{exemplebox}

\courssub{Carte mentale lexicale}

La figure suivante organise la famille de mots autour du centre \zh{遗产} : les biens matériels, le testament, la dimension culturelle et les actions de transmission.

\begin{popfigure}
\begin{tikzpicture}[mindmap, grow cyclic, every node/.style={concept, align=center, font=\small},
root concept/.append style={concept color=popblue, font=\bfseries},
level 1/.append style={level distance=3.4cm, sibling angle=90},
level 2/.append style={level distance=2.6cm, sibling angle=60},
level 1 concept/.append style={concept color=poporangeL}]
\node[root concept, align=center]{遗产\\ yíchǎn\\ héritage}
child { node[concept color=popgreenL, align=center]{财产\\ cáichǎn\\ biens}
  child { node[concept color=popgreenL, align=center]{分 fēn\\ partager} }
}
child { node[concept color=poppinkL, align=center]{遗嘱\\ yízhǔ\\ testament}
  child { node[concept color=poppinkL, align=center]{留下 liúxià\\ laisser} }
}
child { node[concept color=popgoldL, align=center]{文化遗产\\ patrimoine\\ culturel}
  child { node[concept color=popgoldL, align=center]{宝贵 bǎoguì\\ précieux} }
}
child { node[concept color=poppurpleL, align=center]{继承\\ jìchéng\\ hériter}
  child { node[concept color=poppurpleL, align=center]{传下去\\ chuán xiàqu\\ transmettre} }
};
\end{tikzpicture}
\legende{Carte mentale du champ lexical de l'héritage : chaque branche relie un nom à l'action qui lui correspond.}
\end{popfigure}

\courssub{Texte de réinvestissement lexical}

Lisons maintenant un court texte original qui réemploie presque tout le glossaire dans un contexte camerounais. C'est la meilleure façon de vérifier que les mots sont vraiment acquis.

\begin{textebox}[Le bracelet de grand-mère]
我的奶奶是杜阿拉人。她去世以前，写了一份简单的遗嘱。\\
\emph{Wǒ de nǎinai shì Dù'ālā rén. Tā qùshì yǐqián, xiě le yī fèn jiǎndān de yízhǔ.}\\
« Ma grand-mère était de Douala. Avant de décéder, elle a rédigé un testament simple. »\\[4pt]
她的财产不多：一所小房子，一些钱，还有一个金手镯。\\
\emph{Tā de cáichǎn bù duō : yī suǒ xiǎo fángzi, yīxiē qián, hái yǒu yī ge jīn shǒuzhuó.}\\
« Ses biens n'étaient pas nombreux : une petite maison, un peu d'argent, et un bracelet en or. »\\[4pt]
房子和钱分给了她的三个孩子，手镯留给了我。\\
\emph{Fángzi hé qián fēn gěi le tā de sān ge háizi, shǒuzhuó liú gěi le wǒ.}\\
« La maison et l'argent ont été partagés entre ses trois enfants, et le bracelet m'a été laissé. »\\[4pt]
妈妈说，最宝贵的遗产不是这些东西，而是奶奶的善良和勤劳。\\
\emph{Māma shuō, zuì bǎoguì de yíchǎn bú shì zhèxiē dōngxi, ér shì nǎinai de shànliáng hé qínláo.}\\
« Maman dit que l'héritage le plus précieux, ce ne sont pas ces objets, mais la bonté et le goût du travail de grand-mère. »\\[4pt]
我们要继承她的精神，把这个家风传下去。\\
\emph{Wǒmen yào jìchéng tā de jīngshén, bǎ zhège jiāfēng chuán xiàqu.}\\
« Nous devons perpétuer son esprit et transmettre cet exemple familial aux générations futures. »
\end{textebox}

\courssub{Exercice résolu et entraînement}

\begin{exoresolu}[Complétez avec le mot du glossaire qui convient]
Énoncé — Complétez chaque phrase avec le mot approprié : \zh{遗产}, \zh{继承}, \zh{遗嘱}, \zh{去世}, \zh{分}, \zh{宝贵}, \zh{传下去}.\\
1. 爷爷 \_\_\_\_ 以后，家里人才看到他的 \_\_\_\_ 。\\
2. 长城是中国最 \_\_\_\_ 的文化 \_\_\_\_ 。\\
3. 哥哥 \_\_\_\_ 了父亲的农场，把土地 \_\_\_\_ 成了两份。\\
4. 这个传统一定要 \_\_\_\_ 。
\tcblower
Solution détaillée :\\
1. \zh{爷爷去世以后，家里人才看到他的遗嘱。} — « Après le décès de grand-père, la famille a seulement alors vu son testament. » Le premier blanc exige un verbe de registre poli (\zh{去世}), le second un nom désignant le document (\zh{遗嘱}).\\
2. \zh{长城是中国最宝贵的文化遗产。} — « La Grande Muraille est le patrimoine culturel le plus précieux de la Chine. » L'adjectif \zh{宝贵} précède le nom avec \zh{的}.\\
3. \zh{哥哥继承了父亲的农场，把土地分成了两份。} — « Le grand frère a hérité de la ferme de son père et a divisé le terrain en deux parts. » Attention : \zh{继承} sans préposition, et \zh{分} au 1\textsuperscript{er} ton.\\
4. \zh{这个传统一定要传下去。} — « Cette tradition doit absolument être transmise. » Le verbe directionnel \zh{传下去} convient car il s'agit des générations futures.
\end{exoresolu}

\begin{exercice}[À vous de jouer]
Traduisez en chinois, en utilisant le glossaire :\\
1. « Ma grand-mère a laissé un testament. »\\
2. « L'héritage le plus précieux n'est pas l'argent. »\\
3. « Les enfants ont partagé les biens de leur père. »\\
4. « Nous devons perpétuer cette tradition et la transmettre. »
\end{exercice}

\begin{corrige}[Exercice « À vous de jouer »]
1. \zh{我的奶奶留下了一份遗嘱。} \emph{Wǒ de nǎinai liúxià le yī fèn yízhǔ.} — N'oubliez pas le classificateur \zh{份} devant \zh{遗嘱}.\\
2. \zh{最宝贵的遗产不是钱。} \emph{Zuì bǎoguì de yíchǎn bú shì qián.} — L'adjectif \zh{宝贵} exige la particule \zh{的} avant le nom.\\
3. \zh{孩子们分了父亲的财产。} ou avec la structure \zh{把} : \zh{孩子们把父亲的财产分了。} — Les deux sont correctes ; la structure en \zh{把} insiste sur le résultat du partage.\\
4. \zh{我们要继承这个传统，把它传下去。} \emph{Wǒmen yào jìchéng zhège chuántǒng, bǎ tā chuán xiàqu.} — Notez l'ordre : \zh{把 + objet + verbe directionnel}.
\end{corrige}

\begin{savaistu}[贝, le coquillage qui valait de l'argent]
Le caractère \zh{财} \emph{cái} « richesse » contient à gauche la clé \zh{贝} \emph{bèi}, qui signifie « coquillage ». Pourquoi ? Dans la Chine antique, il y a plus de trois mille ans, les coquillages servaient de \textbf{monnaie} ! C'est pourquoi de nombreux caractères liés à l'argent portent cette clé : \zh{财} \emph{cái} « richesse », \zh{贵} \emph{guì} « cher, précieux » (dans \zh{宝贵}), \zh{货} \emph{huò} « marchandise ». Retenez l'astuce : \textbf{coquillage = argent}. Au Cameroun et dans plusieurs régions d'Afrique, les cauris (de petits coquillages) ont aussi servi de monnaie avant les pièces et les billets : les deux civilisations se rejoignent dans l'écriture et dans l'histoire !
\end{savaistu}

\begin{aretenir}[L'essentiel du glossaire]
\begin{itemize}
\item \zh{遗产} désigne l'héritage matériel \textbf{et} immatériel ; \zh{文化遗产} est le patrimoine culturel.
\item \zh{继承} prend son complément \textbf{sans préposition} : \zh{继承遗产}, \zh{继承传统}.
\item \zh{分} \emph{fēn} (verbe, 1\textsuperscript{er} ton) et \zh{份} \emph{fèn} (part, 4\textsuperscript{e} ton) : deux tons, deux natures, deux emplois.
\item \zh{去世} est le mot poli pour « décéder » ; \zh{留下} exprime ce qu'on laisse derrière soi.
\item \zh{传下去} est un verbe directionnel : la tradition « descend » vers les générations futures.
\end{itemize}
\end{aretenir}

Le glossaire est désormais structuré et travaillé. Dans la prochaine étape de la leçon, nous observerons les \textbf{clés (radicaux) et l'ordre des traits} de quelques caractères du thème, afin d'apprendre à les écrire correctement, trait après trait.

\coursec{Clés (radicaux) et ordre des traits}

Dans la section précédente, nous avons dressé le glossaire du thème de l'héritage : \zh{遗产} \emph{yíchǎn} « héritage », \zh{继承} \emph{jìchéng} « hériter », \zh{财产} \emph{cáichǎn} « biens, fortune », \zh{遗嘱} \emph{yízhǔ} « testament ». Ces mots ne sont pas des dessins mystérieux : chacun est construit selon une logique que l'on peut lire. Cette section vous apprend à \cle{décomposer} les caractères, à reconnaître leur \cle{clé} (radical) qui donne le sens, leur \cle{composant phonétique} qui suggère le son, et à les écrire dans le bon \cle{ordre des traits}. C'est la compétence qui transforme un élève qui « copie des dessins » en un élève qui \emph{lit} l'écriture chinoise.

\courssub{Observation : cinq caractères, cinq histoires}

Regardez d'abord ces cinq caractères extraits de notre texte support. Ne cherchez pas encore à les écrire : \emph{observez} leurs parties.

\begin{center}
\begin{tabularx}{\linewidth}{|X|X|X|X|}\hline
\rowcolor{popblueL} Caractère & Pinyin & Traduction & Décomposition \\ \hline
遗 & yí & laisser, léguer & 辶 + 贵 \\ \hline
产 & chǎn & produire, bien & 亠 + partie inférieure \\ \hline
继 & jì & continuer, succéder & 纟 + 㡭 \\ \hline
财 & cái & richesse & 贝 + 才 \\ \hline
嘱 & zhǔ & recommander, léguer & 口 + 属 \\ \hline
\end{tabularx}
\end{center}

Que remarquez-vous ? Chaque caractère contient une petite pièce qui revient dans beaucoup d'autres caractères : \zh{辶}, \zh{纟}, \zh{贝}, \zh{口}. Ces pièces sont les \cle{clés} (ou radicaux, en chinois \zh{部首} \emph{bùshǒu}). La clé indique la \emph{famille de sens} ; l'autre partie, le \cle{composant phonétique}, donne souvent une indication sur la prononciation. C'est le principe des \cle{caractères sémantico-phonétiques}, qui représentent la grande majorité des sinogrammes.

\begin{definition}[Clé et composant phonétique]
La \cle{clé} (部首 \emph{bùshǒu}) est la partie du caractère qui classe le mot dans une famille de sens et qui sert à le retrouver dans un dictionnaire. Le \cle{composant phonétique} est la partie qui suggère la lecture (le son), souvent de façon approximative. Exemple : dans \zh{财} \emph{cái}, la clé \zh{贝} donne le sens (l'argent) et \zh{才} \emph{cái} donne le son (identique, ton inclus).
\end{definition}

\courssub{Étude détaillée des cinq caractères du thème}

\textbf{1. 遗 \emph{yí} — « laisser, léguer » : clé 辶 (12 traits)}\par

Le caractère \zh{遗} se compose de la clé \zh{辶}, appelée \emph{chuò} ou \emph{zǒuzhīdǐ}, la clé du \cle{mouvement}, de la marche (on la retrouve dans \zh{进} \emph{jìn} « entrer », \zh{送} \emph{sòng} « offrir, accompagner », \zh{远} \emph{yuǎn} « loin »), et du composant phonétique \zh{贵} \emph{guì} « précieux ». Image mémorable : ce qui est \emph{précieux} (\zh{贵}) et qui \emph{s'en va, se transmet} (\zh{辶}) d'une génération à l'autre, c'est précisément l'héritage. Ordre des traits : on écrit d'abord tout l'intérieur \zh{贵} (9 traits, de haut en bas), puis seulement à la fin la clé du mouvement \zh{辶} en 3 traits : le point, le virage, puis la grande vague finale qui « porte » le caractère.

\textbf{2. 产 \emph{chǎn} — « produire, bien » : clé 亠 (7 traits)}\par

La clé \zh{亠}, appelée \emph{tóu} (« tête »), est le petit chapeau formé d'un point sur un trait horizontal. Ordre des traits : d'abord le point, puis le trait horizontal, ensuite les traits inférieurs (de haut en bas, comme toujours). On le trouve dans \zh{生产} \emph{shēngchǎn} « produire », \zh{产品} \emph{chǎnpǐn} « produit », \zh{财产} \emph{cáichǎn} « biens ». \textbf{Attention :} \zh{产} compte \textbf{7 traits} au total (2 pour 亠 + 5 pour la partie basse), et non 6.

\textbf{3. 继 \emph{jì} — « continuer, succéder » : clé 纟 (10 traits)}\par

La clé \zh{纟} est la clé de la \cle{soie}, du fil (forme simplifiée de \zh{糸}, \emph{mì}). Elle s'écrit en 3 traits, à gauche du caractère. Quelle belle image pour le thème de l'héritage : \zh{继承} \emph{jìchéng} « hériter », c'est littéralement le \emph{fil} qui relie les générations, qui ne se rompt pas. Le composant de droite, \zh{㡭}, donne le son \emph{jì} (correspondance parfaite). Ordre : la clé soie à gauche d'abord (gauche avant droite), puis la partie droite de haut en bas.

\textbf{4. 财 \emph{cái} — « richesse » : clé 贝 (7 traits)}\par

La clé \zh{贝} représente un \cle{coquillage} : dans la Chine ancienne, les coquillages servaient de monnaie, c'est pourquoi presque tous les caractères liés à l'argent contiennent \zh{贝}. À droite, \zh{才} \emph{cái} donne le son (correspondance parfaite, ton inclus). Ordre : \zh{贝} à gauche (4 traits), puis \zh{才} à droite (3 traits) — règle « gauche avant droite ».

\textbf{5. 嘱 \emph{zhǔ} — « recommander, léguer » : clé 口 (15 traits)}\par

La clé \zh{口} est la \cle{bouche}. Un testament, \zh{遗嘱} \emph{yízhǔ}, se \emph{dicte} d'abord oralement : le mourant confie ses dernières volontés de vive voix. La clé bouche, petite, se place à gauche et s'écrit en premier (3 traits) ; la grande partie droite \zh{属} \emph{shǔ} suit, de haut en bas et de l'extérieur vers l'intérieur (structure en « cadre » pour \zh{尸} puis contenu). Le composant phonétique \zh{属} suggère le son \emph{shǔ} (proche de \emph{zhǔ}, même voyelle, consonne initiale différente).

\begin{propriete}[Règle générale de l'ordre des traits]
Quatre principes régissent l'écriture de tous les caractères chinois :
\begin{enumerate}
\item \cle{De haut en bas} : \zh{产} (le point et le trait du haut d'abord).
\item \cle{De gauche à droite} : \zh{继}, \zh{财}, \zh{嘱} (la clé de gauche d'abord).
\item \cle{L'extérieur avant l'intérieur} : on trace d'abord l'enveloppe, puis le contenu (ex. \zh{同}, \zh{月}).
\item \cle{On ferme en dernier} : le trait qui ferme une boîte ou qui « porte » le caractère vient à la fin.
\end{enumerate}
\cle{Exception célèbre} : pour les caractères à enveloppe du mouvement \zh{辶} (et \zh{辶} seul), on écrit \emph{l'intérieur d'abord} et la clé \zh{辶} \textbf{en dernier}.
\end{propriete}

\begin{popfigure}
\begin{tikzpicture}[node distance=0.5cm]
\node[draw=popblue, fill=popblueL, rounded corners, minimum width=1.1cm, minimum height=1.1cm, font=\Huge] (yi) {遗};
\node[right=0.6cm of yi, font=\Large] (eq1) {=};
\node[draw=popgreen, fill=popgreenL, rounded corners, right=0.6cm of eq1, minimum width=1cm, minimum height=1cm, font=\Huge] (chuo) {辶};
\node[right=0.4cm of chuo, font=\Large] (plus1) {+};
\node[draw=poporange, fill=poporangeL, rounded corners, right=0.4cm of plus1, minimum width=1cm, minimum height=1cm, font=\Huge] (gui) {贵};
\node[below=0.25cm of chuo, align=center, font=\small, text=popdark] {clé : mouvement\\(écrite en dernier)};
\node[below=0.25cm of gui, align=center, font=\small, text=popdark] {phonétique\\guì « précieux »};
\node[draw=popblue, fill=popblueL, rounded corners, minimum width=1.1cm, minimum height=1.1cm, font=\Huge, below=1.6cm of yi] (ji) {继};
\node[right=0.6cm of ji, font=\Large] (eq2) {=};
\node[draw=popgreen, fill=popgreenL, rounded corners, right=0.6cm of eq2, minimum width=1cm, minimum height=1cm, font=\Huge] (si) {纟};
\node[right=0.4cm of si, font=\Large] (plus2) {+};
\node[draw=poporange, fill=poporangeL, rounded corners, right=0.4cm of plus2, minimum width=1cm, minimum height=1cm, font=\Huge] (ji2) {㡭};
\node[below=0.25cm of si, align=center, font=\small, text=popdark] {clé : soie\\(le fil, le lien)};
\node[below=0.25cm of ji2, align=center, font=\small, text=popdark] {phonétique\\jì};
\node[draw=popblue, fill=popblueL, rounded corners, minimum width=1.1cm, minimum height=1.1cm, font=\Huge, below=1.6cm of ji] (cai) {财};
\node[right=0.6cm of cai, font=\Large] (eq3) {=};
\node[draw=popgreen, fill=popgreenL, rounded corners, right=0.6cm of eq3, minimum width=1cm, minimum height=1cm, font=\Huge] (bei) {贝};
\node[right=0.4cm of bei, font=\Large] (plus3) {+};
\node[draw=poporange, fill=poporangeL, rounded corners, right=0.4cm of plus3, minimum width=1cm, minimum height=1cm, font=\Huge] (cai2) {才};
\node[below=0.25cm of bei, align=center, font=\small, text=popdark] {clé : coquillage\\(l'argent)};
\node[below=0.25cm of cai2, align=center, font=\small, text=popdark] {phonétique\\cái};
\end{tikzpicture}
\legende{Décomposition de trois caractères du thème : la clé (sens) en vert, le composant phonétique (son) en orange.}
\end{popfigure}

\courssub{Familles de sens : apprendre un caractère, en reconnaître dix}

\begin{exemplebox}[La famille de 贝 : la « valeur »]
Le coquillage-monnaie \zh{贝} apparaît dans toute une famille de caractères liés à la valeur :
\begin{itemize}
\item \zh{财} \emph{cái} : richesse → \zh{财产} \emph{cáichǎn} « les biens ».
\item \zh{贵} \emph{guì} : précieux, cher → \zh{很贵} \emph{hěn guì} « très cher ».
\item \zh{贵} se retrouve à l'intérieur de \zh{遗} \emph{yí} : ce qui est précieux et transmis.
\item \zh{买} \emph{mǎi} « acheter » et \zh{卖} \emph{mài} « vendre » (forme simplifiée contenant l'idée d'échange).
\item \zh{货} \emph{huò} « marchandise », \zh{账} \emph{zhàng} « compte, facture ».
\end{itemize}
Retenir la famille, c'est retenir six caractères pour le prix d'un.
\end{exemplebox}

\begin{center}
\begin{tabularx}{\linewidth}{|X|X|X|X|}\hline
\rowcolor{popblueL} Clé & Sens de la clé & Caractères du thème & Autres exemples connus \\ \hline
辶 & mouvement, marche & 遗 yí (léguer) & 进 jìn (entrer), 送 sòng (offrir), 远 yuǎn (loin) \\ \hline
亠 & tête (chapeau) & 产 chǎn (produire) & 京 jīng (capitale), 高 gāo (haut) \\ \hline
纟 & soie, fil & 继 jì (continuer) & 结 jié (nouer), 红 hóng (rouge), 纸 zhǐ (papier) \\ \hline
贝 & coquillage, argent & 财 cái (richesse) & 贵 guì (cher), 货 huò (marchandise) \\ \hline
口 & bouche & 嘱 zhǔ (recommander) & 吃 chī (manger), 说 shuō (parler), 问 wèn (demander) \\ \hline
\end{tabularx}
\end{center}

\begin{methode}[Comment mémoriser un caractère en quatre étapes]
\begin{enumerate}
\item \cle{Identifier la clé} : demandez-vous « quelle est la famille de sens ? ». Dans \zh{嘱}, c'est la bouche \zh{口} : le sens touche à la parole.
\item \cle{Repérer le composant phonétique} : « quelle partie me donne le son ? ». Dans \zh{财}, c'est \zh{才} \emph{cái} (correspondance exacte) ; dans \zh{遗}, c'est \zh{贵} \emph{guì} (proche).
\item \cle{Se raconter l'histoire du caractère} : « le fil (\zh{纟}) qui continue (\zh{继}) de génération en génération ». Une image vaut dix répétitions.
\item \cle{Écrire 5 fois en comptant les traits à voix haute} : « un, deux, trois… douze » pour \zh{遗}, en respectant l'ordre. La main mémorise autant que l'œil.
\end{enumerate}
\end{methode}

\begin{attention}[Le piège de 辶 : la clé qui s'écrit en dernier]
Erreur classique des francophones : écrire d'abord la clé du mouvement \zh{辶}, parce qu'elle « entoure » le caractère, puis caser \zh{贵} dedans. C'est faux. Dans \zh{遗}, on écrit \textbf{d'abord} tout l'intérieur \zh{贵} (9 traits), \textbf{puis} la clé \zh{辶} en 3 traits : 1) le point, 2) le virage, 3) la grande vague finale qui passe sous le caractère. Autre piège : ne pas confondre \zh{遗} \emph{yí} « léguer » avec \zh{贵} \emph{guì} « cher » seul — la présence de \zh{辶} change tout le sens. Enfin, attention à \zh{继} : la clé soie \zh{纟} compte 3 traits, pas 2 ; le deuxième trait est un pli, le troisième une montée.
\end{attention}

\begin{savaistu}[Le fil de soie, image du lien]
La clé \zh{纟} (soie) se retrouve dans \zh{继} \emph{jì} « continuer », \zh{结} \emph{jié} « nouer », \zh{红} \emph{hóng} « rouge » (la couleur des teintures de soie) et \zh{线} \emph{xiàn} « le fil ». Pour les Chinois, le fil est l'image même du lien qui ne se rompt pas : on parle du « fil rouge du destin » (\zh{姻缘红线} \emph{yīnyuán hóngxiàn}) qui relie les personnes. Au Cameroun, on dirait que la famille est un tissu : chaque génération est un fil du même pagne. Hériter, \zh{继承}, c'est recevoir le fil et le continuer — la métaphore est universelle, mais le chinois l'a littéralement écrite dans le caractère.
\end{savaistu}

\begin{exoresolu}[Décomposer et écrire]
Énoncé — Pour chaque caractère : a) donne la clé et son sens ; b) donne le nombre total de traits ; c) indique quelle partie s'écrit en premier. Caractères : \zh{遗}, \zh{财}, \zh{嘱}.
\tcblower
Solution détaillée :
\begin{enumerate}
\item \zh{遗} \emph{yí} : a) clé \zh{辶}, le mouvement ; b) 12 traits au total ; c) on écrit d'abord l'intérieur \zh{贵} (9 traits), la clé \zh{辶} vient en dernier (point, virage, grande vague).
\item \zh{财} \emph{cái} : a) clé \zh{贝}, le coquillage, symbole de l'argent ; b) 7 traits ; c) la clé \zh{贝} à gauche d'abord (4 traits), puis \zh{才} à droite (3 traits) — règle « gauche avant droite ».
\item \zh{嘱} \emph{zhǔ} : a) clé \zh{口}, la bouche (le testament se dicte oralement) ; b) 15 traits ; c) la petite bouche \zh{口} à gauche en premier, puis la partie droite \zh{属} de haut en bas.
\end{enumerate}
\end{exoresolu}

\begin{exercice}[À vous de tracer]
1. Écris cinq fois \zh{继} en comptant les traits à voix haute, puis entoure la clé.\\
2. Retrouve la clé commune à \zh{财} et \zh{贵} et explique le lien de sens avec le thème de l'héritage.\\
3. Classe ces caractères selon leur clé : \zh{嘱}, \zh{远}, \zh{红}, \zh{遗}, \zh{吃}, \zh{结} (clés : \zh{口}, \zh{辶}, \zh{纟}).
\end{exercice}

\begin{corrige}[Exercice « À vous de tracer »]
1. \zh{继} compte 10 traits : la clé \zh{纟} (3 traits) à gauche, puis la partie droite (7 traits). La clé à entourer est \zh{纟}.\\
2. La clé commune est \zh{贝} (coquillage, monnaie ancienne) : \zh{财} « richesse » et \zh{贵} « précieux » appartiennent à la famille de la valeur — au cœur du vocabulaire de l'héritage.\\
3. Clé \zh{口} : \zh{嘱}, \zh{吃}. Clé \zh{辶} : \zh{远}, \zh{遗}. Clé \zh{纟} : \zh{红}, \zh{结}.
\end{corrige}

\begin{aretenir}[L'essentiel de la section]
\begin{itemize}
\item Un caractère = une \cle{clé} (le sens) + souvent un \cle{composant phonétique} (le son).
\item \zh{遗} (clé \zh{辶}, 12 traits, clé écrite en dernier), \zh{产} (clé \zh{亠}, 7 traits), \zh{继} (clé \zh{纟}, 10 traits), \zh{财} (clé \zh{贝}, 7 traits), \zh{嘱} (clé \zh{口}, 15 traits).
\item Ordre des traits : haut → bas, gauche → droite, extérieur → intérieur, on ferme en dernier.
\item Apprendre par familles de clés multiplie le vocabulaire mémorisé.
\end{itemize}
\end{aretenir}

\coursec{Compréhension du texte et collocations}

Dans la section précédente, nous avons appris à écrire les caractères clés du thème de l'héritage (\zh{遗}, \zh{产}, \zh{嘱}, \zh{继}) en respectant les clés et l'ordre des traits. Savoir écrire, c'est bien ; savoir \emph{comprendre} et \emph{réutiliser}, c'est encore mieux. Nous allons maintenant vérifier notre compréhension du texte support, apprendre à construire des réponses complètes, puis mémoriser les \cle{collocations} (associations de mots) indispensables pour parler de l'héritage en chinois.

\courssub{Rappel du texte support}

Avant de répondre aux questions, relisons le texte de la leçon, que voici dans sa version intégrale.

\begin{textebox}[爷爷的遗产 — L'héritage de grand-père]
爷爷去年去世了。他给我们留下了很多遗产：一套房子、一些钱和一块地。\\
Yéye qùnián qùshì le. Tā gěi wǒmen liúxià le hěn duō yíchǎn : yí tào fángzi, yìxiē qián hé yí kuài dì.\\
Grand-père est décédé l'année dernière. Il nous a laissé beaucoup de biens : une maison, un peu d'argent et un lopin de terre.\\[4pt]
爷爷生前立了遗嘱。遗嘱里说：房子给爸爸，钱给我和妹妹上学用，地不能卖，要留给孩子们。\\
Yéye shēngqián lì le yízhǔ. Yízhǔ li shuō : fángzi gěi bàba, qián gěi wǒ hé mèimei shàngxué yòng, dì bù néng mài, yào liúgěi háizimen.\\
De son vivant, grand-père a rédigé un testament. Le testament dit : la maison revient à papa, l'argent sert à payer nos études, à ma petite sœur et à moi, et la terre ne doit pas être vendue, elle doit être gardée pour les enfants.\\[4pt]
可是奶奶说：\zh{«}最重要的遗产不是钱，而是爷爷的书法。书法是我们家的文化，要一代一代传下去。\zh{»}\\
Kěshì nǎinai shuō : « Zuì zhòngyào de yíchǎn bú shì qián, ér shì yéye de shūfǎ. Shūfǎ shì wǒmen jiā de wénhuà, yào yí dài yí dài chuán xiàqù. »\\
Mais grand-mère dit : « L'héritage le plus important n'est pas l'argent, mais la calligraphie de grand-père. La calligraphie est la culture de notre famille, il faut la transmettre de génération en génération. »\\[4pt]
我觉得奶奶说得对。钱会用完，可是文化不会。我要学书法，也要把家里的故事传下去。\\
Wǒ juéde nǎinai shuō de duì. Qián huì yòng wán, kěshì wénhuà bù huì. Wǒ yào xué shūfǎ, yě yào bǎ jiā li de gùshi chuán xiàqù.\\
Je trouve que grand-mère a raison. L'argent s'épuise, mais pas la culture. Je veux apprendre la calligraphie et transmettre aussi les histoires de notre famille.
\end{textebox}

\courssub{Méthode : répondre à une question de compréhension}

En chinois comme en français, une bonne réponse de compréhension est une \cle{phrase complète}, pas un mot isolé. La technique la plus sûre consiste à \emph{reprendre les mots de la question} et à remplacer le mot interrogatif par l'information trouvée dans le texte.

\begin{methode}[Répondre à une question de compréhension en chinois]
\begin{enumerate}
\item \textbf{Souligner le mot interrogatif} : \zh{谁} (shéi, qui), \zh{什么} (shénme, quoi), \zh{什么时候} (shénme shíhou, quand), \zh{为什么} (wèishénme, pourquoi), \zh{哪儿} (nǎr, où).
\item \textbf{Localiser la phrase du texte} qui contient la réponse : le mot interrogatif indique le type d'information à chercher (temps, personne, objet, raison).
\item \textbf{Recopier le squelette de la question} en remplaçant le mot interrogatif par la réponse : \zh{爷爷什么时候去世的？} → \zh{爷爷去年去世的。}
\item \textbf{Vérifier} : la réponse doit être une phrase complète, avec sujet et verbe, sans mot interrogatif.
\end{enumerate}
\end{methode}

\begin{attention}[La place de \zh{什么} : le piège numéro un des francophones]
En français, le mot interrogatif se place \emph{au début} de la question : « \textbf{Qu'}a-t-il laissé ? ». En chinois, le mot interrogatif reste \emph{à la place du mot qu'il remplace}, souvent \textbf{en fin de phrase} : \zh{他留下了什么？} (Tā liúxià le shénme ?) — littéralement « Il a laissé quoi ? ».\\
Ne dites jamais \zh{什么他留下了？} : c'est une erreur de calque du français. L'ordre chinois est : Sujet + Verbe + \zh{什么} ?
\end{attention}

\courssub{Questions de compréhension : réponses guidées}

Appliquons la méthode aux cinq questions de la leçon.

\begin{exoresolu}[Questions de compréhension sur le texte]
Énoncé : répondez en phrases complètes aux cinq questions suivantes.\\
1) \zh{爷爷什么时候去世的？} 2) \zh{他留下了什么遗产？} 3) \zh{遗嘱里说了什么？} 4) \zh{奶奶认为最重要的遗产是什么？} 5) \zh{为什么文化遗产比钱更宝贵？}
\tcblower
\textbf{Solutions détaillées :}\\[3pt]
1) Mot interrogatif : \zh{什么时候} (quand) → on cherche le temps dans le texte : \zh{去年}.\\
\zh{爷爷去年去世的。} (Yéye qùnián qùshì de.) — Grand-père est décédé l'année dernière.\\[3pt]
2) Mot interrogatif : \zh{什么} (quoi) → on cherche la liste des biens.\\
\zh{他留下了一套房子、一些钱和一块地。} (Tā liúxià le yí tào fángzi, yìxiē qián hé yí kuài dì.) — Il a laissé une maison, un peu d'argent et un lopin de terre.\\[3pt]
3) On localise la phrase \zh{遗嘱里说…} et on la reprend intégralement.\\
\zh{遗嘱里说：房子给爸爸，钱给我和妹妹上学用，地不能卖，要留给孩子们。} (Yízhǔ li shuō : fángzi gěi bàba, qián gěi wǒ hé mèimei shàngxué yòng, dì bù néng mài, yào liúgěi háizimen.) — Le testament dit : la maison revient à papa, l'argent sert à nos études, la terre ne doit pas être vendue et doit être laissée aux enfants.\\[3pt]
4) On reprend la phrase de grand-mère.\\
\zh{奶奶认为最重要的遗产不是钱，而是爷爷的书法。} (Nǎinai rènwéi zuì zhòngyào de yíchǎn bú shì qián, ér shì yéye de shūfǎ.) — Grand-mère considère que l'héritage le plus important n'est pas l'argent, mais la calligraphie de grand-père.\\[3pt]
5) Mot interrogatif : \zh{为什么} (pourquoi) → on répond avec \zh{因为}.\\
\zh{因为钱会用完，可是文化不会。} (Yīnwèi qián huì yòng wán, kěshì wénhuà bù huì.) — Parce que l'argent s'épuise, mais pas la culture.
\end{exoresolu}

\begin{exemplebox}[La structure \zh{不是…而是…} (ce n'est pas… mais…)]
Cette structure de correction est très utile pour exprimer une opposition :\\
\zh{最重要的遗产不是钱，而是书法。}\\
\emph{Zuì zhòngyào de yíchǎn bú shì qián, ér shì shūfǎ.}\\
« L'héritage le plus important n'est pas l'argent, mais la calligraphie. »\\[4pt]
Autre exemple : \zh{他不是我的老师，而是我的叔叔。}\\
\emph{Tā bú shì wǒ de lǎoshī, ér shì wǒ de shūshu.} — « Ce n'est pas mon professeur, mais mon oncle. »\\[4pt]
Schéma : Sujet + \zh{不是} + A + \zh{，而是} + B. Attention au changement de ton : \zh{不} passe au 2\textsuperscript{e} ton (\zh{bú}) devant \zh{是} (4\textsuperscript{e} ton).
\end{exemplebox}

\courssub{Organigramme de compréhension du texte}

La figure suivante résume la logique du texte : qui donne, quoi, à qui, et quelle leçon en tirer.

\begin{popfigure}
\begin{tikzpicture}[
  node distance=0.55cm and 0.9cm,
  box/.style={draw=popblue, fill=popblueL, rounded corners=3pt, align=center, inner sep=5pt, font=\small},
  lab/.style={font=\small\bfseries, text=popdark},
  arr/.style={-{Stealth[length=2.5mm]}, thick, poporange}
]
\node[box, align=center] (qui){爷爷\\ \emph{yéye}\\ grand-père};
\node[box, right=of qui, align=center] (quoi){房子、钱、地、书法\\ \emph{fángzi, qián, dì, shūfǎ}\\ maison, argent, terre, calligraphie};
\node[box, right=of quoi, align=center] (pourqui){爸爸、我和妹妹\\ \emph{bàba, wǒ hé mèimei}\\ papa, moi et ma sœur};
\node[box, below=1.1cm of quoi, fill=popgreenL, draw=popgreen, align=center] (lecon){文化比钱更宝贵\\ \emph{wénhuà bǐ qián gèng bǎoguì}\\ la culture est plus précieuse que l'argent};
\node[lab, above=0.15cm of qui] {Qui ?};
\node[lab, above=0.15cm of quoi] {Quoi ?};
\node[lab, above=0.15cm of pourqui] {Pour qui ?};
\node[lab, below=0.15cm of lecon] {Leçon ?};
\draw[arr] (qui) -- (quoi);
\draw[arr] (quoi) -- (pourqui);
\draw[arr] (quoi) -- (lecon);
\end{tikzpicture}
\legende{Organigramme de compréhension : de la personne aux biens, des biens aux héritiers, des biens à la leçon du texte.}
\end{popfigure}

\courssub{Les collocations clés du thème}

En chinois, les mots ne s'emploient pas isolément : chaque verbe a ses partenaires habituels. Ces associations fixes s'appellent des \cle{collocations}. Les connaître par cœur, c'est parler un chinois naturel.

\begin{center}
\begin{tabularx}{\linewidth}{|X|X|X|X|}
\hline
\rowcolor{popblueL} Collocation & Pinyin & Structure & Traduction \\
\hline
\zh{留下遗产} & liúxià yíchǎn & verbe + objet & laisser un héritage \\
\hline
\zh{继承遗产} & jìchéng yíchǎn & verbe + objet & hériter (d'un patrimoine) \\
\hline
\zh{立遗嘱} & lì yízhǔ & verbe + objet & rédiger un testament \\
\hline
\zh{分财产} & fēn cáichǎn & verbe + objet & partager les biens \\
\hline
\zh{传下去} & chuán xiàqù & verbe + complément de direction & transmettre (aux générations suivantes) \\
\hline
\zh{用完} & yòng wán & verbe + complément de résultat & épuiser, finir d'utiliser \\
\hline
\end{tabularx}
\end{center}

\begin{exemplebox}[Les collocations en contexte]
\zh{爷爷去世后，孩子们继承了他的遗产。}\\
\emph{Yéye qùshì hòu, háizimen jìchéng le tā de yíchǎn.} — « Après le décès de grand-père, les enfants ont hérité de son patrimoine. »\\[4pt]
\zh{很多老人生前立遗嘱，免得孩子们以后分财产的时候吵架。}\\
\emph{Hěn duō lǎorén shēngqián lì yízhǔ, miǎnde háizimen yǐhòu fēn cáichǎn de shíhou chǎojià.} — « Beaucoup de personnes âgées rédigent un testament de leur vivant, pour éviter que les enfants se disputent plus tard au moment du partage des biens. »\\[4pt]
\zh{我们要把传统文化一代一代传下去。}\\
\emph{Wǒmen yào bǎ chuántǒng wénhuà yí dài yí dài chuán xiàqù.} — « Nous devons transmettre la culture traditionnelle de génération en génération. »
\end{exemplebox}

\begin{attention}[Ne confondez pas \zh{遗产} et \zh{财产}]
\zh{遗产} (yíchǎn) désigne ce qu'on reçoit \emph{après la mort} de quelqu'un (héritage, patrimoine) ; \zh{财产} (cáichǎn) désigne les biens d'une personne \emph{de son vivant} ou à partager. On dit \zh{继承遗产} (hériter) mais \zh{分财产} (partager les biens). D'autre part, \zh{传下去} exige le complément de direction \zh{下去} : on ne dit pas \zh{传文化} seul pour « transmettre la culture aux générations ».
\end{attention}

\begin{savaistu}[La Chine et le patrimoine mondial]
La Chine compte plus de cinquante sites classés au \cle{patrimoine mondial de l'UNESCO}, dont la Grande Muraille (\zh{长城}, Chángchéng), la Cité interdite (\zh{故宫}, Gùgōng) et les grottes de Mogao. Le mot chinois pour « patrimoine culturel » est \zh{文化遗产} (wénhuà yíchǎn) : c'est exactement le même mot \zh{遗产} que pour l'héritage familial ! Pour les Chinois, la culture d'un peuple se « lègue » comme une maison ou une terre. Et au Cameroun ? Les palais des sultans de l'Ouest, les danses traditionnelles bamiléké ou les paysages culturels du pays sont eux aussi un \zh{遗产} précieux à \zh{传下去}.
\end{savaistu}

\courssub{Exercice d'entraînement et mini-débat}

\begin{exercice}[Collocations et compréhension]
\textbf{A. Complétez avec la bonne collocation} (\zh{留下遗产}, \zh{继承遗产}, \zh{立遗嘱}, \zh{分财产}, \zh{传下去}) :\\
1) 爷爷生前\blank{}\blank{}\blank{}，把房子给了大儿子。\\
2) 老人去世后，孩子们要\blank{}\blank{}\blank{}。\\
3) 爸爸\blank{}\blank{}\blank{}了爷爷的房子和地。\\
4) 传统文化要一代一代\blank{}\blank{}\blank{}。\\
\textbf{B. Répondez en phrases complètes} : 1) 奶奶认为最重要的遗产是什么？ 2) 为什么钱不是最重要的遗产？
\end{exercice}

\begin{corrige}[Exercice « Collocations et compréhension »]
\textbf{A.} 1) \zh{立遗嘱} (lì yízhǔ) : grand-père a rédigé un testament de son vivant. 2) \zh{分财产} (fēn cáichǎn) : après le décès, les enfants partagent les biens. 3) \zh{继承} (jìchéng) : papa a hérité de la maison et de la terre. 4) \zh{传下去} (chuán xiàqù) : la culture traditionnelle doit être transmise de génération en génération.\\
\textbf{B.} 1) \zh{奶奶认为最重要的遗产不是钱，而是爷爷的书法。} 2) \zh{因为钱会用完，可是文化不会。}
\end{corrige}

\begin{experience}[Mini-débat guidé : \zh{钱重要还是文化重要？}]
Par groupes de quatre, préparez un mini-débat en cinq minutes. Deux élèves défendent l'argent (\zh{钱重要}), deux défendent la culture (\zh{文化重要}). Utilisez les structures de la leçon :\\
— \zh{我认为…} (wǒ rènwéi…, je pense que…)\\
— \zh{…比…更…} (… bǐ … gèng …, … est plus … que …)\\
— \zh{不是…而是…} (bú shì… ér shì…)\\
— \zh{因为…所以…} (yīnwèi… suǒyǐ…, parce que… donc…)\\[4pt]
Exemples d'arguments : \zh{我认为文化比钱更重要，因为钱会用完，文化不会。} / \zh{我觉得钱也重要，因为没有钱，孩子们不能上学。}\\[4pt]
Pour enrichir le débat, pensez au Cameroun : les palais des sultans de l'Ouest, les traditions bamiléké, la langue et les danses de votre village sont-ils un \zh{遗产} à protéger ? Que vaut-il mieux laisser à ses enfants : \zh{钱} ou \zh{文化} ?
\end{experience}

\begin{aretenir}[À retenir de cette section]
\begin{itemize}
\item Pour répondre à une question de compréhension : reprendre les mots de la question et remplacer le mot interrogatif par la réponse.
\item En chinois, \zh{什么}, \zh{谁}, \zh{什么时候} restent à la place du mot remplacé, souvent en fin de phrase — jamais au début comme en français.
\item La structure \zh{不是…而是…} corrige une idée : « ce n'est pas A, mais B ».
\item Les collocations à connaître par cœur : \zh{留下遗产}, \zh{继承遗产}, \zh{立遗嘱}, \zh{分财产}, \zh{传下去}.
\item \zh{遗产} (héritage reçu après un décès) ≠ \zh{财产} (biens à partager).
\end{itemize}
\end{aretenir}

\coursec{Production guidée et réinvestissement}

Après avoir vérifié votre compréhension du texte sur l'héritage du grand-père et mémorisé les collocations essentielles (\zh{留下遗产} « laisser un héritage », \zh{宝贵的财富} « une richesse précieuse », \zh{传给下一代} « transmettre à la génération suivante »), il est temps de passer à l'acte : \textbf{vous allez produire vous-mêmes}. C'est l'étape où le vocabulaire devient vraiment le vôtre. Nous procéderons en trois temps : une production écrite guidée à partir d'une trame, une présentation orale d'une minute, puis une reformulation du texte support à la troisième personne. Nous terminerons par une correction collective et une carte mentale de synthèse.

\courssub{Le modèle de paragraphe guidé : parler de l'héritage dans sa famille}

\textbf{Observer d'abord : un paragraphe modèle}\par

Avant d'écrire, lisez ce paragraphe rédigé par un élève de Terminale à Yaoundé. Observez comment chaque phrase remplit une fonction précise.

\begin{textebox}[Modèle de production : l'héritage de ma grand-mère]
我的奶奶去年去世了。\\
Wǒ de nǎinai qùnián qùshì le.\\
\emph{Ma grand-mère est décédée l'année dernière.}\\[4pt]
但是，我觉得最重要的遗产不是这些东西，而是她教给我们的传统。\\
Dànshì, wǒ juéde zuì zhòngyào de yíchǎn bú shì zhèxiē dōngxi, ér shì tā jiāo gěi wǒmen de chuántǒng.\\
\emph{Mais je pense que l'héritage le plus important, ce ne sont pas ces choses, mais les traditions qu'elle nous a enseignées.}\\[4pt]
她常常说：「家庭是最重要的。」这句话是宝贵的财富。\\
Tā chángcháng shuō : « Jiātíng shì zuì zhòngyào de. » Zhè jù huà shì bǎoguì de cáifù.\\
\emph{Elle disait souvent : « La famille est ce qu'il y a de plus important. » Cette phrase est une richesse précieuse.}\\[4pt]
我想把这个遗产传给我的孩子。\\
Wǒ xiǎng bǎ zhège yíchǎn chuán gěi wǒ de háizi.\\
\emph{Je veux transmettre cet héritage à mes enfants.}
\end{textebox}

Observez la construction : le paragraphe suit toujours le même chemin en quatre étapes — \textbf{qui} est décédé, \textbf{quels biens} ont été laissés, quel est l'\textbf{héritage immatériel}, et quelle est l'\textbf{opinion personnelle} de l'auteur. C'est exactement la trame que vous allez utiliser.

\begin{popfigure}
\begin{tikzpicture}[node distance=0.5cm and 0.9cm,
  etape/.style={rectangle, rounded corners=4pt, draw=popblue, fill=popblueL, align=center, minimum width=3.4cm, minimum height=1.1cm, font=\small},
  fleche/.style={-{Stealth[length=3mm]}, thick, poporange}]
\node[etape, align=center] (e1){\textbf{1. Introduction}\\\zh{我的爷爷去世了。}\\qui ?};
\node[etape, right=of e1, draw=poppurple, fill=poppurpleL, align=center] (e2){\textbf{2. Biens laissés}\\\zh{他留下了……}\\quoi ?};
\node[etape, below=0.6cm of e2, draw=popgreen, fill=popgreenL, align=center] (e3){\textbf{3. Héritage immatériel}\\\zh{最重要的遗产是……}\\traditions, valeurs};
\node[etape, left=of e3, draw=popgold, fill=popgoldL, align=center] (e4){\textbf{4. Opinion personnelle}\\\zh{我觉得……，因为……}\\mon avis};
\draw[fleche] (e1) -- (e2);
\draw[fleche] (e2) -- (e3);
\draw[fleche] (e3) -- (e4);
\end{tikzpicture}
\legende{La trame de production en quatre étapes : de la personne aux biens, puis à l'héritage immatériel, et enfin à votre opinion.}
\end{popfigure}

\textbf{La trame à trous}\par

Voici la trame officielle de cette production. Les blancs signalés par \zh{＿＿＿} sont les seuls éléments que vous devez remplacer dans un premier temps.

\begin{aretenir}[La trame de production guidée]
\begin{enumerate}
\item \zh{我的＿＿＿去世了。} (Wǒ de \_\_\_ qùshì le.) — Mon/ma \_\_\_ est décédé(e). \emph{[qui]}
\item \zh{他/她留下了＿＿＿。} (Tā liúxià le \_\_\_.) — Il/elle a laissé \_\_\_. \emph{[biens]}
\item \zh{我觉得最重要的遗产是＿＿＿，因为＿＿＿。} (Wǒ juéde zuì zhòngyào de yíchǎn shì \_\_\_, yīnwèi \_\_\_.) — Je pense que l'héritage le plus important est \_\_\_, parce que \_\_\_. \emph{[opinion + justification]}
\end{enumerate}
Pour atteindre 5 ou 6 phrases, ajoutez une citation de la personne (\zh{他常常说：「……」}) et une phrase de conclusion (\zh{我想把这个遗产传给……}).
\end{aretenir}

\begin{methode}[Comment réussir votre production guidée en quatre pas]
\begin{enumerate}
\item \textbf{Recopiez la trame} intégralement, caractère par caractère, sans rien modifier de ce qui n'est pas un blanc. La trame est grammaticalement correcte : toute modification risque de créer une faute.
\item \textbf{Remplacez uniquement les blancs} avec le vocabulaire du glossaire de la leçon : \zh{爷爷} (yéye) grand-père, \zh{奶奶} (nǎinai) grand-mère, \zh{父亲} (fùqīn) père, \zh{母亲} (mǔqīn) mère ; \zh{一所房子} (yí suǒ fángzi) une maison, \zh{一些钱} (yìxiē qián) de l'argent, \zh{一块地} (yí kuài dì) un terrain ; \zh{传统} (chuántǒng) traditions, \zh{文化} (wénhuà) culture, \zh{语言} (yǔyán) langue.
\item \textbf{Enrichissez avec un adjectif} du glossaire : placez \zh{宝贵的} (bǎoguì de, précieux) ou \zh{重要的} (zhòngyào de, important) devant le nom, sans oublier \zh{的} : \zh{宝贵的遗产} un héritage précieux, \zh{重要的传统} une tradition importante.
\item \textbf{Relisez à voix haute} en vérifiant trois points : les tons de chaque mot nouveau, la présence de \zh{了} après \zh{去世} et \zh{留下} (action accomplie), et la place de \zh{因为} devant la justification.
\end{enumerate}
\end{methode}

\courssub{La grammaire en action : 觉得 pour exprimer son opinion}

Dans la trame, la phrase clé est \zh{我觉得……}. C'est l'outil indispensable pour donner un avis personnel en chinois, et les francophones font souvent des erreurs ici.

\begin{definition}[觉得 juéde]
Verbe modal d'opinion : « penser, trouver, avoir l'impression que ». Il introduit un jugement personnel, toujours suivi d'une proposition complète (sujet + verbe ou adjectif). Structure : \textbf{Sujet + 觉得 + proposition}. Exemple : \zh{我觉得传统很重要。} (Wǒ juéde chuántǒng hěn zhòngyào.) « Je trouve que les traditions sont importantes. »
\end{definition}

\begin{exemplebox}[Donner son avis avec 觉得]
\zh{我觉得最重要的遗产是我们的传统。}\\
\emph{Wǒ juéde zuì zhòngyào de yíchǎn shì wǒmen de chuántǒng.}\\
« Je pense que l'héritage le plus important, ce sont nos traditions. »\\[4pt]
\zh{他觉得家庭的语言很宝贵。} \emph{(Tā juéde jiātíng de yǔyán hěn bǎoguì.)} « Il trouve que la langue de la famille est précieuse. »\\[4pt]
\zh{我们都觉得爷爷的话是宝贵的财富。} \emph{(Wǒmen dōu juéde yéye de huà shì bǎoguì de cáifù.)} « Nous pensons tous que les paroles de grand-père sont une richesse précieuse. »
\end{exemplebox}

\begin{attention}[Deux pièges pour les francophones]
\begin{itemize}
\item \textbf{Ne pas utiliser \zh{是} seul pour donner un avis.} En français on dit « c'est important », mais \zh{是重要} est incorrect : devant un adjectif, le chinois exige un adverbe de degré, le plus souvent \zh{很} : \zh{遗产很重要} (et non \zh{遗产是重要}). \zh{是} ne s'emploie que devant un \textbf{nom} : \zh{遗产是传统} « l'héritage, ce sont les traditions ».
\item \textbf{Ne pas traduire « je pense que » par \zh{我想} suivi d'une opinion.} \zh{我想} signifie « je veux / j'ai l'intention de » (\zh{我想去中国} « je veux aller en Chine »). Pour une opinion, utilisez \zh{我觉得} ou \zh{我认为} (plus formel).
\end{itemize}
\end{attention}

\begin{center}
\begin{tabularx}{\linewidth}{|X|X|X|}
\hline
\rowcolor{popblueL} \textbf{Français} & \textbf{Chinois correct} & \textbf{Faute fréquente} \\
\hline
Je pense que c'est important. & \zh{我觉得这很重要。} & \zh{我觉得这是重要。} \\
\hline
L'héritage est précieux. & \zh{遗产很宝贵。} & \zh{遗产是宝贵。} \\
\hline
L'héritage, c'est la culture. & \zh{遗产是文化。} & \zh{遗产很文化。} \\
\hline
Je trouve que les traditions sont importantes. & \zh{我觉得传统很重要。} & \zh{我想传统很重要。} \\
\hline
\end{tabularx}
\end{center}

\courssub{Phrase d'ouverture pour l'examen}

\begin{textebox}[Phrase d'ouverture utile à l'examen]
每个家庭都有自己的遗产。\\
Měi ge jiātíng dōu yǒu zìjǐ de yíchǎn.\\
\emph{Chaque famille a son propre héritage.}\\[4pt]
有的家庭留下房子和钱，有的家庭留下传统和故事。\\
Yǒu de jiātíng liúxià fángzi hé qián, yǒu de jiātíng liúxià chuántǒng hé gùshi.\\
\emph{Certaines familles laissent des maisons et de l'argent, d'autres laissent des traditions et des histoires.}
\end{textebox}

Cette phrase d'ouverture est un excellent début de rédaction ou d'exposé : elle est générale, correcte, et elle annonce naturellement votre exemple personnel. La structure \zh{有的……，有的……} (« les uns…, les autres… ») est très appréciée à l'examen car elle montre que vous savez construire une opposition élégante.

\courssub{Exercice résolu : compléter la trame}

\begin{exoresolu}[Production guidée de Marie, élève à Douala]
Énoncé : Marie complète la trame avec les éléments suivants : sa grand-mère (\zh{奶奶}), une maison et des histoires (\zh{一所房子和很多故事}), la langue de sa famille (\zh{我们家的语言}), parce qu'elle unit la famille (\zh{因为它让我们一家人在一起}). Rédigez le paragraphe complet.
\tcblower
Solution détaillée :
\begin{enumerate}
\item \zh{我的奶奶去世了。} (Wǒ de nǎinai qùshì le.) — On place le possessif \zh{我的} devant \zh{奶奶} et \zh{了} après \zh{去世}.
\item \zh{她留下了一所房子和很多故事。} (Tā liúxià le yì suǒ fángzi hé hěn duō gùshi.) — Attention au classificateur : \zh{一所房子} avec \zh{所} pour les bâtiments ; \zh{很多} devant \zh{故事}.
\item \zh{我觉得最重要的遗产是我们家的语言，因为它让我们一家人在一起。} (Wǒ juéde zuì zhòngyào de yíchǎn shì wǒmen jiā de yǔyán, yīnwèi tā ràng wǒmen yì jiā rén zài yìqǐ.) — \zh{觉得} + proposition complète ; \zh{因为} introduit la justification.
\end{enumerate}
Vérification : tons corrects (\zh{nǎinai} deux syllabes en 3\textsuperscript{e} ton dont la seconde neutre), \zh{的} présent dans \zh{最重要的遗产} et \zh{我们家的语言}, collocations respectées (\zh{留下}, \zh{最重要的遗产}).
\end{exoresolu}

\courssub{Variante orale : la présentation d'une minute}

\begin{experience}[Mon héritage familial en une minute]
Par binôme, puis devant la classe :
\begin{enumerate}
\item \textbf{Préparation (5 min)} : chaque élève rédige son paragraphe guidé (5 à 6 phrases) à partir de la trame, puis le lit à voix haute à son voisin qui vérifie les tons.
\item \textbf{Présentation (1 min)} : l'élève se lève, commence par la phrase d'ouverture \zh{每个家庭都有自己的遗产。}, puis enchaîne son paragraphe sans lire mot à mot (un plan de quatre mots-clés sur un papier est autorisé : \zh{谁 — 留下 — 遗产 — 觉得}).
\item \textbf{Écoute active} : la classe note pour chaque exposé une collocation correcte entendue (\zh{留下}, \zh{宝贵的}, \zh{传给}…) et une question à poser en chinois : \zh{你觉得什么最重要？} « Qu'est-ce qui est le plus important pour toi ? »
\item \textbf{Critères d'évaluation} : exactitude des tons (surtout \zh{yíchǎn}, \zh{liúxià}, \zh{juéde}), présence de \zh{了}, fluidité, regard vers l'auditoire.
\end{enumerate}
\end{experience}

\courssub{Reformulation du texte support à la troisième personne}

Transformer un texte de la première à la troisième personne est un exercice classique de l'examen (version et reprise de texte). La règle est simple mais exige de la rigueur.

\begin{propriete}[Passage de la 1\textsuperscript{re} à la 3\textsuperscript{e} personne]
\begin{itemize}
\item \zh{我} (wǒ, je) devient \zh{他} (tā, il) ou \zh{她} (tā, elle) selon le narrateur ;
\item \zh{我的爷爷} (mon grand-père) devient \zh{他的爷爷} (son grand-père) ;
\item les verbes et les collocations ne changent \textbf{pas} : \zh{留下了}, \zh{觉得}, \zh{传给} restent identiques ;
\item les citations directes peuvent être conservées entre guillemets « 》 ou rapportées.
\end{itemize}
Exemple : \zh{我的爷爷留下了一所房子。} → \zh{他的爷爷留下了一所房子。} « Son grand-père a laissé une maison. »
\end{propriete}

\begin{exercice}[À vous de reformuler]
Reformulez à la troisième personne (narrateur masculin) : \zh{我觉得最重要的遗产是我们家的传统。我想把它传给我的孩子。}
\end{exercice}

\begin{corrige}[Exercice de reformulation]
\zh{他觉得最重要的遗产是他们家的传统。他想把它传给自己的孩子。}\\
\emph{Tā juéde zuì zhòngyào de yíchǎn shì tāmen jiā de chuántǒng. Tā xiǎng bǎ tā chuán gěi zìjǐ de háizi.}\\
« Il pense que l'héritage le plus important, ce sont les traditions de sa famille. Il veut les transmettre à ses propres enfants. »\\
Points de vigilance : \zh{我} → \zh{他} (deux fois) ; \zh{我们家} → \zh{他们家} ; \zh{我的孩子} → \zh{自己的孩子} (\zh{自己} « soi-même » évite la répétition de \zh{他的}).
\end{corrige}

\courssub{Correction collective : la grille de vérification}

Après la production, la classe corrige ensemble deux ou trois paragraphes écrits au tableau. La correction porte toujours sur les mêmes trois points, dans cet ordre :

\begin{center}
\begin{tabularx}{\linewidth}{|X|X|X|}
\hline
\rowcolor{popgreenL} \textbf{Point de contrôle} & \textbf{Question à se poser} & \textbf{Exemple corrigé} \\
\hline
Les tons & Chaque mot nouveau est-il prononcé avec les bons tons ? & \zh{yíchǎn} (2\textsuperscript{e} + 3\textsuperscript{e} ton), non \emph{yìchǎn} \\
\hline
La particule \zh{的} & Y a-t-il \zh{的} entre l'adjectif ou le possessif et le nom ? & \zh{宝贵的遗产}, \zh{我们家的语言} \\
\hline
Les collocations & Les verbes vont-ils avec les bons noms ? & \zh{留下遗产}, \zh{传给下一代}, \zh{学习传统} \\
\hline
\end{tabularx}
\end{center}

\begin{attention}[Erreurs relevées fréquemment en classe]
\begin{itemize}
\item \zh{我的爷爷去世} sans \zh{了} : l'action est accomplie, \zh{了} est obligatoire → \zh{去世了}.
\item \zh{最重要遗产} sans \zh{的} : le superlatif \zh{最 + adjectif} exige \zh{的} devant le nom → \zh{最重要的遗产}.
\item \zh{他留下了一个遗产} : calque du français « il a laissé un héritage » ; en chinois on dit simplement \zh{他留下了遗产} ou on précise le bien : \zh{他留下了一所房子}.
\end{itemize}
\end{attention}

\courssub{Synthèse de la leçon : la carte mentale à compléter}

Pour clore la leçon, complétez sur votre cahier la carte mentale suivante en ajoutant, à chaque branche, deux mots du glossaire et une phrase complète de votre production.

\begin{popfigure}
\begin{tikzpicture}[mindmap, grow cyclic, every node/.style={concept, align=center, font=\small},
  root concept/.append style={concept color=popblue, fill=popblueL, font=\small\bfseries},
  level 1/.append style={level distance=3.2cm, sibling angle=90},
  level 1 concept/.append style={concept color=poporange, fill=poporangeL}]
\node[root concept, align=center]{\zh{遗产}\\héritage}
  child { node[align=center]{\zh{人}\\les personnes\\\zh{爷爷 奶奶}\\……} }
  child { node[align=center]{\zh{留下}\\les biens\\\zh{房子 钱}\\……} }
  child { node[align=center]{\zh{传统}\\l'immatériel\\\zh{文化 语言}\\……} }
  child { node[align=center]{\zh{我觉得}\\mon avis\\\zh{宝贵 重要}\\……} };
\end{tikzpicture}
\legende{Carte mentale de synthèse : complétez chaque branche avec deux mots du glossaire et une phrase personnelle.}
\end{popfigure}

\begin{aretenir}[L'essentiel de la leçon]
\begin{itemize}
\item Une production guidée réussie = trame recopiée fidèlement + blancs remplis avec le glossaire + un adjectif enrichisseur (\zh{宝贵的}, \zh{重要的}).
\item Pour donner votre avis : \zh{我觉得 + proposition} ; jamais \zh{是} seul devant un adjectif, jamais \zh{我想} pour une opinion.
\item À l'oral comme à l'écrit, ouvrez avec \zh{每个家庭都有自己的遗产。} et concluez avec \zh{我想把这个遗产传给……}.
\item Trois contrôles systématiques : les tons, la particule \zh{的}, les collocations \zh{留下 / 传给 / 学习}.
\end{itemize}
\end{aretenir}

\newpage

\coursec{Activité d'intégration}

\courssub{Objectifs de l'activité}

À l'issue de cette activité d'intégration, l'élève sera capable de :

\begin{itemize}
\item mobiliser dans une production personnelle le \cle{vocabulaire de l'héritage} étudié dans la leçon : \zh{遗产} (yíchǎn, héritage), \zh{继承} (jìchéng, hériter de), \zh{留下} (liúxià, laisser en héritage), \zh{传下去} (chuán xiàqu, transmettre aux générations suivantes), \zh{传统} (chuántǒng, tradition), \zh{祖先} (zǔxiān, ancêtres) ;
\item employer correctement les \cle{collocations} verbales du thème : \zh{留下遗产} (laisser un héritage), \zh{继承传统} (hériter d'une tradition), \zh{把文化传下去} (transmettre la culture) ;
\item utiliser la structure de contraste \cle{不是……而是……} (bú shì... ér shì..., « ce n'est pas... mais plutôt... ») pour opposer héritage matériel et héritage culturel ;
\item rédiger un texte cohérent de 6 à 8 phrases en caractères simplifiés correctement tracés ;
\item présenter oralement son texte avec une prononciation soignée des tons et répondre à des questions de compréhension en chinois.
\end{itemize}

\courssub{Situation}

Votre lycée de Yaoundé participe à la \cle{Semaine du patrimoine culturel} organisée en partenariat avec l'Institut Confucius. Le thème de cette année est : « Ce que nos familles nous transmettent ». Chaque classe de Terminale doit préparer un stand d'exposition : un court texte rédigé en chinois, présenté oralement devant un jury composé d'élèves et d'enseignants de chinois. Le meilleur texte sera publié sur le site web du lycée.

Pour vous inspirer, le professeur vous distribue le modèle suivant, rédigé par une ancienne élève :

\begin{textebox}[Modèle d'exposé : « L'héritage de ma grand-mère »]
我奶奶去世的时候，没有留下很多钱，而是留下了更宝贵的东西。\\
\emph{Wǒ nǎinai qùshì de shíhou, méiyǒu liúxià hěn duō qián, ér shì liúxià le gèng bǎoguì de dōngxi.}\\
« Quand ma grand-mère est décédée, elle n'a pas laissé beaucoup d'argent, mais elle a laissé quelque chose de plus précieux. »\\[4pt]
她教会我们全家做传统的菜，还给我们讲了很多祖先的故事。\\
\emph{Tā jiāohuì wǒmen quánjiā zuò chuántǒng de cài, hái gěi wǒmen jiǎng le hěn duō zǔxiān de gùshi.}\\
« Elle a appris à toute notre famille à préparer les plats traditionnels, et elle nous a raconté beaucoup d'histoires de nos ancêtres. »\\[4pt]
我们家最重要的遗产不是房子，而是这些故事和菜。\\
\emph{Wǒmen jiā zuì zhòngyào de yíchǎn bú shì fángzi, ér shì zhèxiē gùshi hé cài.}\\
« L'héritage le plus important de notre famille, ce n'est pas la maison, mais ces histoires et ces plats. »\\[4pt]
我要把这个传统传下去，让我的孩子也继承它。\\
\emph{Wǒ yào bǎ zhège chuántǒng chuán xiàqu, ràng wǒ de háizi yě jìchéng tā.}\\
« Je veux transmettre cette tradition, afin que mes enfants l'héritent aussi. »
\end{textebox}

\courssub{Consignes}

Travail en groupes de trois élèves. Suivez les étapes dans l'ordre.

\begin{enumerate}
\item \textbf{Compréhension du modèle.} Lisez le texte modèle à voix haute à tour de rôle, puis répondez par écrit en chinois aux deux questions suivantes : \zh{奶奶留下了什么？} (Nǎinai liúxià le shénme ? — Qu'a laissé la grand-mère ?) et \zh{为什么她说遗产不是房子？} (Wèishénme tā shuō yíchǎn bú shì fángzi ? — Pourquoi dit-elle que l'héritage n'est pas la maison ?).
\item \textbf{Recherche de lexique.} Soulignez dans le modèle les trois verbes \zh{留下}, \zh{继承} et \zh{传下去}, et recopiez chaque collocation complète (verbe + complément). Ajoutez deux autres mots du glossaire de la leçon que vous voulez utiliser.
\item \textbf{Rédaction guidée.} Rédigez ensemble un texte de 6 à 8 phrases sur l'héritage d'une famille (réelle ou imaginaire, camerounaise ou chinoise) en suivant cette trame :
\begin{itemize}
\item phrase 1 : présenter la personne qui a transmis l'héritage (\zh{我的……});
\item phrases 2-3 : décrire l'héritage matériel avec \zh{留下};
\item phrases 4-5 : décrire l'héritage culturel (langue, plats, contes, musique, valeurs) avec \zh{继承} ou \zh{教会};
\item phrase 6 : une phrase de contraste avec \zh{不是……而是……};
\item phrases 7-8 : conclure avec \zh{把……传下去} et un souhait pour l'avenir.
\end{itemize}
\item \textbf{Vérification de l'écrit.} Relisez votre texte : vérifiez l'ordre des traits et les clés des caractères du thème (\zh{遗}, \zh{继}, \zh{留}, \zh{传}), la place des compléments et la ponctuation chinoise.
\item \textbf{Présentation orale.} Chaque groupe présente son texte devant la classe en 2 minutes maximum. Un membre écrit au tableau les deux phrases-clés du groupe.
\item \textbf{Écoute active.} Pendant chaque présentation, les autres groupes préparent deux questions de compréhension commençant par \zh{什么} ou \zh{为什么}, et les posent au groupe qui vient de parler.
\end{enumerate}

\courssub{Ressources à mobiliser}

\begin{aretenir}[La structure 不是……而是……]
Structure : Sujet + \zh{不是} + A + \zh{而是} + B. Elle sert à corriger ou à opposer deux idées : « ce n'est pas A, mais plutôt B ». Le second élément est celui que le locuteur juge vrai ou plus important.\\
Exemple : \zh{最宝贵的遗产不是钱，而是爱。} \emph{Zuì bǎoguì de yíchǎn bú shì qián, ér shì ài.} « L'héritage le plus précieux, ce n'est pas l'argent, mais l'amour. »
\end{aretenir}

\begin{center}
\begin{tabularx}{\linewidth}{|X|X|X|X|}
\hline
\rowcolor{popblueL} Caractères & Pinyin & Nature & Traduction \\
\hline
遗产 & yíchǎn & n. & héritage, patrimoine \\
\hline
留下 & liúxià & v. & laisser (en héritage) \\
\hline
继承 & jìchéng & v. & hériter de, perpétuer \\
\hline
传下去 & chuán xiàqu & v. + compl. & transmettre (aux générations suivantes) \\
\hline
传统 & chuántǒng & n./adj. & tradition ; traditionnel \\
\hline
祖先 & zǔxiān & n. & ancêtres \\
\hline
宝贵 & bǎoguì & adj. & précieux \\
\hline
教会 & jiāohuì & v. & enseigner, apprendre à (qqn) \\
\hline
\end{tabularx}
\end{center}

\begin{attention}[Erreurs fréquentes à éviter]
\begin{itemize}
\item Ne dites pas \zh{是……而是……} : la structure exige la négation \zh{不是} dans la première partie.
\item \zh{传下去} se construit souvent avec \zh{把} : \zh{把传统传下去} (bǎ chuántǒng chuán xiàqu), et non \zh{传下去传统} à la fin de phrase.
\item Attention aux tons de \zh{继承} (jìchéng, 4\up{e} + 2\up{e} ton) : ne pas confondre avec \zh{机场} (jīchǎng, aéroport).
\item Le caractère \zh{遗} contient la clé du mouvement \zh{辶}, qui se trace en dernier ; \zh{继} commence par la clé de la soie \zh{纟} à gauche.
\end{itemize}
\end{attention}

\courssub{Proposition de production et corrigé commenté}

\begin{exoresolu}[Production modèle : « L'héritage de mon grand-père »]
Rédigez un texte de 6 à 8 phrases sur l'héritage matériel et culturel d'une famille, en utilisant au moins une fois \zh{留下}, \zh{继承}, \zh{传下去} et la structure \zh{不是……而是……}.
\tcblower
\textbf{Texte modèle :}\\[4pt]
我爷爷是杜阿拉的一个农民。\\
\emph{Wǒ yéye shì Dù'ālā de yí ge nóngmín.} « Mon grand-père était un paysan de Douala. »\\[3pt]
他去世以前，给我们留下了一小块地和一所老房子。\\
\emph{Tā qùshì yǐqián, gěi wǒmen liúxià le yì xiǎo kuài dì hé yì suǒ lǎo fángzi.} « Avant de mourir, il nous a laissé un petit lopin de terre et une vieille maison. »\\[3pt]
他还教会我们说我们民族的语言，唱传统的歌。\\
\emph{Tā hái jiāohuì wǒmen shuō wǒmen mínzú de yǔyán, chàng chuántǒng de gē.} « Il nous a aussi appris à parler la langue de notre ethnie et à chanter les chants traditionnels. »\\[3pt]
对我来说，最宝贵的遗产不是那块地，而是我们的语言和文化。\\
\emph{Duì wǒ láishuō, zuì bǎoguì de yíchǎn bú shì nà kuài dì, ér shì wǒmen de yǔyán hé wénhuà.} « Pour moi, l'héritage le plus précieux, ce n'est pas ce lopin de terre, mais notre langue et notre culture. »\\[3pt]
房子会老，钱会用完，可是祖先的故事不会消失。\\
\emph{Fángzi huì lǎo, qián huì yòngwán, kěshì zǔxiān de gùshi bú huì xiāoshī.} « La maison vieillira, l'argent s'épuisera, mais les histoires des ancêtres ne disparaîtront pas. »\\[3pt]
我要把这些故事写下来，把爷爷的传统传下去。\\
\emph{Wǒ yào bǎ zhèxiē gùshi xiě xiàlai, bǎ yéye de chuántǒng chuán xiàqu.} « Je veux écrire ces histoires et transmettre la tradition de mon grand-père. »\\[3pt]
我希望我的孩子也能继承这个宝贵的遗产。\\
\emph{Wǒ xīwàng wǒ de háizi yě néng jìchéng zhège bǎoguì de yíchǎn.} « J'espère que mes enfants pourront aussi hériter de ce précieux patrimoine. »\\[6pt]
\textbf{Commentaire des choix de langue :}
\begin{itemize}
\item La phrase 2 utilise \zh{留下} avec deux compléments introduits par des classificateurs corrects : \zh{一块地} (une parcelle de terre) et \zh{一所房子} (une maison) — vérifier les classificateurs est un réflexe essentiel.
\item La phrase 4 applique la structure \zh{不是……而是……} avec la négation obligatoire ; l'expression \zh{对我来说} (« pour moi ») introduit naturellement l'opinion personnelle.
\item La phrase 6 montre l'emploi correct de \zh{把} + objet + \zh{传下去}, conformément à la règle : l'objet (\zh{爷爷的传统}) est placé avant le verbe grâce à \zh{把}.
\item Le verbe \zh{继承} clôt le texte avec un complément nominal complet (\zh{这个宝贵的遗产}), ce qui évite la phrase incomplète \zh{继承它}, moins naturelle ici.
\item Le vocabulaire reste de niveau HSK 3-4 : \zh{农民}, \zh{民族}, \zh{消失}, \zh{希望} sont des mots de fréquence élevée adaptés à la Terminale.
\end{itemize}
\end{exoresolu}

\courssub{Bilan de l'activité}

Cette activité a permis de réinvestir l'ensemble des savoir-faire de la leçon dans une tâche de communication authentique :

\begin{itemize}
\item \textbf{Compréhension de l'écrit} : lecture et exploitation du texte modèle, repérage des collocations et de la structure de contraste ;
\item \textbf{Lexique} : emploi actif des mots de l'héritage (\zh{遗产}, \zh{留下}, \zh{继承}, \zh{传下去}, \zh{传统}, \zh{祖先}) dans un contexte personnel ;
\item \textbf{Grammaire} : production de phrases avec \zh{不是……而是……} et avec la construction en \zh{把};
\item \textbf{Écriture des caractères} : attention portée aux clés (\zh{辶}, \zh{纟}, \zh{田}) et à l'ordre des traits lors de la mise au tableau ;
\item \textbf{Expression orale} : présentation publique avec soin des tons, écoute active et interaction par questions en \zh{什么} et \zh{为什么}.
\end{itemize}

L'enseignant évalue chaque groupe selon quatre critères : la justesse de la prononciation des tons, l'emploi correct des trois collocations imposées, la présence et la correction de la structure \zh{不是……而是……}, et l'exactitude des caractères écrits au tableau. Le texte gagnant, choisi par vote de la classe, sera affiché au stand de la Semaine du patrimoine : une belle manière de montrer que, comme le dit le proverbe chinois \zh{前人栽树，后人乘凉} (qiánrén zāi shù, hòurén chéngliáng — « les générations précédentes plantent les arbres, les générations suivantes profitent de leur ombre »), l'héritage se reçoit et se transmet.

\newpage

\coursec{Méthodes et exercices résolus}

\coursec{Leçon 5 — Vocabulaire et compréhension de texte : héritage}

\courssub{Mise en route : l'héritage entre Cameroun et Chine}

\begin{textebox}[Situation de communication : discussion familiale à Yaoundé]
\zh{—— 小李，你爷爷不是住在杜阿拉郊区吗？} \\
\emph{—— Xiǎo Lǐ, nǐ yéye bù shì zhù zài Dùālā jiāoqū ma?} \\
« — Petit Li, ton grand-père n'habite-t-il pas en banlieue de Douala ? »

\zh{—— 是的。他前年把在家乡的一块地和做木雕的手艺都传给了我爸爸。} \\
\emph{—— Shì de. Tā qiánnián bǎ zài jiāxiāng de yí kuài dì hé zuò mùdiāo de shǒuyì dōu chuán gěi le wǒ bàba.} \\
« — Oui. L'année d'avant, il a transmis à mon père un terrain au village natal et son savoir-faire de sculpture sur bois. »

\zh{—— 这真是宝贵的遗产啊！不仅是财产，更是文化。} \\
\emph{—— Zhè zhēn shì bǎoguì de yíchǎn a! Bù jǐn shì cáichǎn, gèng shì wénhuà.} \\
« — C'est vraiment un héritage précieux ! Non seulement un bien matériel, mais surtout la culture. »
\end{textebox}

\begin{savaistu}[Héritage : regards croisés Cameroun / Chine]
Au \textbf{Cameroun}, l'héritage coutumier (\emph{coutume}) prime souvent sur le droit écrit : la terre appartient à la lignée, le chef de famille gère le \cle{foncier} ; la transmission orale (généalogies, proverbes, danses, recettes de \zh{ndolé} \emph{ndolé}) est un devoir sacré. En \textbf{Chine}, la \cle{piété filiale} (\zh{孝 xiào}) structure l'héritage : le fils aîné avait historiquement la part principale (\zh{嫡长子继承制}), mais le Code civil actuel (2021) garantit l'égalité entre héritiers ; les \cle{biens culturels immatériels} (\zh{非物质文化遗产 fēiwùzhì wénhuà yíchǎn}) — calligraphie, médecine traditionnelle, opéra de Pékin — sont protégés par l'État comme « héritage de la nation ». Dans les deux pays, \textbf{la terre et le savoir-faire} sont perçus comme l'âme de la famille, bien au-delà de leur valeur marchande.
\end{savaistu}

\courssub{Texte support : lecture, pinyin, traduction}

\begin{textebox}[Texte original : « La promesse au grand-père » — Niveau HSK 3-4]
\zh{我的爷爷住在雅温得郊区的一个小村庄里。他今年七十五岁了，身体还很硬朗。} \\
\emph{Wǒ de yéye zhù zài Yǎwēndé jiāoqū de yí gè xiǎo cūnzhuāng lǐ. Tā jīnnián qīshíwǔ suì le, shēntǐ hái hěn yìnglǎng.} \\
« Mon grand-père habite dans un petit village en périphérie de Yaoundé. Il a 75 ans cette année, et sa santé est encore robuste. »

\zh{爷爷常常对我说：「孩子，钱花完了还可以赚，但我们祖先留下的传统和手艺可不能丢。」} \\
\emph{Yéye chángcháng duì wǒ shuō : « Háizi, qián huā wán le hái kěyǐ zhuàn, dàn wǒmen zǔxiān liú xià de chuántǒng hé shǒuyì kě bù néng diū. »} \\
« Grand-père me dit souvent : "Enfant, l'argent dépensé peut se regagner, mais la tradition et le savoir-faire que nos ancêtres nous ont laissés, on ne doit surtout pas les perdre." »

\zh{去年，他决定把家族的历史、做木雕的手艺以及他在杜阿拉附近的一块地传给我爸爸。} \\
\emph{Qùnián, tā juédìng bǎ jiāzú de lìshǐ, zuò mùdiāo de shǒuyì yǐjí tā zài Dùālā fùjìn de yí kuài dì chuán gěi wǒ bàba.} \\
« L'an dernier, il a décidé de transmettre à mon père l'histoire de la famille, l'art de la sculpture sur bois, ainsi qu'un terrain qu'il possède près de Douala. »

\zh{他说：「这块地不只是财产，更是我们家族的根。」} \\
\emph{Tā shuō : « Zhè kuài dì bù zhǐ shì cáichǎn, gèng shì wǒmen jiāzú de gēn. »} \\
« Il a dit : "Ce terrain n'est pas seulement un bien, c'est surtout la racine de notre famille." »

\zh{我答应爷爷，一定要好好继承这份宝贵的遗产，传给下一代。} \\
\emph{Wǒ dāying yéye, yídìng yào hǎohǎo jìchéng zhè fèn bǎoguì de yíchǎn, chuán gěi xià yī dài.} \\
« J'ai promis à grand-père de bien hériter de cet héritage précieux et de le transmettre à la génération suivante. »
\end{textebox}

\begin{methode}[Lire un texte sur l'héritage : la méthode en 5 étapes]
\begin{enumerate}
\item \textbf{Lecture globale silencieuse.} Lis le texte une première fois sans dictionnaire. Repère le thème général (\zh{遗产 yíchǎn} — héritage) et les mots connus (\zh{家、文化、老人、孩子、雅温得、杜阿拉}).
\item \textbf{Repérage des mots-clés du thème.} Souligne le lexique de la transmission : \zh{传给} (transmettre à), \zh{继承} (hériter), \zh{留下} (laisser), \zh{传统} (tradition), \zh{手艺} (savoir-faire), \zh{祖先} (ancêtres), \zh{财产} (biens), \zh{宝贵} (précieux).
\item \textbf{Analyse syntaxique (ordre des mots).} Identifie la structure \cle{Sujet + Temps/Lieu + Verbe + Complément} et la structure \cle{\zh{把} (bǎ)} pour l'objet déplacé : \zh{他 把 地 传给 爸爸}.
\item \textbf{Lecture à voix haute avec le pinyin.} Vérifie chaque ton : \zh{yí chuán} (2 puis 2\^{e} ton modifié) vs \zh{yí kuài} ; \zh{bǎoguì} (3-4) ; \zh{jiāzú} (1-2). Une erreur de ton change le sens (\zh{mǎi} 买 ≠ \zh{mài} 卖).
\item \textbf{Reformulation et résumé.} Traduis phrase par phrase, puis résume en français en 2 phrases : Qui ? Qu'est-ce qui est transmis ? Pourquoi est-ce vital pour le grand-père ?
\end{enumerate}
\end{methode}

\courssub{Glossaire du thème : l'héritage (tableau de vocabulaire)}

\begin{center}
\begin{tabularx}{\linewidth}{|X|X|X|X|}
\hline
\rowcolor{popblueL} \textbf{Caractères} & \textbf{Pinyin} & \textbf{Nature} & \textbf{Traduction} \\ \hline
\zh{遗产} & yíchǎn & n. & héritage, legs \\ \hline
\zh{继承} & jìchéng & v. & hériter, succéder à \\ \hline
\zh{传给} & chuán gěi & v. & transmettre à (qn) \\ \hline
\zh{留下} & liú xià & v. & laisser, léguer \\ \hline
\zh{祖先} & zǔxiān & n. & ancêtres, aïeux \\ \hline
\zh{后代} & hòudài & n. & descendance, générations futures \\ \hline
\zh{家族} & jiāzú & n. & clan, grande famille \\ \hline
\zh{传统} & chuántǒng & n. & tradition ; adj. traditionnel \\ \hline
\zh{文化} & wénhuà & n. & culture \\ \hline
\zh{手艺} & shǒuyì & n. & savoir-faire, métier manuel, artisanat \\ \hline
\zh{木雕} & mùdiāo & n. & sculpture sur bois \\ \hline
\zh{财产} & cáichǎn & n. & biens, patrimoine, fortune \\ \hline
\zh{土地} & tǔdì & n. & terre, terrain, foncier \\ \hline
\zh{房子} & fángzi & n. & maison, bâtiment \\ \hline
\zh{根} & gēn & n. & racine ; fig. origine, fondement \\ \hline
\zh{宝贵} & bǎoguì & adj. & précieux, inestimable \\ \hline
\zh{硬朗} & yìnglǎng & adj. & robuste, en bonne santé (personne âgée) \\ \hline
\zh{答应} & dāying & v. & promettre, accepter de \\ \hline
\zh{决定} & juédìng & v./n. & décider / décision \\ \hline
\zh{郊区} & jiāoqū & n. & banlieue, périphérie \\ \hline
\end{tabularx}
\end{center}

\begin{propriete}[Collocations indispensables (mots « collés » à apprendre par cœur)]
\begin{center}
\begin{tabularx}{\linewidth}{|X|X|X|}
\hline
\rowcolor{popgreenL} \textbf{Collocation (chinois + pinyin)} & \textbf{Structure} & \textbf{Traduction} \\ \hline
\zh{继承遗产} (jìchéng yíchǎn) & V + N & hériter d'un héritage \\ \hline
\zh{留下遗产} (liú xià yíchǎn) & V + N & laisser un héritage \\ \hline
\zh{传给后代} (chuán gěi hòudài) & V + N & transmettre à la descendance \\ \hline
\zh{宝贵的财富} (bǎoguì de cáifù) & Adj + \zh{的} + N & un trésor précieux \\ \hline
\zh{一块地} (yí kuài dì) & Num + Class + N & un terrain (classif. \zh{块}) \\ \hline
\zh{一座房子} (yí zuò fángzi) & Num + Class + N & une maison (classif. \zh{座}) \\ \hline
\zh{一门手艺} (yí mén shǒuyì) & Num + Class + N & un métier/artisanat (classif. \zh{门}) \\ \hline
\zh{对……来说} (duì … lái shuō) & Préposition & pour …, du point de vue de … \\ \hline
\end{tabularx}
\end{center}
\end{propriete}

\courssub{Clés (radicaux) et ordre des traits : décoder les caractères}

\begin{methode}[Stratégie « Clé + Phonétique + Contexte » pour 3 caractères du texte]
\begin{enumerate}
\item \textbf{\zh{传} (chuán) — Transmettre} : Clé \zh{亻} (personne debout) → action humaine. Phonétique \zh{专} (zhuān) → son proche. Ordre : 亻(2 traits) → 专 (6 traits : haut → bas, gauche → droite). \emph{Mnémo} : Une personne \zh{亻} transmet son savoir \zh{专} spécialisé.
\item \textbf{\zh{遗} (yí) — Laisser, héritage} : Clé \zh{辶} (marche, mouvement) → action de laisser derrière soi. Phonétique \zh{贵} (guì) → son proche (yí/guì). Ordre : partie gauche \zh{贵} (12 traits : haut → bas) → \zh{辶} (3 traits : point,横折,捺) \textbf{en dernier} (règle : l'enveloppement final se ferme en dernier).
\item \textbf{\zh{家} (jiā) — Famille, maison} : Clé \zh{宀} (toit) → lieu, habitation. Intérieur \zh{豕} (porc/cochon) → dans la Chine ancienne, avoir un porc sous son toit = prospérité. Ordre : \zh{宀} (3 traits : 点, 竖, 横折) → \zh{豕} (7 traits : 内部 de gauche à droite, haut en bas). Règle : \cle{Extérieur avant intérieur, on ferme en dernier}.
\end{enumerate}
\end{methode}

\begin{popfigure}
\begin{tikzpicture}[scale=0.7, every node/.style={font=\small}]
% Stroke order diagram for 家 (Jia) - simplified static representation
\draw[popblue, very thick] (0,0) grid (4,4);
% Roof 宀
\draw[poporange, line width=2pt, ->] (2,3.8) -- (2,3.2) node[midway, right] {1. 点};
\draw[poporange, line width=2pt, ->] (2,3.2) -- (0.5,1.5) node[midway, left] {2. 竖};
\draw[poporange, line width=2pt, ->] (0.5,1.5) -- (3.5,1.5) node[midway, above] {3. 横折};
% Pig 豕 inside
\draw[popgreen, line width=2pt, ->] (1,1.3) -- (3,1.3) node[midway, below] {4. 横};
\draw[popgreen, line width=2pt, ->] (1,1.3) -- (1,0.5) node[midway, left] {5. 竖};
\draw[popgreen, line width=2pt, ->] (1,0.5) -- (3,0.5) node[midway, below] {6. 横};
\draw[popgreen, line width=2pt, ->] (2,1.3) -- (2,0.5) node[midway, right] {7. 竖 (milieu)};
\node[align=center, popdark] at (2, -0.8) {Ordre des traits : \zh{家} (10 traits)\\\textbf{Extérieur (宀) $\rightarrow$ Intérieur (豕)}};
\end{tikzpicture}
\legende{Ordre des traits du caractère \zh{家} (jiā) : le toit (3 traits) puis le porc (7 traits). Règle : on ferme le cadre en dernier (le trait horizontal bas du toit est le 3\^{e} trait, mais le « crochet » final du toit couvre le haut de l'intérieur).}
\end{popfigure}

\courssub{Compréhension du texte : grammaire, structures et questions}

\begin{propriete}[Structure \zh{把} (bǎ) : action sur objet déterminé — Transmettre / Laisser]
\emph{Schéma :} \textbf{Sujet + \zh{把} + Objet + Verbe + Complément de résultat / direction / bénéficiaire}. \\
L'objet est \textbf{connu, précis, déterminé} (souvent avec un classificateur). Le verbe \textbf{ne peut pas être seul} ; il doit être suivi d'une particule (\zh{了}, \zh{给}, \zh{完}, \zh{好}...) ou d'un complément.
\begin{center}
\begin{tabularx}{\linewidth}{|X|X|X|}
\hline
\rowcolor{poppurpleL} \textbf{Phrase du texte} & \textbf{Analyse} & \textbf{Traduction} \\ \hline
\zh{他把…地传给我爸爸} & Suj: \zh{他} | Obj: \zh{一块地} | V: \zh{传给} | Bénéf: \zh{我爸爸} & Il a transmis le terrain à mon père. \\ \hline
\zh{祖先留下的传统} & \zh{留下} sans \zh{把} (sujet inanimé/abstrait, focus sur le résultat) & La tradition laissée par les ancêtres. \\ \hline
\zh{把家族的历史…传给} & Objets multiples coordonnés par \zh{、} : histoire, artisanat, terrain. & Transmettre l'histoire, l'artisanat, le terrain. \\ \hline
\end{tabularx}
\end{center}
\emph{Attention :} On ne dit pas \zh{*他传给了我爸爸一块地} (ordre français SVOI). La structure \zh{把} est \textbf{obligatoire} si l'objet est déterminé et que l'action le « manipule » (transmettre, donner, vendre, jeter).
\end{propriete}

\begin{propriete}[Différence cruciale : \zh{传给} vs \zh{继承} vs \zh{留下}]
\begin{center}
\begin{tabularx}{\linewidth}{|X|X|X|X|}
\hline
\rowcolor{poporangeL} \textbf{Verbe} & \textbf{Agent (Sujet)} & \textbf{Moment} & \textbf{Exemple} \\ \hline
\zh{传给} (chuán gěi) & Donateur (vivant) & De son vivant & \zh{爷爷把手艺传给孙子。} \\ \hline
\zh{留下} (liú xià) & Donateur / Ancêtres & De son vivant / Posthume & \zh{祖先留下了宝贵的财富。} \\ \hline
\zh{继承} (jìchéng) & Héritier / Receveur & \textbf{Après décès} (généralement) & \zh{儿子继承了父亲的土地。} \\ \hline
\end{tabularx}
\end{center}
\cle{Indice contextuel} : \zh{去世后} (après le décès) $\rightarrow$ \zh{继承} ; \zh{决定把…传给} (décide de transmettre) $\rightarrow$ \zh{传给}.
\end{propriete}

\begin{propriete}[Classificateurs obligatoires pour les biens hérités]
En chinois, \textbf{Nombre + Classificateur + Nom} est obligatoire. Pas de « nu » (nom nu) après un nombre/démonstratif.
\begin{center}
\begin{tabularx}{\linewidth}{|X|X|X|}
\hline
\rowcolor{poppinkL} \textbf{Bien} & \textbf{Classificateur} & \textbf{Exemple correct} \\ \hline
Terrain, champ, parcelle & \zh{块} (kuài) & \zh{一块地、这块地、那块地} \\ \hline
Maison, bâtiment, montagne & \zh{座} (zuò) & \zh{一座房子、两座大楼} \\ \hline
Histoire, affaire, chose abstraite & \zh{个} (gè) / \zh{件} (jiàn) & \zh{一个故事、一件事} \\ \hline
Métier, art, technique, langue & \zh{门} (mén) & \zh{一门手艺、两门外语} \\ \hline
Héritage, legs (terme juridique/abstrait) & \zh{份} (fèn) / \zh{项} (xiàng) & \zh{一份遗产、三项专利} \\ \hline
\end{tabularx}
\end{center}
\end{propriete}

\begin{exoresolu}[Compréhension détaillée du texte support — Type Bac]
Énoncé — Relis le texte « La promesse au grand-père ». Réponds en français (sauf Q4 en chinois).
\begin{enumerate}
\item Quel âge a le grand-père et où habite-t-il ? Cite la phrase chinoise.
\item Quels sont les trois éléments que le grand-père décide de transmettre ? (Vocabulaire précis).
\item Explique la phrase : « 这块地不只是财产，更是我们家族的根。 » (Structure \zh{不只是…更是…}).
\item Réponds en chinois : \zh{作者答应爷爷做什么？}
\end{enumerate}
\tcblower
Solution détaillée :
\begin{enumerate}
\item \textbf{Âge et lieu} : 75 ans, village en périphérie de Yaoundé. \\
Citation : \zh{他今年七十五岁了，身体还很硬朗。} \emph{Tā jīnnián qīshíwǔ suì le, shēntǐ hái hěn yìnglǎng.} + \zh{我的爷爷住在雅温得郊区的一个小村庄里。}
\item \textbf{Trois éléments} :
\begin{enumerate}
\item L'histoire de la famille (\zh{家族的历史 jiāzú de lìshǐ}).
\item Le savoir-faire de la sculpture sur bois (\zh{做木雕的手艺 zuò mùdiāo de shǒuyì}).
\item Un terrain près de Douala (\zh{在杜阿拉附近的一块地 zài Dùālā fùjìn de yí kuài dì}).
\end{enumerate}
\item \textbf{Analyse « Racine »} : Structure \zh{不只是 A，更是 B} = « Ce n'est pas seulement A, c'est surtout B ».
\begin{itemize}
\item A = \zh{财产} (biens matériels, valeur marchande).
\item B = \zh{家族的根} (métaphore : origine, identité, mémoire collective, ancrage).
\end{itemize}
Le grand-père oppose valeur \textbf{économique} (jetable, \zh{钱花完了还可以赚}) et valeur \textbf{identitaire} (irremplaçable, \zh{不能丢}).
\item \textbf{Réponse en chinois} : \zh{作者答应爷爷，一定要好好继承这份宝贵的遗产，传给下一代。} \emph{Zuòzhě dāying yéye, yídìng yào hǎohǎo jìchéng zhè fèn bǎoguì de yíchǎn, chuán gěi xià yī dài.}
\end{enumerate}
\end{exoresolu}

\begin{exoresolu}[Grammaire : Transformation de phrases (Structure \zh{把} et Classificateurs)]
Énoncé — Transforme les phrases en utilisant la structure \zh{把} + Objet (avec classificateur correct) + Verbe + Complément. Traduis.
1. 父亲把那座老房子留给了我。 (Focus sur l'action du père sur la maison)
2. 爷爷传给了孙子一门手艺。 (Corrige l'ordre des mots : objet déterminé $\rightarrow$ \zh{把})
3. 我们要继承祖先的传统。 (Pas de \zh{把} ici ! Pourquoi ?)
\tcblower
Solution détaillée :
\begin{enumerate}
\item \textbf{Phrase source déjà correcte} : \zh{父亲把那座老房子留给了我。} \emph{Fùqin bǎ nà zuò lǎo fángzi liú gěi le wǒ.} — « Mon père m'a laissé cette vieille maison. » Classificateur \zh{座} pour maison.
\item \textbf{Correction ordre} : \zh{爷爷把一门手艺传给了孙子。} \emph{Yéye bǎ yí mén shǒuyì chuán gěi le sūnzi.} — « Grand-père a transmis un métier à son petit-fils. » Classificateur \zh{门} pour \zh{手艺}. L'objet \zh{一门手艺} est déterminé $\rightarrow$ structure \zh{把} obligatoire.
\item \textbf{Pas de \zh{把}} : \zh{我们要继承祖先的传统。} \emph{Wǒmen yào jìchéng zǔxiān de chuántǒng.} — « Nous devons hériter de la tradition des ancêtres. » \textbf{Raison} : \zh{继承} (hériter) n'est pas une action de « manipulation/transfert » opérée par le sujet \emph{sur} l'objet de son vivant ; c'est une réception légale/morale. Le sujet \emph{reçoit}, il ne \emph{dispose pas} de l'objet via \zh{把}. De plus, l'objet \zh{祖先的传统} est générique/abstrait, pas un objet concret déterminé qu'on « prend en main ».
\end{enumerate}
\end{exoresolu}

\courssub{Production guidée et réinvestissement}

\begin{experience}[Atelier oral : « L'objet de famille » — 15 min]
\textbf{Consigne} : En binôme. L'élève A est le grand-parent, l'élève B le petit-enfant.
\begin{enumerate}
\item A choisit un objet imaginaire (terrain, maison, recette, bijou, livre de comptes, outil) et son classificateur.
\item A dit : \zh{« 我想把[Classif+Objet]传给你，因为… »} (Je veux te transmettre… car…).
\item B répond : \zh{« 我一定好好继承[Objet]，并传给下一代。 »} (Je l'hériterai bien et le transmettrai).
\item Inversion des rôles. Le prof circule, note les erreurs de tons (\zh{chuán} vs \zh{zhuàn}), classificateurs (\zh{个} au lieu de \zh{块/座/门}), et place de \zh{把}.
\end{enumerate}
\textbf{Évaluation} : Fluidité, tons corrects, classificateur juste, structure \zh{把} / \zh{继承} appropriée.
\end{experience}

\begin{exercice}[Entraînement écrit : Vocabulaire, Grammaire, Traduction]
\textbf{Exercice 1 — Complétion (Classificateurs et Verbes de transmission)} \\
Complète avec le bon classificateur (\zh{块、座、门、份、个}) et le bon verbe (\zh{传给、继承、留下、答应}) :
\begin{enumerate}
\item 爷爷想把做竹编的\_\_\_\_手艺\_\_\_\_孙子。
\item 父亲去世前，把那\_\_\_\_房子\_\_\_\_了我。
\item 这是祖先\_\_\_\_的\_\_\_\_宝贵遗产，我们要\_\_\_\_。
\item 我\_\_\_\_爷爷，一定不丢家族的传统。
\end{enumerate}

\textbf{Exercice 2 — Traduction Thème (Français $\rightarrow$ Chinois)} \\
Traduis en caractères + pinyin :
\begin{enumerate}
\item Mon arrière-grand-père a laissé un terrain à mon père. (Utilise \zh{留下} + \zh{给} ou \zh{传给}).
\item Cette tradition n'est pas seulement une coutume, c'est l'âme de notre famille. (\zh{不只是…更是…} ; \zh{灵魂 línghún}).
\item Nous promettons d'hériter de cet artisanat et de le transmettre aux enfants. (\zh{答应} + \zh{继承} + \zh{传给}).
\end{enumerate}

\textbf{Exercice 3 — Expression écrite (50-80 caractères)} \\
Rédige un court message WeChat à ton cousin(ne) pour lui raconter la décision de ton grand-père de transmettre la maison familiale (\zh{老宅 lǎozhái}) et la recette du \zh{ndolé}. Utilise : \zh{把、传给、继承、宝贵、答应、根}.
\end{exercice}

\begin{corrige}[Exercice 1]
\begin{enumerate}
\item \zh{一门 / 传给} — \zh{爷爷想把做竹编的\textbf{一门}手艺\textbf{传给}孙子。} (Classif. \zh{门} pour métier ; verbe \zh{传给} : action du donneur vivant).
\item \zh{座 / 传给} (ou \zh{留给}) — \zh{父亲去世前，把那\textbf{座}房子\textbf{传给/留给}了我。} (Classif. \zh{座} pour maison ; \zh{传给} insiste sur l'acte de transmission, \zh{留给} sur le legs testamentaire).
\item \zh{留下 / 一份 / 继承} — \zh{这是祖先\textbf{留下}的\textbf{一份}宝贵遗产，我们要\textbf{继承}。} (Verbe \zh{留下} pour ancêtres ; classif. \zh{份} pour héritage abstrait ; verbe \zh{继承} pour nous, receveurs).
\item \zh{答应} — \zh{我\textbf{答应}爷爷，一定不丢家族的传统。} (Structure \zh{答应} qn + phrase engagement).
\end{enumerate}
\end{corrige}

\begin{corrige}[Exercice 2]
\begin{enumerate}
\item \zh{我的曾祖父把一块地留给了/传给了我父亲。} \emph{Wǒ de zēngzǔfù bǎ yí kuài dì liú gěi le / chuán gěi le wǒ fùqin.}
\item \zh{这传统不只是习俗，更是我们家族的灵魂。} \emph{Zhè chuántǒng bù zhǐ shì xísú, gèng shì wǒmen jiāzú de línghún.}
\item \zh{我们答应要好好继承这门手艺，并传给孩子们。} \emph{Wǒmen dāying yào hǎohǎo jìchéng zhè mén shǒuyì, bìng chuán gěi háizimen.}
\end{enumerate}
\end{corrige}

\begin{corrige}[Exercice 3 — Proposition de rédaction]
\zh{表弟，爷爷决定把老宅和做ndolé的手艺传给爸爸。他说老宅是我们家族的根，手艺是宝贵的遗产。我答应爷爷，一定继承好，将来传给下一代。你也要记得哦！} \\
\emph{Biaodi, yéye juédìng bǎ lǎozhái hé zuò ndolé de shǒuyì chuán gěi bàba. Tā shuō lǎozhái shì wǒmen jiāzú de gēn, shǒuyì shì bǎoguì de yíchǎn. Wǒ dāying yéye, yídìng jìchéng hǎo, jiānglái chuán gěi xià yī dài. Nǐ yě yào jìde o!} \\
« Cousin, grand-père a décidé de transmettre la maison ancestrale et le savoir-faire du ndolé à papa. Il dit que la maison est la racine de notre famille, l'artisanat un héritage précieux. J'ai promis à grand-père de bien l'hériter et de le transmettre plus tard. Toi aussi, souviens-toi ! »
\end{corrige}

\begin{attention}[Les 5 erreurs qui coûtent des points au Bac (Rappel synthétique)]
\begin{enumerate}
\item \textbf{Confondre \zh{传给} (donneur, vivant) et \zh{继承} (receveur, post-mortem).} \cle{Test} : Qui est le sujet ? \zh{爷爷传给} / \zh{儿子继承}.
\item \textbf{Oublier le classificateur.} \zh{*一地、一房子、一手艺} $\rightarrow$ \textbf{Faux}. Correct : \zh{一块地、一座房子、一门手艺}.
\item \textbf{Mauvais ordre des mots (calque français) :} \zh{*他传给我一块地} $\rightarrow$ \textbf{Faux}. Correct : \zh{他把一块地传给我。} (Objet déterminé + \zh{把}).
\item \textbf{Tons approximatifs sur vocabulaire-clé :} \zh{yí \textbf{ch}án} (non, \zh{chuán} 2), \zh{bǎo \textbf{guì}} (3-4), \zh{jiā \textbf{zú}} (1-2), \zh{mù \textbf{diāo}} (4-1). Dictee pinyin = points faciles si révisés.
\item \textbf{Caractères proches (pièges visuels) :} \zh{传} (transmettre, 亻) $\neq$ \zh{转} (tourner, 车) ; \zh{遗} (laisser, 辶) $\neq$ \zh{贵} (précieux, 贝) ; \zh{宝} (trésor, 宀) $\neq$ \zh{实} (réel, 宀+贯). Vérifie la \cle{clé} avant d'écrire.
\end{enumerate}
\end{attention}

\newpage

\coursec{Exercices}

\courssub{Je vérifie mes connaissances}

\begin{exercice}[QCM : le vocabulaire de l'héritage]
Pour chaque question, entoure la bonne réponse.

\begin{enumerate}
\item \zh{遗产} (yíchǎn) signifie :
\begin{enumerate}
\item la maison \item l'héritage \item le mariage
\end{enumerate}
\item \zh{继承} (jìchéng) signifie :
\begin{enumerate}
\item hériter \item acheter \item vendre
\end{enumerate}
\item \zh{遗嘱} (yízhǔ) signifie :
\begin{enumerate}
\item le testament \item la fête \item le champ
\end{enumerate}
\item \zh{宝贵} (bǎoguì) signifie :
\begin{enumerate}
\item pauvre \item précieux \item ancien
\end{enumerate}
\item \zh{财产} (cáichǎn) signifie :
\begin{enumerate}
\item les biens, la fortune \item les voisins \item les enfants
\end{enumerate}
\end{enumerate}
\end{exercice}

\begin{exercice}[Vrai ou faux ? Justifie ta réponse]
D'après le texte support de la leçon, indique si chaque phrase est vraie (V) ou fausse (F) et justifie en français.

\begin{enumerate}
\item Le grand-père a laissé une maison et des champs à sa famille.
\item Pour le narrateur, l'argent est plus important que la culture.
\item Le grand-père a appris au narrateur à parler chinois.
\item Le narrateur veut transmettre cet héritage à ses enfants.
\end{enumerate}
\end{exercice}

\begin{exercice}[Complète avec le bon mot]
Complète chaque phrase avec le mot qui convient : \zh{遗产}, \zh{继承}, \zh{遗嘱}, \zh{财产}, \zh{宝贵}.

\begin{enumerate}
\item 爷爷去世以前立了\_\_\_\_\_\_。(Yéye qùshì yǐqián lì le \_\_\_\_\_\_.)
\item 他留给我们的\_\_\_\_\_\_是一栋房子。(Tā liú gěi wǒmen de \_\_\_\_\_\_ shì yí dòng fángzi.)
\item 我要\_\_\_\_\_\_爸爸的事业。(Wǒ yào \_\_\_\_\_\_ bàba de shìyè.)
\item 中国的文化非常\_\_\_\_\_\_。(Zhōngguó de wénhuà fēicháng \_\_\_\_\_\_.)
\item 他们分好了家里的\_\_\_\_\_\_。(Tāmen fēn hǎo le jiā li de \_\_\_\_\_\_.)
\end{enumerate}
\end{exercice}

\begin{exercice}[Associe le pinyin aux caractères]
Relie chaque caractère ou mot à son pinyin.

\begin{center}
\begin{tabularx}{\linewidth}{|X|X|}
\hline
\rowcolor{popblueL} Caractères & Pinyin \\
\hline
1. 遗产 & a. yízhǔ \\
\hline
2. 遗嘱 & b. chuán \\
\hline
3. 传 & c. jìchéng \\
\hline
4. 继承 & d. yíchǎn \\
\hline
5. 宝贵 & e. bǎoguì \\
\hline
\end{tabularx}
\end{center}
\end{exercice}

\courssub{Je m'entraîne}

\begin{exercice}[Traduis en français]
Traduis les phrases suivantes en français.

\begin{enumerate}
\item \zh{爷爷给我们留下了一栋房子和几块地。}\\ (Yéye gěi wǒmen liúxià le yí dòng fángzi hé jǐ kuài dì.)
\item \zh{我觉得文化比钱更宝贵。}\\ (Wǒ juéde wénhuà bǐ qián gèng bǎoguì.)
\item \zh{父亲去世以前立了遗嘱。}\\ (Fùqin qùshì yǐqián lì le yízhǔ.)
\item \zh{我们要把传统传下去。}\\ (Wǒmen yào bǎ chuántǒng chuán xiàqu.)
\end{enumerate}
\end{exercice}

\begin{exercice}[Remets les mots dans l'ordre]
Remets les mots dans le bon ordre pour former une phrase correcte, puis traduis-la.

\begin{enumerate}
\item \zh{遗产 / 了 / 爷爷 / 留下 / 很多}
\item \zh{宝贵 / 是 / 文化 / 财富 / 的 / 最}
\item \zh{继承 / 想 / 我 / 的 / 事业 / 父亲}
\item \zh{把 / 房子 / 他们 / 分 / 好 / 了}
\end{enumerate}
\end{exercice}

\begin{exercice}[Texte à trous]
Complète le texte suivant avec les mots de la banque : \zh{遗产}, \zh{继承}, \zh{宝贵}, \zh{传}, \zh{财产}.

\begin{textebox}[Un héritage précieux]
我的爷爷是农民。他去世以后，我们\_\_\_\_\_\_了他的房子和地。这些\_\_\_\_\_\_不多，但是爷爷还给我们留下了更\_\_\_\_\_\_的东西：他教我们的歌和故事。爸爸常说：« 钱会用完，文化不会。» 所以我要把这些故事\_\_\_\_\_\_下去，让我的孩子也知道。\\
(Wǒ de yéye shì nóngmín. Tā qùshì yǐhòu, wǒmen \_\_\_\_\_\_ le tā de fángzi hé dì. Zhèxiē \_\_\_\_\_\_ bù duō, dànshì yéye hái gěi wǒmen liúxià le gèng \_\_\_\_\_\_ de dōngxi : tā jiāo wǒmen de gē hé gùshi. Bàba cháng shuō : « Qián huì yòng wán, wénhuà bú huì. » Suǒyǐ wǒ yào bǎ zhèxiē gùshi \_\_\_\_\_\_ xiàqu, ràng wǒ de háizi yě zhīdao.)\\
\emph{Mon grand-père était paysan. Après sa mort, nous avons hérité de sa maison et de ses champs. Ces biens ne sont pas nombreux, mais grand-père nous a aussi laissé quelque chose de plus précieux : les chansons et les histoires qu'il nous a apprises. Papa dit souvent : « L'argent s'épuise, pas la culture. » Je veux donc transmettre ces histoires, pour que mes enfants les connaissent aussi.}
\end{textebox}
\end{exercice}

\begin{exercice}[Associe les collocations à leur traduction]
Relie chaque collocation chinoise à sa traduction française.

\begin{center}
\begin{tabularx}{\linewidth}{|X|X|}
\hline
\rowcolor{popblueL} Collocation & Traduction \\
\hline
1. 立遗嘱 (lì yízhǔ) & a. transmettre (aux générations suivantes) \\
\hline
2. 分财产 (fēn cáichǎn) & b. laisser un héritage \\
\hline
3. 传下去 (chuán xiàqu) & c. faire son testament \\
\hline
4. 留下遗产 (liúxià yíchǎn) & d. partager les biens \\
\hline
5. 继承事业 (jìchéng shìyè) & e. reprendre l'affaire (le métier) de quelqu'un \\
\hline
\end{tabularx}
\end{center}
\end{exercice}

\courssub{Je me prépare au Bac}

\begin{exercice}[Type Bac — Compréhension écrite]
Lis le texte suivant, puis réponds aux questions \textbf{en chinois, par des phrases complètes}.

\begin{textebox}[L'héritage de grand-mère]
我的奶奶住在雅温得。她今年八十岁了。奶奶没有很多钱，也没有大房子，但是她给了我们最宝贵的遗产。她年轻的时候是老师，教过很多学生。她常常对我们说：« 知识是谁也拿不走的财产。» 奶奶会讲我们民族的故事，会做传统的菜，还会唱老歌。每个星期天，我们都去她家，听她讲故事。去年，奶奶立了遗嘱：她把她的小房子留给了家里最小的孩子，把她的书留给了我。我觉得这些书比房子更宝贵，因为书里有奶奶的一生。我要好好保存这些书，以后把奶奶的故事传下去，让我的孩子也认识他们的曾祖母。\\
(Wǒ de nǎinai zhù zài Yǎwēndé. Tā jīnnián bāshí suì le. Nǎinai méiyǒu hěn duō qián, yě méiyǒu dà fángzi, dànshì tā gěi le wǒmen zuì bǎoguì de yíchǎn. Tā niánqīng de shíhou shì lǎoshī, jiāo guo hěn duō xuésheng. Tā chángcháng duì wǒmen shuō : « Zhīshi shì shéi yě ná bù zǒu de cáichǎn. » Nǎinai huì jiǎng wǒmen mínzú de gùshi, huì zuò chuántǒng de cài, hái huì chàng lǎo gē. Měi ge xīngqītiān, wǒmen dōu qù tā jiā, tīng tā jiǎng gùshi. Qùnián, nǎinai lì le yízhǔ : tā bǎ tā de xiǎo fángzi liú gěi le jiā li zuì xiǎo de háizi, bǎ tā de shū liú gěi le wǒ. Wǒ juéde zhèxiē shū bǐ fángzi gèng bǎoguì, yīnwèi shū li yǒu nǎinai de yìshēng. Wǒ yào hǎohao bǎocún zhèxiē shū, yǐhòu bǎ nǎinai de gùshi chuán xiàqu, ràng wǒ de háizi yě rènshi tāmen de zēngzǔmǔ.)\\
\emph{Ma grand-mère habite à Yaoundé. Elle a quatre-vingts ans cette année. Grand-mère n'a pas beaucoup d'argent ni de grande maison, mais elle nous a donné l'héritage le plus précieux. Dans sa jeunesse, elle était institutrice et a enseigné à beaucoup d'élèves. Elle nous disait souvent : « Le savoir est un bien que personne ne peut prendre. » Grand-mère sait raconter les histoires de notre ethnie, préparer les plats traditionnels et chanter de vieilles chansons. Chaque dimanche, nous allons chez elle l'écouter raconter des histoires. L'année dernière, grand-mère a fait son testament : elle a laissé sa petite maison au plus jeune enfant de la famille, et ses livres à moi. Je pense que ces livres sont plus précieux que la maison, car ils contiennent la vie de grand-mère. Je veux bien conserver ces livres, puis transmettre les histoires de grand-mère, pour que mes enfants connaissent aussi leur arrière-grand-mère.}
\end{textebox}

\begin{enumerate}
\item 奶奶年轻的时候做什么工作？(Nǎinai niánqīng de shíhou zuò shénme gōngzuò ?)
\item 奶奶常常对孩子们说什么？(Nǎinai chángcháng duì háizimen shuō shénme ?)
\item 奶奶把房子留给了谁？(Nǎinai bǎ fángzi liú gěi le shéi ?)
\item 为什么作者觉得书比房子更宝贵？(Wèishénme zuòzhě juéde shū bǐ fángzi gèng bǎoguì ?)
\item 作者以后想做什么？(Zuòzhě yǐhòu xiǎng zuò shénme ?)
\end{enumerate}

\textbf{Questions de langue :}
\begin{enumerate}
\item Trouve dans le texte une phrase avec la structure comparatives \zh{比} et recopie-la.
\item Explique en français le sens de la phrase : \zh{知识是谁也拿不走的财产。}
\end{enumerate}
\end{exercice}

\begin{exercice}[Type Bac — Thème et version]
\textbf{A. Thème.} Traduis en chinois (caractères simplifiés) :

\begin{enumerate}
\item Mon grand-père a laissé une maison et des champs ; je pense que notre culture est plus précieuse que l'argent.
\item Avant de mourir, mon père a fait son testament.
\item Nous devons transmettre nos traditions à nos enfants.
\end{enumerate}

\textbf{B. Version.} Traduis en français :

\begin{enumerate}
\item \zh{奶奶给我们留下了最宝贵的遗产。}\\ (Nǎinai gěi wǒmen liúxià le zuì bǎoguì de yíchǎn.)
\item \zh{他们分好了父亲留下的财产。}\\ (Tāmen fēn hǎo le fùqin liúxià de cáichǎn.)
\item \zh{我想继承父亲的事业。}\\ (Wǒ xiǎng jìchéng fùqin de shìyè.)
\end{enumerate}
\end{exercice}

\begin{exercice}[Type Bac — Expression écrite et écriture des caractères]
\textbf{A. Expression écrite.} Rédige en chinois un court texte intitulé \zh{我家的遗产} (Wǒ jiā de yíchǎn, « L'héritage de ma famille »). Tu raconteras ce qu'un membre de ta famille (grand-père, grand-mère, oncle...) a laissé à ta famille : des biens matériels (maison, champs, argent) ou immatériels (histoires, chansons, traditions, savoir-faire). Tu diras ce que tu veux en faire. \textbf{Consignes :} écris entre 80 et 120 caractères chinois ; utilise au moins quatre mots du vocabulaire de la leçon (\zh{遗产}, \zh{继承}, \zh{遗嘱}, \zh{财产}, \zh{宝贵}, \zh{传下去}) et une structure comparative avec \zh{比}.

\textbf{B. Écriture des caractères.} Trace les caractères \zh{遗}, \zh{继}, \zh{财}, \zh{嘱} en respectant l'ordre des traits (de haut en bas, de gauche à droite, l'extérieur avant l'intérieur pour les caractères à enclos), puis indique la clé (radical) de chacun d'eux.
\end{exercice}

\newpage

\coursec{Corrigés des exercices}

\begin{corrige}[Exercice 1 — QCM : le vocabulaire de l'héritage]
\begin{enumerate}
\item \textbf{b.} l'héritage. \zh{遗产} (yíchǎn) = héritage, ce qui est laissé par les défunts. La maison se dit \zh{房子} (fángzi), le mariage \zh{结婚} (jiéhūn).
\item \textbf{a.} hériter. \zh{继承} (jìchéng) = hériter, succéder. Acheter = \zh{买} (mǎi), vendre = \zh{卖} (mài).
\item \textbf{a.} le testament. \zh{遗嘱} (yízhǔ) = testament (document laissé avant la mort).
\item \textbf{b.} précieux. \zh{宝贵} (bǎoguì) = précieux, de grande valeur. Pauvre = \zh{穷} (qióng), ancien = \zh{老} / \zh{旧} (lǎo / jiù).
\item \textbf{a.} les biens, la fortune. \zh{财产} (cáichǎn) = biens, patrimoine. Les voisins = \zh{邻居} (línjū), les enfants = \zh{孩子} (háizi).
\end{enumerate}
\textbf{Astuce mémorisation :} les mots de ce thème partagent souvent le caractère \zh{遗} (yí, « laisser après soi ») : \zh{遗产} (héritage), \zh{遗嘱} (testament), \zh{留下} (laisser).
\end{corrige}

\begin{corrige}[Exercice 2 — Vrai ou faux ? Justifie ta réponse]
\begin{enumerate}
\item \textbf{Vrai.} Le texte dit que le grand-père a laissé à sa famille une maison et des champs (\zh{爷爷给我们留下了一栋房子和几块地}).
\item \textbf{Faux.} C'est le contraire : pour le narrateur, la culture est plus précieuse que l'argent (\zh{我觉得文化比钱更宝贵} — « je pense que la culture est plus précieuse que l'argent »).
\item \textbf{Faux.} Le grand-père a appris au narrateur des chansons et des histoires (\zh{歌和故事}), pas à parler chinois.
\item \textbf{Vrai.} Le narrateur déclare vouloir transmettre cet héritage à ses enfants (\zh{我要把这些故事传下去，让我的孩子也知道}).
\end{enumerate}
\textbf{Méthode :} pour justifier, repère dans le texte la phrase exacte qui confirme ou contredit l'affirmation, et cite-la ou reformule-la en français.
\end{corrige}

\begin{corrige}[Exercice 3 — Complète avec le bon mot]
\begin{enumerate}
\item 爷爷去世以前立了\textbf{遗嘱}。(Yéye qùshì yǐqián lì le \textbf{yízhǔ}.) — « Avant de mourir, grand-père a fait son testament. » Collocation obligatoire : \zh{立遗嘱} (lì yízhǔ).
\item 他留给我们的\textbf{遗产}是一栋房子。(Tā liú gěi wǒmen de \textbf{yíchǎn} shì yí dòng fángzi.) — « L'héritage qu'il nous a laissé est une maison. »
\item 我要\textbf{继承}爸爸的事业。(Wǒ yào \textbf{jìchéng} bàba de shìyè.) — « Je veux reprendre l'affaire de mon père. » Ici il faut un verbe : \zh{继承}.
\item 中国的文化非常\textbf{宝贵}。(Zhōngguó de wénhuà fēicháng \textbf{bǎoguì}.) — « La culture chinoise est très précieuse. » Ici il faut un adjectif : \zh{宝贵}.
\item 他们分好了家里的\textbf{财产}。(Tāmen fēn hǎo le jiā li de \textbf{cáichǎn}.) — « Ils ont bien partagé les biens de la famille. » Collocation : \zh{分财产} (fēn cáichǎn).
\end{enumerate}
\textbf{Point de vigilance :} avant de choisir, identifie la nature grammaticale attendue (nom, verbe ou adjectif) d'après la structure de la phrase.
\end{corrige}

\begin{corrige}[Exercice 4 — Associe le pinyin aux caractères]
\begin{center}
\begin{tabularx}{\linewidth}{|X|X|X|}
\hline
\rowcolor{popblueL} Caractères & Pinyin & Traduction \\
\hline
1. 遗产 & \textbf{d.} yíchǎn & l'héritage \\
\hline
2. 遗嘱 & \textbf{a.} yízhǔ & le testament \\
\hline
3. 传 & \textbf{b.} chuán & transmettre \\
\hline
4. 继承 & \textbf{c.} jìchéng & hériter \\
\hline
5. 宝贵 & \textbf{e.} bǎoguì & précieux \\
\hline
\end{tabularx}
\end{center}
\textbf{Attention aux tons :} \zh{遗} (yí, 2\up{e} ton) dans \zh{遗产} et \zh{遗嘱} ; \zh{传} se lit \emph{chuán} (2\up{e} ton) quand il signifie « transmettre ». Ne confonds pas \zh{继} (jì, 4\up{e} ton) avec \zh{记} (jì, « se souvenir ») : même pinyin, sens différents.
\end{corrige}

\begin{corrige}[Exercice 5 — Traduis en français]
\begin{enumerate}
\item \zh{爷爷给我们留下了一栋房子和几块地。} — « Grand-père nous a laissé une maison et quelques champs (parcelles). »\\ Structure : \zh{给} + personne + \zh{留下} + \zh{了} + objet = « laisser qqch à qqn ». \zh{栋} (dòng) est le classificateur des bâtiments, \zh{块} (kuài) celui des parcelles de terre.
\item \zh{我觉得文化比钱更宝贵。} — « Je pense que la culture est plus précieuse que l'argent. »\\ Structure comparative : A + \zh{比} + B + \zh{更} + adjectif.
\item \zh{父亲去世以前立了遗嘱。} — « Avant de mourir, mon père a fait son testament. »\\ \zh{去世以前} = « avant de mourir » (le complément de temps se place avant le verbe principal).
\item \zh{我们要把传统传下去。} — « Nous devons transmettre les traditions (aux générations suivantes). »\\ Construction en \zh{把} : Sujet + \zh{把} + objet + verbe + complément (\zh{下去} = « vers le bas / dans la continuité »).
\end{enumerate}
\end{corrige}

\begin{corrige}[Exercice 6 — Remets les mots dans l'ordre]
\begin{enumerate}
\item \zh{爷爷留下了很多遗产。} (Yéye liúxià le hěn duō yíchǎn.) — « Grand-père a laissé beaucoup d'héritage. » Ordre : Sujet + verbe + \zh{了} + complément de quantité + objet.
\item \zh{文化是最宝贵的财富。} (Wénhuà shì zuì bǎoguì de cáifù.) — « La culture est le bien le plus précieux. » Superlatif : \zh{最} + adjectif + \zh{的} + nom.
\item \zh{我想继承父亲的事业。} (Wǒ xiǎng jìchéng fùqin de shìyè.) — « Je veux reprendre l'affaire de mon père. » Le verbe modal \zh{想} précède le verbe principal \zh{继承}.
\item \zh{他们把房子分好了。} (Tāmen bǎ fángzi fēn hǎo le.) — « Ils ont bien partagé la maison. » Construction en \zh{把} : l'objet \zh{房子} avance avant le verbe ; \zh{好} est le complément de résultat (« de façon achevée »).
\end{enumerate}
\textbf{Point de vigilance :} en chinois, l'ordre est Sujet — (complément de temps/lieu) — Verbe — Objet. Ne place jamais le complément de temps après le verbe comme en français.
\end{corrige}

\begin{corrige}[Exercice 7 — Texte à trous]
\begin{enumerate}
\item 我们\textbf{继承}了他的房子和地。(jìchéng) — verbe « hériter de », suivi de \zh{了} (action achevée).
\item 这些\textbf{遗产}不多。(yíchǎn) — nom, « cet héritage n'est pas important ». (\zh{财产} serait possible grammaticalement, mais \zh{遗产} convient mieux au sens : ce qui est laissé par le défunt ; on réserve \zh{财产} pour le troisième trou non — voir ci-dessous : la banque impose ici \zh{遗产}.)
\item 更\textbf{宝贵}的东西 (bǎoguì) — adjectif après \zh{更} : « quelque chose de plus précieux ».
\item 把这些故事\textbf{传}下去 (chuán) — verbe dans la structure \zh{把} + objet + \zh{传下去} : « transmettre ».
\end{enumerate}
Texte complet : \zh{我的爷爷是农民。他去世以后，我们继承了他的房子和地。这些遗产不多，但是爷爷还给我们留下了更宝贵的东西：他教我们的歌和故事。爸爸常说：« 钱会用完，文化不会。» 所以我要把这些故事传下去，让我的孩子也知道。}\\
\textbf{Note :} le mot \zh{财产} de la banque est un distracteur : il désigne les biens matériels en général, alors que le texte insiste sur ce qui est « laissé » par le grand-père (\zh{遗产}).
\end{corrige}

\begin{corrige}[Exercice 8 — Associe les collocations à leur traduction]
\begin{center}
\begin{tabularx}{\linewidth}{|X|X|}
\hline
\rowcolor{popblueL} Collocation & Traduction \\
\hline
1. 立遗嘱 (lì yízhǔ) & \textbf{c.} faire son testament \\
\hline
2. 分财产 (fēn cáichǎn) & \textbf{d.} partager les biens \\
\hline
3. 传下去 (chuán xiàqu) & \textbf{a.} transmettre (aux générations suivantes) \\
\hline
4. 留下遗产 (liúxià yíchǎn) & \textbf{b.} laisser un héritage \\
\hline
5. 继承事业 (jìchéng shìyè) & \textbf{e.} reprendre l'affaire (le métier) de quelqu'un \\
\hline
\end{tabularx}
\end{center}
\textbf{À retenir :} en chinois, un verbe se combine avec un objet précis ; apprends toujours la collocation entière (\zh{立遗嘱}, jamais \zh{做遗嘱}).
\end{corrige}

\begin{corrige}[Exercice 9 — Type Bac — Compréhension écrite]
\textbf{Barème indicatif (10 pts) :} 1 pt par réponse de compréhension correcte et complète (5 pts) ; 2 pts pour la qualité linguistique des réponses en chinois ; 1,5 pt par question de langue (3 pts).

\textbf{Questions de compréhension — réponses attendues :}
\begin{enumerate}
\item 奶奶年轻的时候是老师，教过很多学生。(Nǎinai niánqīng de shíhou shì lǎoshī, jiāo guo hěn duō xuésheng.) — « Dans sa jeunesse, grand-mère était institutrice ; elle a enseigné à beaucoup d'élèves. »
\item 她常常对孩子们说：« 知识是谁也拿不走的财产。» — « Elle disait souvent : le savoir est un bien que personne ne peut prendre. »
\item 奶奶把房子留给了家里最小的孩子。(Nǎinai bǎ fángzi liú gěi le jiā li zuì xiǎo de háizi.) — « Elle a laissé la maison au plus jeune enfant de la famille. »
\item 因为书里有奶奶的一生。(Yīnwèi shū li yǒu nǎinai de yìshēng.) — « Parce que les livres contiennent la vie de grand-mère. »
\item 作者想好好保存这些书，以后把奶奶的故事传下去，让自己的孩子也认识他们的曾祖母。 — « L'auteur veut bien conserver les livres et transmettre les histoires de grand-mère pour que ses enfants connaissent leur arrière-grand-mère. »
\end{enumerate}

\textbf{Questions de langue :}
\begin{enumerate}
\item Phrase comparative avec \zh{比} : \zh{我觉得这些书比房子更宝贵。} (Wǒ juéde zhèxiē shū bǐ fángzi gèng bǎoguì.) Structure : A + \zh{比} + B + (\zh{更}) + adjectif.
\item \zh{知识是谁也拿不走的财产。} signifie : « Le savoir est un bien que personne ne peut (nous) prendre / enlever. » Structure : \zh{谁也} + verbe négatif = « personne ne... » ; \zh{拿不走} = « ne pas pouvoir emporter » (complément de possibilité négatif).
\end{enumerate}
\textbf{Points de vigilance :} réponds par des \textbf{phrases complètes} en reprenant les éléments de la question ; ne recopie pas tout le texte ; vérifie les tons dans le pinyin et la place de \zh{把} dans la question 3.
\end{corrige}

\begin{corrige}[Exercice 10 — Type Bac — Thème et version]
\textbf{Barème indicatif (10 pts) :} 2 pts par phrase (1 pt pour l'exactitude du vocabulaire, 1 pt pour la grammaire et l'ordre des mots).

\textbf{A. Thème — propositions de traduction :}
\begin{enumerate}
\item \zh{爷爷留下了一栋房子和几块地，我觉得我们的文化比钱更宝贵。}\\ (Yéye liúxià le yí dòng fángzi hé jǐ kuài dì, wǒ juéde wǒmen de wénhuà bǐ qián gèng bǎoguì.)\\ Points clés : classificateur \zh{栋} pour la maison, \zh{块} pour les champs ; comparatif \zh{比} + \zh{更} + adjectif.
\item \zh{父亲去世以前立了遗嘱。} (Fùqin qùshì yǐqián lì le yízhǔ.)\\ Point clé : « avant de mourir » se traduit par \zh{去世以前}, placé \emph{avant} le verbe.
\item \zh{我们要把传统传给我们的孩子。} ou \zh{我们要把传统传下去，让孩子们知道。}\\ (Wǒmen yào bǎ chuántǒng chuán gěi wǒmen de háizi.)\\ Point clé : construction en \zh{把} avec \zh{传给} + personne.
\end{enumerate}

\textbf{B. Version — traductions :}
\begin{enumerate}
\item \zh{奶奶给我们留下了最宝贵的遗产。} — « Grand-mère nous a laissé l'héritage le plus précieux. »
\item \zh{他们分好了父亲留下的财产。} — « Ils ont bien partagé les biens laissés par leur père. »
\item \zh{我想继承父亲的事业。} — « Je veux reprendre l'affaire (le métier) de mon père. »
\end{enumerate}
\textbf{Points de vigilance :} \zh{留下} = « laisser (en héritage) » ; \zh{分好} = « partager de manière achevée » ; ne traduis pas \zh{事业} par « carrière » mais par « affaire / entreprise / œuvre » selon le contexte.
\end{corrige}

\begin{corrige}[Exercice 11 — Type Bac — Expression écrite et écriture des caractères]
\textbf{Barème indicatif (10 pts) :} respect de la longueur 80–120 caractères (1 pt) ; au moins quatre mots du thème (2 pts) ; une comparaison avec \zh{比} (1 pt) ; correction grammaticale et ordre des mots (3 pts) ; exactitude des caractères (2 pts) ; cohérence du récit (1 pt).

\textbf{A. Proposition de rédaction modèle :}

\begin{textebox}[我家的遗产 — modèle de rédaction]
我的爷爷是农民，他去世以前立了遗嘱，把房子和几块地留给了我们。这些财产不多，但是我觉得爷爷留下的最宝贵的遗产不是钱，而是他教我们的故事和歌。爷爷常说：« 文化比钱更宝贵。» 我同意他的话。所以我要好好学习我们民族的传统，把这些故事传下去。以后我有了孩子，我也要教他们唱爷爷的歌，让他们知道自己的根。这就是我家的遗产。\\
(Wǒ de yéye shì nóngmín, tā qùshì yǐqián lì le yízhǔ, bǎ fángzi hé jǐ kuài dì liú gěi le wǒmen. Zhèxiē cáichǎn bù duō, dànshì wǒ juéde yéye liúxià de zuì bǎoguì de yíchǎn bú shì qián, ér shì tā jiāo wǒmen de gùshi hé gē. Yéye cháng shuō : « Wénhuà bǐ qián gèng bǎoguì. » Wǒ tóngyì tā de huà. Suǒyǐ wǒ yào hǎohao xuéxí wǒmen mínzú de chuántǒng, bǎ zhèxiē gùshi chuán xiàqu. Yǐhòu wǒ yǒu le háizi, wǒ yě yào jiāo tāmen chàng yéye de gē, ràng tāmen zhīdao zìjǐ de gēn. Zhè jiù shì wǒ jiā de yíchǎn.)\\
\emph{Mon grand-père était paysan ; avant de mourir, il a fait son testament et nous a laissé la maison et quelques champs. Ces biens sont modestes, mais je pense que l'héritage le plus précieux qu'il nous a laissé n'est pas l'argent, mais les histoires et les chansons qu'il nous a apprises. Grand-père disait souvent : « La culture est plus précieuse que l'argent. » Je suis d'accord avec lui. Je veux donc bien apprendre les traditions de notre ethnie et transmettre ces histoires. Plus tard, quand j'aurai des enfants, je leur apprendrai aussi les chansons de grand-père, pour qu'ils connaissent leurs racines. Voilà l'héritage de ma famille.}
\end{textebox}

Ce modèle contient environ 150 caractères : à l'examen, raccourcis-le légèrement pour rester entre 80 et 120 caractères. Il utilise \zh{遗产}, \zh{遗嘱}, \zh{财产}, \zh{宝贵}, \zh{传下去} et la comparaison \zh{文化比钱更宝贵}.

\textbf{Points de vigilance :} annonce le thème dès la première phrase ; utilise \zh{把} pour « laisser qqch à qqn » (\zh{把房子留给了我们}) ; termine par une phrase de conclusion ; compte tes caractères.

\textbf{B. Écriture des caractères :}
\begin{itemize}
\item \zh{遗} (yí) : clé \zh{辶} (chuò, « clé du mouvement / marcher »). On écrit d'abord la partie intérieure \zh{贵} (de haut en bas), puis la clé \zh{辶} en dernier (point, puis la courbe en trois temps). Règle : l'intérieur avant la clé de mouvement qui enveloppe.
\item \zh{继} (jì) : clé \zh{纟} (sī, « clé de la soie »), à gauche. On écrit d'abord la clé de soie (trois traits, de haut en bas), puis la partie droite, de gauche à droite et de haut en bas.
\item \zh{财} (cái) : clé \zh{贝} (bèi, « clé du coquillage / de l'argent »), à gauche. On écrit \zh{贝} d'abord (le vertical, le cadre, puis les deux traits intérieurs et le point final), puis \zh{才} à droite (horizontal, puis le crochet vertical, puis le trait tombant).
\item \zh{嘱} (zhǔ) : clé \zh{口} (kǒu, « clé de la bouche »), à gauche. On écrit d'abord la petite bouche \zh{口} à gauche (vertical, horizontal-crochet, horizontal de fermeture), puis la partie droite \zh{属} de haut en bas et de l'extérieur vers l'intérieur.
\end{itemize}
\textbf{Règle générale rappelée :} de haut en bas, de gauche à droite, l'extérieur avant l'intérieur, et on ferme les enclos en dernier.
\end{corrige}

\newpage

\coursec{Fiche bilan}

\begin{popfigure}
\begin{tikzpicture}[mindmap, grow cyclic, every node/.style={concept, text=white, font=\small}, root concept/.append style={concept color=popblue, font=\bfseries}, level 1/.append style={level distance=3.4cm, sibling angle=51}, level 1 concept/.append style={font=\footnotesize}]
\node[root concept, align=center]{L'héritage\\ \zh{遗产} yíchǎn}
  child[concept color=poporange]{ node[align=center]{Famille\\ \zh{父母} fùmǔ\\ \zh{子女} zǐnǚ} }
  child[concept color=popgreen]{ node[align=center]{Biens\\ \zh{房子} fángzi\\ \zh{土地} tǔdì\\ \zh{钱} qián} }
  child[concept color=poppurple]{ node[align=center]{Actes\\ \zh{继承} jìchéng\\ \zh{留给} liúgěi\\ \zh{分} fēn} }
  child[concept color=poppink]{ node[align=center]{Culture\\ \zh{传统} chuántǒng\\ \zh{文化} wénhuà} }
  child[concept color=popgold]{ node[align=center]{Grammaire\\ \zh{把} bǎ\\ \zh{被} bèi\\ \zh{给} gěi} }
  child[concept color=popmuted]{ node[align=center]{Écriture\\ clé 宀 toit\\ clé 子 enfant\\ clé 田 champ} }
  child[concept color=popgreen!70!black]{ node[align=center]{Cameroun / Chine\\ succession\\ coutume et loi} };
\end{tikzpicture}
\legende{Carte mentale de la leçon : le vocabulaire et les notions autour de l'héritage.}
\end{popfigure}

\begin{bilan}
\textbf{1. Lexique clé à connaître par cœur}
\begin{itemize}
  \item \zh{遗产} yíchǎn : héritage, patrimoine ; \zh{继承} jìchéng : hériter (de) ; \zh{留给} liúgěi : laisser à, léguer.
  \item \zh{父母} fùmǔ : les parents ; \zh{子女} zǐnǚ : les enfants ; \zh{孙子} sūnzi : petit-fils.
  \item \zh{房子} fángzi : maison ; \zh{土地} tǔdì : terre, terrain ; \zh{财产} cáichǎn : biens, fortune.
  \item \zh{传统} chuántǒng : tradition ; \zh{文化} wénhuà : culture ; \zh{习惯} xíguàn : coutume, habitude.
  \item \zh{分} fēn : partager, diviser ; \zh{公平} gōngpíng : juste, équitable ; \zh{法律} fǎlǜ : loi.
\end{itemize}

\textbf{2. Structures grammaticales}
\begin{center}
\begin{tabularx}{\linewidth}{|X|X|X|}
\hline
\rowcolor{popblueL} Structure & Exemple & Traduction \\
\hline
Sujet + 把 + objet + verbe + complément & \zh{爷爷把房子留给了我。} Yéye bǎ fángzi liúgěi le wǒ. & Grand-père m'a laissé la maison. \\
\hline
Sujet + 被 + agent + verbe & \zh{土地被分给了子女。} Tǔdì bèi fēngěi le zǐnǚ. & La terre a été partagée entre les enfants. \\
\hline
Verbe + 给 + personne & \zh{他留给我一本书。} Tā liúgěi wǒ yì běn shū. & Il m'a laissé un livre. \\
\hline
Classificateurs & \zh{一块土地、一座房子} yí kuài tǔdì, yí zuò fángzi & une terre, une maison \\
\hline
\end{tabularx}
\end{center}

\textbf{3. Clés et écriture}
\begin{itemize}
  \item \zh{家} : clé 宀 (toit) ; on écrit le toit d'abord, de haut en bas.
  \item \zh{房} : clé 户 (porte) ; extérieur avant intérieur.
  \item \zh{留} : clé 田 (champ) ; partie haute avant partie basse.
  \item \zh{孩} : clé 子 (enfant) ; gauche avant droite.
\end{itemize}

\textbf{4. Méthode express : comprendre un texte}
\begin{enumerate}
  \item Lire une première fois sans dictionnaire pour saisir le thème général.
  \item Repérer les mots connus (personnes, lieux, verbes d'action).
  \item Souligner les mots-outils : \zh{把}, \zh{被}, \zh{给}, \zh{了} pour identifier qui fait quoi à qui.
  \item Répondre aux questions en recopiant la structure de la phrase du texte.
\end{enumerate}
\end{bilan}

\begin{aretenir}[Checklist avant le Bac]
\begin{itemize}
  \item $\square$ Je sais écrire et prononcer \zh{遗产}, \zh{继承}, \zh{留给} avec les bons tons.
  \item $\square$ Je connais le vocabulaire de la famille : \zh{父母}, \zh{子女}, \zh{爷爷}, \zh{孙子}.
  \item $\square$ Je sais utiliser la structure avec \zh{把} pour parler d'un objet transmis.
  \item $\square$ Je sais former le passif avec \zh{被}.
  \item $\square$ J'emploie le bon classificateur : \zh{块} pour la terre, \zh{座} pour la maison, \zh{本} pour le livre.
  \item $\square$ Je sais placer \zh{了} après le verbe pour une action accomplie.
  \item $\square$ Je reconnais les clés 宀, 户, 田, 子 et l'ordre des traits.
  \item $\square$ Je sais lire un court texte et repérer qui hérite de quoi.
  \item $\square$ Je peux comparer succession au Cameroun et en Chine en deux phrases simples.
  \item $\square$ Je réponds aux questions de compréhension avec des phrases complètes en chinois.
\end{itemize}
\end{aretenir}

\findecours

\end{document}
